%% Note de communication aux encadrants de stage — LIED.
%%
%% Compilation :
%%   cd recherche/note_communication_encadrants/latex
%%   pdflatex note_encadrants && pdflatex note_encadrants
%%
%% Les figures et les macros de nombres sont RASSEMBLÉES par
%% ../scripts/collect_figures.py. Aucune figure n'est dessinée ici, aucun
%% nombre de M4.4 n'est écrit à la main : `numbers_m4_4.tex` est la copie de
%% `m4_4_rebond_credit_soc/report/numbers.tex`, elle-même engendrée depuis
%% results/analysis/ et confrontée aux figures par tests/test_figures.py.
\documentclass[11pt,a4paper]{article}
\usepackage[T1]{fontenc}
\usepackage[utf8]{inputenc}
\usepackage[french]{babel}
\usepackage{lmodern}
\usepackage{amsmath,amssymb}
\usepackage[margin=2.4cm]{geometry}
\usepackage{booktabs,longtable,tabularx,array}
\usepackage{enumitem}
\usepackage{graphicx}
\usepackage{xcolor}
\usepackage{microtype}
\usepackage{float}
\PassOptionsToPackage{hyphens}{url}
\usepackage{hyperref}

\DeclareUnicodeCharacter{2264}{\ensuremath{\le}}
\DeclareUnicodeCharacter{2265}{\ensuremath{\ge}}
\DeclareUnicodeCharacter{03C3}{\ensuremath{\sigma}}
\DeclareUnicodeCharacter{03B3}{\ensuremath{\gamma}}
\DeclareUnicodeCharacter{03BB}{\ensuremath{\lambda}}
\DeclareUnicodeCharacter{03C1}{\ensuremath{\rho}}
\DeclareUnicodeCharacter{03B4}{\ensuremath{\delta}}
\DeclareUnicodeCharacter{03B1}{\ensuremath{\alpha}}
\DeclareUnicodeCharacter{0394}{\ensuremath{\Delta}}
\DeclareUnicodeCharacter{03B5}{\ensuremath{\varepsilon}}
\DeclareUnicodeCharacter{2080}{\ensuremath{_0}}
%% Espace fine insécable : séparateur de milliers des macros de M4.4.
\DeclareUnicodeCharacter{202F}{\,}
%% Caractères présents dans provenance.tex, engendré par collect_figures.py.
%% Les déclarer ici plutôt que d'appauvrir les légendes à la source.
\DeclareUnicodeCharacter{221D}{\ensuremath{\propto}}
\DeclareUnicodeCharacter{2212}{\ensuremath{-}}
\DeclareUnicodeCharacter{00BD}{\ensuremath{\tfrac{1}{2}}}

\hypersetup{colorlinks=true,linkcolor=blue!45!black,urlcolor=blue!45!black,
  pdftitle={Note de communication — modélisation multi-agents de l'effet rebond},
  pdfauthor={Anatole Joseph-Stouls}}
\setlist{nosep}
\setcounter{tocdepth}{2}

\newcommand{\code}[1]{\texttt{#1}}
%% Chemins longs : \texttt ne coupe pas, \path si — et sans échappement.
\newcommand{\chemin}[1]{\texttt{\path{#1}}}
\urlstyle{tt}
%% Marquage du niveau de preuve, repris de la discipline des rapports de
%% campagne. Un lecteur doit pouvoir savoir, phrase par phrase, s'il lit une
%% mesure, une déduction, une proposition à tester ou une zone d'ombre.
\newcommand{\fait}{\textbf{[Fait observé]}}
\newcommand{\inference}{\textbf{[Inférence]}}
\newcommand{\hyp}{\textbf{[Hypothèse]}}
\newcommand{\negatif}{\textbf{[Résultat négatif]}}
\newcommand{\incertitude}{\textbf{[Incertitude]}}

\graphicspath{{figures/}}
\newcommand{\figurine}[3]{%
  \begin{figure}[H]\centering
  \includegraphics[width=#2\linewidth]{#1}
  \caption{#3}\label{fig:#1}
  \end{figure}}

\input{numbers_m4_4}

\title{Une société d'agents collaboratifs\\
peut-elle produire un effet rebond ?\\
\large Rapport de stage --- raisonnement, résultats et questions ouvertes}
\author{Anatole Joseph-Stouls\\
\small LIED --- Laboratoire interdisciplinaire des énergies de demain\\
\small ENS Rennes, département mécatronique, parcours \emph{Recherche aux interfaces}\\
\small Encadrement : Éric Herbert, Petros Chatzimpiros, Christophe Goupil,
Jean-Philippe Brunetton}
\date{Campagnes : mars--août 2026 ; rédaction et vérifications : 14 septembre 2026}

\begin{document}
\maketitle

\begin{abstract}
\noindent Cette note rend compte, de façon exhaustive, d'un programme de
modélisation mené de mars à août 2026 : construire un modèle multi-agents
minimal d'une société d'entités échangeant une ressource unique comptée en
joules, vérifier quels phénomènes économiques réels s'y retrouvent, et
chercher si un effet rebond peut y émerger.

\textbf{Sur le premier volet}, le modèle retrouve endogènement une dizaine de
traits qualitatifs : un régime stationnaire borné avec renouvellement, une
différenciation durable entre créancières et débitrices sans qu'aucun type
n'ait été déclaré, une inégalité qui se loge dans les bilans et non dans la
taille productive, des cascades de faillites causales de toutes tailles avec
loi de puissance tronquée, une survie qui croît avec la taille, et une
activité agrégée à asymétrie négative dont la mémoire suit la durée de vie
des entités. Les limites restent importantes : queue de Pareto non démontrée,
absence de longs épisodes aux horizons étudiés et absence de filiation.
La recherche du corps exponentiel est conservée comme étape historique ;
elle n'est plus une cible imposée au modèle.

\textbf{Sur le second volet}, une amplification qualifiée de rebond au sens
interne est observée, sans calibration énergétique empirique. Le caractère
sous- ou super-proportionnel de la réponse à long terme dépend d'une convention
de modélisation --- la dotation de naissance suit-elle ou non l'échelle
technologique --- et ce n'est pas un détail : sous dotation fixe la réponse
est \emph{sous}-pro\-por\-tion\-nelle ($\varepsilon = \Epsilon$), sous dotation
compensée elle vaut exactement $1/(1-\gamma)$, donc fortement
super-proportionnelle. La dernière campagne, M4.4, montre où va le surcroît
de production dans les distributions --- certains quantiles extrêmes
augmentent moins que les médianes, sans suivi d'un groupe fixe --- et
découvre que l'institution de crédit employée par toute la lignée, qui paraît
partager le surplus à parts égales, asservit en réalité sur la valeur : par
contrat elle partage à $\PartageParContrat$, par joule à $\PartagePondere$,
et l'écart est exactement $\operatorname{Cov}(p,\Delta)/\mathbb{E}[\Delta]$.

\textbf{Le verdict externe reste suspendu.} Non par manque de résultats, mais
parce que plusieurs des observables du modèle n'ont pas de cible empirique
directement comparable. Cette absence est elle-même documentée comme un
résultat du travail.
\end{abstract}

\tableofcontents
\clearpage

\section*{Avertissement de lecture}
\addcontentsline{toc}{section}{Avertissement de lecture}

\paragraph{Ce que cette note est.} Un proto-rapport de stage autonome : elle décrit
le modèle, sa raison d'être, sa généalogie, et l'intégralité des résultats
obtenus, y compris ceux hérités des versions antérieures. Elle ne suppose la
lecture d'aucun autre document. Elle retrace le raisonnement, les hypothèses,
les impasses et les résultats négatifs. Le document frère
\chemin{recherche/note_resultats/} est un proto-article scientifique consacré
aux conclusions du travail ; il ne remplace pas ce récit.

\paragraph{Ce qu'elle couvre.} Le récit va des premiers modèles aux
campagnes M4.3Live-v2 et M4.4. Les résultats positifs, négatifs et les
changements d'hypothèses y restent datés ; les vérifications de rédaction
complètent les campagnes sans réécrire leur histoire.

\paragraph{Le marquage des énoncés.} Le stage s'impose une règle de confiance
explicite. Chaque affirmation porte son niveau :

\begin{itemize}
\item \fait{} --- mesuré dans les données ou lu dans le code, protocole
      documenté ;
\item \inference{} --- déduit de faits observés, sans démonstration causale
      complète ;
\item \hyp{} --- proposition falsifiable qui motive une expérience ;
\item \negatif{} --- effet cherché et non trouvé, ou lecture antérieure
      invalidée ;
\item \incertitude{} --- zone d'ombre assumée.
\end{itemize}

Les interrogations, incomplétudes et anomalies documentaires ne sont pas
effacées : elles sont en notes de bas de page, au point du texte qu'elles
concernent.

\paragraph{Les nombres et les figures.} Les 268 macros numériques de M4.4
sont copiées du rapport de campagne, qui les engendre depuis les
fichiers d'analyse. Cette garantie ne couvre pas les valeurs littérales
encore présentes dans le texte et les tableaux. Un test de la suite confronte les valeurs annotées sur
une figure à la macro correspondante\footnote{Ce test a immédiatement trouvé
une dérive réelle : deux documents de suivi portaient 93,4~\% là où la valeur
engendrée dit \PartAnnuleeKzero~\%. La correction est datée dans le journal du
programme.}. Les figures historiques viennent des campagnes ou de Simulation
Lab. Les compléments de rédaction sont reconstruits depuis les données
conservées, avec leurs points sources et leurs incertitudes. Leur provenance --- run source, script générateur, empreinte du
fichier --- est en annexe~\ref{sec:provenance}.

\paragraph{Parcours de lecture.}
Pour comprendre le résultat directeur, lire les parties I et II, puis le volet
rebond et la campagne M4.4. La partie III retrace les bifurcations et les
impasses : elle peut être différée si l'on cherche seulement le modèle final.
La partie IV est un registre de confrontations et résultats historiques,
dont les résultats négatifs sont regroupés. La dernière partie distingue
limites et expériences encore souhaitées. Les annexes servent à retrouver
les configurations, les points représentés et les sources.

\paragraph{Deux temporalités.}
Les campagnes historiques gardent leur date et leur protocole ; les calculs
complémentaires effectués pendant la rédaction sont datés et sourcés
séparément. Le rapport n'exclut pas un résultat nouveau au seul motif qu'il est
postérieur au stage. Les supports de questions ont servi à élaborer le texte ;
leur lecture n'est pas nécessaire pour le comprendre.

\paragraph{Comment lire une figure.}
Une trajectoire brute montre une réalisation ; une moyenne exige de préciser
l'unité de réplication et l'incertitude. Les nouveaux graphiques montrent les
graines ou leurs IC95, sans relier arbitrairement des catégories. Les courbes
analytiques et les trajectoires temporelles sont explicitement distinguées
des ajustements statistiques. Une droite log--log visible motive un test ;
elle ne remplace pas une validation de loi.

\clearpage

\part{La question, et pourquoi ce modèle}

\section{L'effet rebond, et l'hypothèse qu'il soit endogène}

\subsection{La commande}

L'objectif confié au début du stage est de \emph{chercher une explication
générale de l'effet rebond}, en adoptant un point de vue physique sur le
métabolisme des sociétés industrielles : flux d'énergie et de matière,
accumulation, dissipation, réseaux de distribution, organisation collective.

L'effet rebond désigne le fait qu'une amélioration de l'efficacité
énergétique ne se traduise pas par une baisse proportionnelle de la
consommation. Les explications usuelles sont comportementales ou
microéconomiques : le service devient moins cher, on en consomme davantage.

\hyp{} \textbf{L'hypothèse directrice de ce travail est différente : et si le
rebond était un phénomène \emph{endogène à l'organisation} de la société,
c'est-à-dire une propriété de la manière dont les agents économiques sont
couplés entre eux, plutôt qu'une somme de décisions individuelles ?}

Cette hypothèse a une origine datable. La lecture de Hendrick, Rinaldo et
Manoli, \emph{A stochastic theory of urban metabolism} (2025), montre des
distributions urbaines à \emph{scaling} de taille finie qui s'effondrent, après
remise à l'échelle, sur des courbes communes. Leurs auteurs signalent que cette
universalité ressemble à ce que produisent les processus auto-critiques, tout
en précisant qu'ils n'en étudient pas la dynamique. Le passage de là à « une
société économique pourrait être auto-critique, et le rebond en serait une
signature » est une hypothèse formulée à partir de cette lecture, pas un
résultat de l'article\footnote{Le stage ne revendique pas non plus l'idée
générale d'appliquer l'auto-organisation critique à l'économie : Bak, Chen,
Scheinkman et Woodford proposent dès 1992 un modèle-jouet d'économie de
production auto-critique où de petites perturbations sectorielles engendrent de
grandes fluctuations agrégées. L'apport visé ici est le prolongement vers le
métabolisme industriel et vers le rebond.}.

\subsection{Pourquoi une lecture énergétique est légitime}

Le modèle compte tout en joules. Cela demande une justification, car il
manipule aussi des prêts, des créances et des intérêts.

\hyp{} \textbf{Postulat monnaie--exergie.} Un transfert monétaire est lu comme
le transfert d'un droit socialement reconnu à mobiliser de l'exergie. C'est ce
postulat --- et lui seul --- qui autorise à représenter dans un même
modèle-jouet les flux physiques et les positions financières.

Il faut être précis sur son statut. Ce postulat s'inscrit dans la continuité de
l'approche macro-thermodynamique de Herbert, Giraud, Louis-Napoléon et Goupil,
qui fonde une macrodynamique de monde fini sur le couplage de sphères physique
et économique \emph{distinctes} ; leur travail soutient la dépendance
matérielle, pas une identité monnaie--exergie. La forte coévolution du PIB
mondial et de la consommation d'énergie entre 1965 et 2015 l'illustre, mais
cette source mesure l'énergie et non l'exergie, et documente aussi un
découplage relatif\footnote{\incertitude{} Le joule est, dans le moteur, une
unité comptable abstraite. La dynamique peut tourner sans que le postulat soit
vrai ; c'est sa \emph{lecture} comme métabolisme industriel qui en dépend
entièrement. C'est la limite la plus structurante de tout le travail, et elle
n'est pas levée.}.

\subsection{La seconde voie : les distributions}

Une deuxième entrée, économique et distributionnelle, vient de
Drăgulescu et Yakovenko puis de la synthèse de Yakovenko et Rosser. Ils
décrivent, sur les revenus, un corps exponentiel dit « thermique » couvrant la
grande majorité de la population, et une queue de Pareto dite
« superthermique », plus dynamique et hors équilibre.

Cette coexistence de deux régimes suggère qu'une même organisation économique
pourrait faire émerger plusieurs régimes macroscopiques à partir
d'interactions purement locales. C'est cette idée qui fixe le programme :
\textbf{ne coder que des interactions locales, et regarder ce que la
macro-statistique fait toute seule}.
La suite des lectures a déplacé l'objectif : l'auteur ne retient plus
le corps exponentiel comme cible universelle obligatoire. La diversité des
revenus, des unités statistiques et des modèles de distribution demande
une comparaison ouverte, pas un ajustement forcé à un gabarit.

\subsection{La question dorsale de cette note}

Tout ce qui suit répond à une question en deux volets.

\begin{quote}
\emph{Si l'on modélise la société comme un système d'agents collaboratifs,
quels phénomènes réels retrouve-t-on dans les comportements du système --- et
qui viennent donc soutenir l'interprétabilité du modèle ? Et parvient-on à
reproduire un phénomène qui puisse s'identifier à l'effet rebond ?}
\end{quote}

Le premier volet est traité en partie~\ref{part:phenomenes}, sous forme de
registre : chaque phénomène candidat reçoit un statut, y compris « absent » et
« non testable ». Le second est traité en partie~\ref{part:rebond}. La
partie~\ref{part:m44} donne les résultats complets de la dernière campagne.

\subsection{Le cahier des charges, resté stable}

Six exigences, tenues du premier moteur au dernier :

\begin{enumerate}
\item ne coder que des interactions \textbf{locales} ;
\item observer les macro-statistiques \textbf{émergentes}, sans les imposer ;
\item vérifier les \textbf{classes de distribution} du capital, du revenu et
      de la valeur nette ;
\item mesurer la propagation par le \textbf{rapport de branchement} ;
\item rendre les sorties \textbf{pilotables} par les paramètres microscopiques ;
\item rendre chaque version \textbf{comparable} à la précédente.
\end{enumerate}

La sixième exigence a pris, en cours de route, une importance qui n'était pas
prévue : elle est devenue une chaîne de tests de parité bit à bit reliant les
sept moteurs successifs, et c'est ce qui permet aujourd'hui d'hériter des
résultats anciens sans les rejouer (§\ref{sec:parite}).

\clearpage

\part{Le modèle : ce qu'il fait, exactement}

\section{L'entité, le joule, et ce qui circule}

\subsection{Une seule sorte d'agent}

\fait{} Le moteur ne connaît qu'un type d'objet : l'\textbf{entité}. Il n'y a
ni firme, ni ménage, ni banque, ni travailleur déclaré. Toutes les entités
suivent exactement les mêmes règles. Toute différenciation observée est
\emph{émergente}\footnote{Ce point est historiquement important : les premiers
moteurs comportaient un mécanisme bancaire spécial, retiré délibérément. Le
motif était double --- éviter que le modèle repose sur des entités immortelles
et réduire une complexité qui aurait rendu l'étude analytique beaucoup plus
difficile --- et, une fois toutes les entités sous règles communes, la
réapparition ou non de rôles durables devenait une véritable expérience
d'émergence.}.

Chaque entité $i$ possède un capital réel productif $K_i \ge 0$, compté en
joules. Un contrat de prêt lie une prêteuse $\ell$ à une emprunteuse $b$ par un
principal nominal $q$ et un taux par pas $r$. On en déduit les créances $C_i$,
les dettes $D_i$, et la valeur nette
\[
NW_i = K_i + C_i - D_i.
\]

\paragraph{Une précision de vocabulaire qui compte.} Le symbole $K$, appelé
historiquement « capital », désigne d'abord la \textbf{taille intrinsèque
productive} de l'entité. Il ne correspond pas au capital comptable ni au
patrimoine. Pour une comparaison avec un patrimoine net, c'est $NW$ qui est
l'analogue le plus proche --- et encore n'est-ce qu'un proxy, le moteur ne
valorisant ni prix de marché, ni actifs hétérogènes, ni passifs autres que les
prêts internes. Le dictionnaire complet des correspondances est en
annexe~\ref{sec:dictionnaire} ; il est à lire avant toute comparaison externe.

\paragraph{L'asymétrie constitutive.} Le capital est réel, productif et
dépréciable ; le principal est \textbf{nominal, perpétuel et non amorti}. Le
réseau de dettes peut donc rester lourd alors que les actifs réels se
déprécient ou s'effondrent. C'est un choix, et il porte une hypothèse
d'information : la créance reste connue de la prêteuse à sa valeur comptable
tandis que le capital placé chez l'emprunteuse s'érode, si bien que la prêteuse
accumule une fragilité qu'elle ne perçoit pas\footnote{\incertitude{} Ce
mécanisme est scientifiquement assumé pour les premiers moteurs, où il était
central. Sa conservation dans les lignées récentes est de nouveau discutable :
elle n'a pas été réexaminée.}.

\subsection{Stocks, flux et correspondance physique--économie}
\label{sec:physique}
Le parallèle entre physique et économie est le centre de la démarche, non une
analogie ajoutée après calcul. Il faut néanmoins distinguer une correspondance
de dimensions, une hypothèse de représentation et une validation empirique.
Le modèle est une étude fondamentale d'interactions : l'entité reste générique,
et peut éclairer un système concret si les mécanismes pertinents y sont
effectivement présents, sans devenir implicitement une personne puis une firme.

\begin{center}
\begin{tabularx}{\linewidth}{@{}lXX@{}}
\toprule
Objet & Dimension et opération & Lecture économique proposée \\
\midrule
$K_i$ & Stock, en joules du modèle & Taille productive mobilisable \\
$A_iK_i^{\gamma_i}$ & Production ajoutée pendant un pas, en joules & Flux productif, pas un salaire \\
$A_iK_i^{\gamma_i}/\Delta t$ & Puissance moyenne sur un pas & Capacité de production par unité de temps \\
$q$ & Principal nominal, même unité de compte que $K$ & Créance et dette, non amorties \\
$rq$ & Transfert pendant un pas & Service d'intérêt, pas création de ressource \\
$NW_i$ & Stock comptable $K_i+C_i-D_i$ & Position de bilan, distincte du stock physique \\
\bottomrule
\end{tabularx}
\end{center}
Un watt vaut un joule par seconde. Les sorties du moteur sont en joules
\emph{par pas} ; les appeler des watts nécessiterait de fixer $\Delta t$ en
secondes, ce qui n'a pas été calibré.\footnote{Dans l'écriture
$K\leftarrow K+A K^\gamma$, $A$ a pour dimension
$\mathrm{J}^{1-\gamma}$ pour la production d'un pas. Le coefficient d'une
équation différentielle correspondante aurait la dimension
$\mathrm{J}^{1-\gamma}/\mathrm{s}$. Les comparer sans préciser le pas ferait
disparaître une dimension physique. Le taux $r$ est une fraction due par pas ;
$r/\Delta t$ a la dimension inverse d'un temps.}

L'hypothèse interprétative est que certains transferts monétaires représentent
des transferts de droits à mobiliser de l'exergie. Le stock réel et le réseau
nominal sont donc couplés, mais ne sont pas identiques : un prêt déplace $K$ et
inscrit simultanément une créance et une dette. La production représente un
apport au stock modélisé ; le réservoir extérieur qui le permet n'est pas
décrit. Le bilan interne est fermé comptablement, pas thermodynamiquement sur
l'ensemble économie--environnement. En outre, l'observable dite « revenu total »
est $\mathrm{prod}+\mathrm{int\_in}$ : elle additionne production et transferts
reçus et ne doit pas être assimilée sans correction à une production nette
agrégée.

\subsection{Échelles : quels paramètres sont réellement indépendants ?}
\label{sec:parametres_reduits}
En régime homogène, sans bruit ni crédit, l'ordre production puis
dépréciation donne
\[
 K_{t+1}=(K_t+A K_t^\gamma)(1-\delta),\qquad
 K_{\rm aut}=\left[\frac{A(1-\delta)}{\delta}\right]^{1/(1-\gamma)}
\]
pour $A>0$, $0<\delta<1$ et $0<\gamma<1$. Ce point fixe déterministe n'est
pas une moyenne du système bruité. Poser $x_i=K_i/K_{\rm aut}$ transforme
la production en
\[
 x_i\longleftarrow x_i+\frac{\delta}{1-\delta}x_i^\gamma .
\]
À institution fixée, l'échelle commune de $A$ et $K_0$ peut ainsi être remplacée
par le rapport $\kappa_0=K_0/K_{\rm aut}$. Pour un toy-model, éviter de compter
deux fois la même liberté d'échelle est essentiel.\footnote{Une possibilité
équivalente est de choisir $K_0$ comme unité et de retenir
$a=A K_0^{\gamma-1}$. Il s'agit ici d'éliminer \emph{une} liberté d'unité de
capital, pas de démontrer que les autres paramètres sont redondants :
$\gamma$, $\delta$, $\sigma$, $\lambda$ et l'institution restent pertinents.
Une invariance de certaines statistiques vis-à-vis de $\lambda$ n'est pas une
équivalence trajectorielle de populations de tailles différentes.}

Le rééchelonnement $K,q,K_0\mapsto c(K,q,K_0)$ et
$A\mapsto c^{1-\gamma}A$ conserve les équations homogènes et les taux
marginaux. Les seuils numériques absolus et les arrondis du moteur limitent
l'exactitude logicielle ; le recalage du pas de temps n'est pas une symétrie
exacte du marché discret. Une hausse de $A$ à $K_0$ fixe change donc
$\kappa_0$, tandis que sa compensation $K_0\mapsto
K_0\,m^{1/(1-\gamma)}$ le préserve. Les deux sont des expériences distinctes :
la compensation est un contrôle, pas une correction physique imposée.

La tension $T=K_{\rm aut}/K_{\rm eq}$, avec
$K_{\rm eq}=(Y/(N_{\rm prod}A))^{1/\gamma}$, compare l'échelle individuelle
sans crédit à l'échelle productive observée. Elle est utile notamment à
$\delta$ fixé ; elle n'est ni une variable d'état universelle, ni un autre nom
de toute la chaîne crédit--mortalité--population.


\subsection{La chronologie d'un pas}

L'ordre fait partie du modèle : le modifier changerait la date de paiement des
nouveaux prêts, l'ensemble des entités admissibles au marché, et la
construction des cascades.

\begin{enumerate}
\item \textbf{Interventions (moteurs Live et M4.4).} Les modifications
      programmées sont appliquées avant les naissances.
\item \textbf{Naissances.} $B_t \sim \mathrm{Poisson}(\lambda)$ ; chaque
      entrante reçoit une dotation $K_0$.
\item \textbf{Choc multiplicatif.} $K_i \leftarrow K_i e^{\eta_i}$, avec
      $\eta_i \sim \mathcal{N}(-\sigma^2/2, \sigma^2)$ --- de moyenne
      multiplicative unitaire, donc sans dérive imposée.
\item \textbf{Production.} $K_i \leftarrow K_i + A\,K_i^{\gamma}$, avec
      $0 < \gamma < 1$ : la production est \textbf{concave} en la taille. C'est la
      seule source de croissance réelle.
\item \textbf{Intérêts.} Chaque emprunteuse paie $\sum r q$ depuis son
      capital ; si elle ne peut pas, elle paie au prorata et entre en défaut de
      service.
\item \textbf{Dépréciation.} $K_i \leftarrow (1-\delta) K_i$. Les principaux,
      eux, ne se déprécient pas.
\item \textbf{Marché.} Dans M4B, un tour par entité vivante ; parmi $k$
      entités tirées, la plus riche prête à la plus pauvre. Dans M4.4, les
      rencontres se font par paires, sur
      $\lfloor\eta(n,\rho,\beta,n_{\mathrm{ref}})\rfloor$ tours.
      Ici $n$ est l'effectif du bassin admissible et
      $\eta=\rho n_{\mathrm{ref}}(n/n_{\mathrm{ref}})^\beta$ ;
      $\rho=\beta=n_{\mathrm{ref}}=1$ par défaut, soit $n$ tours.
      Le sens libre du prêt (\code{loan\_direction="free"}, par défaut)
      est déterminé par l'optimum productif de la paire. En technologie
      homogène, il retrouve le sens riche vers pauvre ; avec des technologies
      différentes, ce sens n'est plus garanti.
\item \textbf{Faillites.} Entrée en faillite si défaut de service ou
      $NW_i < 0$. Les créances \emph{sur} la faillie sont détruites, ses
      propres créances sont annulées, son capital résiduel est détruit. Les
      nouvelles insolvabilités ainsi produites forment la génération suivante,
      et l'on itère jusqu'au point fixe --- \textbf{dans le même pas}.
\item \textbf{Mesures.} Séries, états individuels, réseau et avalanches sont
      enregistrés sans consommer de tirage aléatoire : les trajectoires avec et
      sans enregistrement sont identiques.
\end{enumerate}

\subsection{Le crédit comme institution de coopération}

C'est le c\oe{}ur du modèle, et le point où le mot « collaboratif » de la
question dorsale prend un sens précis.

La production étant concave, une même quantité de capital produit \emph{plus}
chez une petite que chez une grosse. Un transfert de stock de la riche vers la
pauvre augmente donc la production \emph{jointe} de la paire. Le marché
réalise cela dans sa restriction homogène : la plus riche met une part de son
stock à disposition de la plus pauvre. Le noyau technologique généralisé
optimise la production jointe sans imposer ce sens en régime hétérogène.

\textbf{Il y a donc un surplus coopératif à partager}, et le taux d'intérêt est
le mécanisme de partage. Pour un contrat de principal $q$, deux quantités sont
définies par l'état du contrat : $L$, la perte de puissance extractrice de la
donneuse, et $\Delta$, le surplus coopératif de la paire. Le service vaut
\[
r\,q = L + p\,\Delta.
\]
La règle imposée \code{bargain} explore $p\in[0,1]$ ; le partage implicite
d'une autre règle peut sortir de cet intervalle. À $p=0$ la donneuse fait une opération blanche et la receveuse garde tout le
surplus ; à $p=1$ c'est l'inverse. \textbf{Cette lecture du taux comme
paramètre de partage est une contribution de la dernière campagne} : les
lignées antérieures employaient une formule --- la moyenne géométrique des
rendements marginaux --- sans savoir quel partage elle réalisait. La
section~\ref{sec:partage} montre ce qu'elle réalise, et c'est le résultat le
plus inattendu du programme.

\figurine{m01_reseau_credit}{1.0}{\textbf{Le réseau de crédit, à la fin d'un
run.} Matrice prêteuse--emprunteuse, distribution des degrés, et rôles croisés.
Aucun type n'a été déclaré, et pourtant une structure apparaît : les degrés
sont dispersés d'un à vingt contrats, et surtout \textbf{182 entités sur 228
sont simultanément prêteuses et emprunteuses}. Il n'y a pas de « banques » d'un
côté et d'« emprunteuses » de l'autre : il y a un tissu.
\emph{(Spécimen Simulation Lab, M4B à $\sigma = 0{,}5$, $k=2$, graine 14 ---
distinct du point de référence de campagne.)}}

\subsection{La faillite : la règle qui rend la propagation possible}

\fait{} C'est la simplification la plus décisive de tout le travail. Quand une
entité fait faillite :

\begin{itemize}
\item ses contrats sont \textbf{annulés} --- ses créancières perdent le
      principal, sans recouvrement ni amortisseur ;
\item son capital résiduel est \textbf{détruit} --- il ne va nulle part.
\end{itemize}

\inference{} Un défaut détruit donc directement des actifs nominaux détenus par
d'autres, et peut les rendre insolvables à leur tour. C'est ce qui transforme
une mortalité en \textbf{contagion}. Les versions antérieures, qui
redistribuaient ou amortissaient, produisaient des morts \emph{synchronisées}
mais pas de propagation (§\ref{sec:trajectoire}).

\subsection{Les invariants, et ce qu'ils garantissent}

\fait{} Le moteur maintient exactement, à chaque pas :
$K_i \ge 0$ ; $\sum_i C_i = \sum_i D_i$ ; la somme des valeurs nettes des
survivantes égale leur capital total ; et le bilan réel
\[
\Delta K_{\text{tot}} = \text{injections} + \text{choc réalisé} +
\text{production} - \text{dépréciation} - \text{destruction par faillite}.
\]

\inference{} Il n'y a donc \textbf{ni création ni fuite de joules} hors de ces
canaux nommés. Les prêts et les intérêts sont des transferts internes purs. Une
campagne complète de 594 simulations a tourné sans une seule erreur comptable.

\subsection{Ce que le modèle n'a pas}
\label{sec:manques}

Il faut l'énoncer aussi clairement que ce qu'il a, car c'est la moitié de la
réponse à la question d'interprétabilité.

\fait{} Le moteur n'a \textbf{ni prix, ni travail, ni consommation finale, ni
ménages distincts des firmes, ni filiation, ni héritage, ni transmission}. Il
ne sépare pas l'outil, l'énergie incorporée, la maintenance, le service utile
et le surplus. Les prêts sont perpétuels, nominaux et sans recouvrement. Le
marché sélectionne la plus riche et la plus pauvre d'un échantillon de
taille $k$, ce qui est une règle de rencontre très particulière.

\figurine{m02_vue_macro}{1.0}{\textbf{Le système est borné, et il le reste.}
Capital total, crédit nouveau, population et levier $\sum D / \sum K$ sur
4000 pas. Aucune dérive, aucune extinction, aucun emballement : le régime est
stationnaire avec renouvellement permanent. C'est le préalable à toute
comparaison distributionnelle --- et le premier moteur du stage, lui, divergeait
pour une partie de son espace de paramètres.}

\clearpage

%% Section de définitions — exigence d'autonomie du document.
%% Toute notation et tout terme interne employés dans la note sont ici.
\section{Définitions}
\label{sec:definitions}

Un lecteur qui n'a jamais vu ce projet doit pouvoir lire la note sans rien
chercher ailleurs. Toutes les notations et tous les termes internes employés
dans la suite sont donc rassemblés ici.

\subsection{Paramètres du modèle}

\begin{longtable}{@{}p{0.10\textwidth}p{0.86\textwidth}@{}}
\toprule
\textbf{symbole} & \textbf{sens, et plage explorée} \\
\midrule
\endhead
$\lambda$ & \textbf{intensité de naissance} : nombre moyen d'entités créées
  par pas, paramètre d'une loi de Poisson. Fixe la \emph{taille} du système
  sans changer ses grandeurs intensives. Exploré sur $[1\,;100]$. \\
$K_0$ & \textbf{dotation de naissance} : capital reçu par une entrante, en
  joules. Protège les jeunes. Exploré sur $[2\,;200]$. \\
$\sigma$ & \textbf{volatilité individuelle} : écart-type du choc log-normal
  multiplicatif appliqué au capital à chaque pas. Levier dominant de la
  démographie. Exploré sur $[0\,;1]$. \\
$\delta$ & \textbf{taux de dépréciation} du capital par pas. Fixe l'échelle
  du capital. Exploré sur $[0{,}01\,;0{,}10]$. \\
$k$ & \textbf{taille du bassin de rencontre} : nombre d'entités tirées à
  chaque tour de marché, parmi lesquelles la plus riche prête à la plus
  pauvre. Exploré sur $\{2,\ldots,10\}$ dans M4B ; vaut 2 et devient un choix
  constitutif dans les lignées récentes. \\
$A$ & \textbf{productivité} de la technologie, facteur du terme de
  production. C'est le levier d'« amélioration d'efficacité » de tout le volet
  rebond. \\
$\gamma$ & \textbf{exposant de production}, $0 < \gamma < 1$ : la production
  $A K^\gamma$ est \emph{concave} en la taille. Vaut $1/2$ dans la lignée de
  référence. \\
$\rho$ & \textbf{intensité de la phase de marché} : nombre de tours de
  rencontre par entité et par pas. Vaut 1 par défaut. C'est le seul levier
  qui déplace les exposants de queue (§\ref{sec:controle}). \\
$p$ & \textbf{partage du surplus coopératif}, imposé dans $[0,1]$ pour
  \code{bargain}, mais potentiellement hors de cet intervalle pour une autre
  règle : part du surplus
  de la paire que le service d'intérêt transfère à la prêteuse. $p = 0$
  altruiste, $p = 1$ asservissement total (§\ref{sec:partage}). \\
\bottomrule
\end{longtable}

\subsection{Variables d'état et observables}

\begin{longtable}{@{}p{0.15\textwidth}p{0.81\textwidth}@{}}
\toprule
\textbf{notation} & \textbf{définition} \\
\midrule
\endhead
$K_i$ & \textbf{capital}, ou taille intrinsèque productive de l'entité $i$, en
  joules. Réel, productif, dépréciable. \\
$C_i$, $D_i$ & \textbf{créances} et \textbf{dettes} de $i$ : sommes des
  principaux prêtés et empruntés. Nominaux, perpétuels, non amortis. \\
$NW_i$ & \textbf{valeur nette}, $NW_i = K_i + C_i - D_i$. C'est le prédicat de
  faillite : $NW_i < 0$ met l'entité en défaut. \\
$q$, $r$ & \textbf{principal} et \textbf{taux par pas} d'un contrat. Le service
  versé par pas vaut $rq$. \\
$L$, $\Delta$ & pour un contrat : \textbf{perte de puissance extractrice} de la
  donneuse, et \textbf{surplus coopératif} de la paire --- le gain de production
  jointe permis par le transfert de stock. Le service se décompose en
  $rq = L + p\Delta$. \\
\code{prod} & \textbf{production} d'une entité sur un pas, $A K_i^\gamma$. Le
  terme « puissance extractrice » désigne la même quantité. \\
\code{int\_in}, \code{int\_out} & \textbf{intérêts reçus} et \textbf{versés} par
  pas. \code{int\_in} est, dans l'ontologie retenue, l'analogue du revenu
  monétaire usuel. \\
$E_t$ & \textbf{énergie totale} présente dans le système en fin de pas,
  $E_t = \sum_i K_i(t)$. Distincte de la production totale : la concavité de la
  production et les variations d'effectif changent leur relation. \\
$\varepsilon$ & \textbf{élasticité} de la production totale à la productivité.
  La dérivée locale est $\mathrm{d}\ln(\mathrm{prod})/\mathrm{d}\ln A$ ;
  M4.4 estime un contraste fini : moyenne par graine de
  $\ln(\overline{Y}_{\text{traité}}/\overline{Y}_{\text{contrôle}})/\ln(1{,}5)$,
  sur $3000<t\le4000$. \\
$E$ & \textbf{réponse normalisée} : réponse observée divisée par la réponse
  strictement proportionnelle $(m-1)p_{\mathrm{trait}}$, où $m$ est le
  multiplicateur appliqué à $A$ et $p_{\mathrm{trait}}$ la part de production
  traitée avant intervention. $E > 1$ est qualifié de rebond dans le modèle. \\
\bottomrule
\end{longtable}

\subsection{Termes internes}

\begin{description}[leftmargin=1.3em,style=nextline]
\item[Bassin] L'échantillon de $k$ entités tiré à un tour de marché. Le
  \emph{Gini du bassin} est le coefficient de Gini du capital mesuré sur les
  entités entrant en rencontre, avant appariement.
\item[Carnet] L'ensemble des contrats de prêt vivants à un instant donné. Le
  \emph{taux moyen du carnet} est la moyenne des $r$ des contrats vivants. Un
  contrat quitte le carnet dès que l'une de ses deux extrémités meurt.
\item[Rotation du crédit] Volume de nouveaux contrats conclus par pas, rapporté
  au capital total. Mesure la vitesse à laquelle le stock productif change de
  mains. Ce n'est pas un nombre de contrats : c'est un flux rapporté à un stock.
\item[Tension d'échelle] Rapport $T = K_{\mathrm{aut}}/K_{\mathrm{eq}}$ entre le
  point fixe individuel sans crédit et le capital équivalent du système. Il
  résume bien les effets d'échelle \textbf{à $\delta$ fixé}, et \emph{échoue}
  comme paramètre d'état quand $\delta$ varie\footnote{\incertitude{} Le mot
  « tension » a été employé dans deux sens dans les documents antérieurs :
  celui-ci, et par extension la chaîne causale complète de la contraction de
  population. Dans cette note il désigne \emph{uniquement} le rapport ci-dessus,
  avec sa réserve ; la chaîne causale est nommée « la chaîne »
  (§\ref{sec:chaine}). Cette distinction terminologique est retenue dans
  la présente rédaction.}.
\item[Capital équivalent $K_{\mathrm{eq}}$] Capital qui, s'il était détenu par
  chaque entité, donnerait la production totale observée. Il sépare ce qui
  relève de l'effectif de ce qui relève de l'échelle par entité.
\item[Avalanche] Composante faiblement connexe du graphe des pertes de créances,
  restreinte aux mortes d'une même résolution de cascade
  (§\ref{sec:avalanches}).
\item[Racine] Faillite entrée directement dans la file initiale d'un pas, sans
  avoir été déclenchée par la perte d'une créance.
\item[Rapport de branchement $b_1$] $1 - \text{racines}/\text{morts}$ : fraction
  moyenne de morts qui ne sont pas des racines. Borné par 1 \emph{par
  construction}.
\item[Descendance moyenne $b_2$] Nombre moyen de couples causals par mort, lu
  sur le graphe des cascades. Peut dépasser 1 quand une victime a plusieurs
  parents.
\item[Sur-détermination] Distinguer trois mesures. L'excès d'arêtes par
  morte est $b_2-b_1=\sum_{v\text{ non racine}}(d_v-1)/N$, où $d_v$
  est le nombre de parents retenus à la génération précédente et $N$ le
  nombre total de mortes. Ce n'est pas la proportion de victimes à plusieurs
  parents parmi les non-racines vues (\code{multi\_parent\_share}) :
  une victime à trois parents contribue deux arêtes supplémentaires.
  La suffisance d'une perte pour déclencher une faillite est encore une autre
  propriété, étudiée par le plancher de suffisance du lot J.
\item[Fardeau] Somme des principaux dus par une entité, rapportée à sa
  production. C'est la variable dans laquelle la convexité du risque était
  posée en hypothèse (§\ref{sec:partage}).
\item[Contraste apparié] Comparaison graine par graine entre un bras traité et
  le contrôle qui partage son amorçage et l'état du générateur à la bifurcation.
  Cela ne garantit pas des chocs communs par entité après divergence des
  effectifs. Les intervalles des contrastes moyens de M4.4 sont des
  intervalles de Student à onze degrés de liberté ; les tableaux de dispersion
  des exposants emploient aussi des écarts-types, explicitement distingués.
\item[Bras] Une condition expérimentale d'une campagne : un jeu de paramètres
  appliqué à l'ensemble des graines.
\item[Portée] Étendue d'une intervention : \emph{globale} (toutes les entités),
  \emph{nouvelles} (les entrantes seulement) ou \emph{fraction} (une part des
  vivantes).
\item[Amorçage (\emph{burn-in})] Pas initiaux écartés de toute mesure, le temps
  que le système atteigne son régime stationnaire.
\item[Élasticité] $\mathrm{d}\ln Y / \mathrm{d}\ln X$ : variation relative de
  $Y$ par variation relative de $X$. Sauf mention contraire, $X$ est la
  productivité $A$.
\item[Coefficient de Gini] Mesure d'inégalité entre 0 (tous égaux) et 1 (tout
  chez un seul). Calculé ici sur le capital, la valeur nette, la production ou
  le revenu d'intérêt --- quatre grandeurs qu'il ne faut jamais confondre,
  puisqu'elles donnent des valeurs allant de $0{,}05$ à $0{,}61$ sur le même
  état.
\item[PRCC] Coefficient de corrélation partielle de rang, mesuré sur un plan
  d'expérience par hypercube latin. Entre $-1$ et $1$ ; mesure l'association
  monotone d'un paramètre à une sortie, les autres tenus constants.
\item[Statistique de Vuong] Test de comparaison de deux lois non emboîtées par
  rapport de vraisemblance. $|z| > 1{,}96$ tranche en faveur de l'une ; en deçà
  le test ne conclut pas --- ce qui ne signifie pas que les deux lois s'ajustent
  aussi bien, mais que \emph{ce test} ne les sépare pas.
\item[Estimateur de Hill] Estimateur de l'exposant de queue d'une loi de
  puissance, calculé sur les points au-delà d'un seuil. La \emph{couverture} est
  le nombre de points retenus au-delà de ce seuil ; elle gouverne la précision.
\item[CCDF] Fonction de survie empirique, $P(S \ge s)$. Tracée en échelle
  logarithmique, une loi de puissance y est une droite et une coupure de taille
  finie un décrochement.
\end{description}


\clearpage

\part{Comment on est arrivé là}
\label{part:trajectoire}

\section{Une démarche en accordéon}
\label{sec:trajectoire}

Le programme n'a pas déroulé un plan linéaire. Il a alterné des phases
d'élargissement --- ajouter des mécanismes pour voir si un phénomène apparaît
--- et des phases d'élagage --- retirer tout ce dont on peut se passer, pour
savoir ce qui produit vraiment l'effet. Les résultats négatifs de cette
section ne sont pas des accidents de parcours : ce sont eux qui ont fixé la
forme du modèle final.

\subsection{Le résultat négatif fondateur : les classes étaient des cohortes}

La première génération de moteurs était riche en postes comptables et en règles
de marché. Elle produisait une population apparemment divisée entre
« travailleurs » et « banques », ainsi qu'une queue lourde de tailles et de
revenus. C'était l'observation la plus enthousiasmante des premiers mois.

\negatif{} L'analyse distributionnelle l'a invalidée. Sur 805 simulations du
même moteur, dont 163 retenues comme stationnaires, le faisceau est décisif :
corrélation entre âge et logarithme de la taille de
\textbf{0,88 à 0,95}\footnote{Source :
\code{recherche/analyse\_distributions\_taille\_revenu/latex/rapport.tex},
lignes 166--212. Le critère de stationnarité était calculé par le harnais
lui-même (\code{run\_simulation.py:detect\_regime\_diagnostics}) : queue bornée
\emph{et} flux de faillites équilibré avec les créations. L'ajustement portait
sur un seul instantané stationnaire par run, puis les verdicts étaient agrégés
--- afin de ne pas fabriquer une multimodalité en mélangeant des temps ou des
paramètres.},
distribution des âges bimodale avec un trou intermédiaire, association des deux
composantes ajustées aux fondatrices et aux entrantes, et alourdissement de la
queue avec l'horizon --- l'exposant dérivant de 5,4 à 3,1.

\inference{} La séparation observée était donc portée par deux
\textbf{cohortes d'âge} --- une cohorte fondatrice avantagée et des entrantes à
vie courte --- et non par deux classes économiques stables\footnote{Ce résultat
porte sur cette lignée et son protocole d'amorçage ; il ne doit pas être
présenté comme une propriété universelle de toute ancienne simulation. Il a été
établi par un rapport d'agent fortement étayé par les scripts et les sorties,
puis adopté dans la suite du travail.}.

\figurine{t01_artefact_cohorte}{1.0}{\textbf{Le résultat négatif fondateur.}
Ce qui ressemblait à deux classes économiques est une structure d'âge. C'est ce
constat qui a imposé toute la discipline méthodologique de la suite : pas de
mise en commun entre régimes ni entre graines, contrôle explicite des cohortes,
test de stationnarité avant toute mesure de forme, et comparaison de lois avec
troncature correcte.}

\figurine{t02_derive_queue}{0.85}{\textbf{Et la queue n'était pas
stationnaire.} L'exposant ajusté dérive avec l'horizon d'observation. Une
régression log-log à bon $r^2$ ne suffit donc pas à identifier une loi de
puissance --- leçon qui vaut pour tout le reste de la note.}

\subsection{M2 et M3 : le crédit devient causal, mais ne propage pas}

\negatif{} Dans \textbf{M2}, la démographie et les distributions étaient
essentiellement celles d'un processus de croissance-renouvellement de type
Reed. L'ablation complète du crédit ne changeait pas les formes observées : le
crédit était \emph{distributionnellement neutre}.

\fait{} \textbf{M3} sépare liquidité et capital productif, et rend le crédit
causal sur le niveau démographique : il divise approximativement par deux la
population et l'espérance de vie. Mais il ne modifie toujours pas de manière
convaincante les familles distributionnelles\footnote{Cette neutralité était
elle-même conditionnelle à la règle de comportement : elle disparaissait sous
une règle de revenu alternative, sans pour autant faire apparaître
d'auto-organisation critique.}.

\negatif{} Surtout, M3 produit de grands groupes de morts sous chocs
sectoriels corrélés --- et l'analyse de leur structure montre que ce sont des
morts \textbf{synchronisées}, pas une propagation : rapport racines/taille égal
à 1, profondeur au plus 2. \textbf{Synchroniser n'est pas contaminer}, et la
distinction n'était visible qu'en construisant explicitement le graphe causal.

\subsection{M4 : la percée}

\fait{} La percée consiste à \emph{simplifier davantage} la règle de faillite :
annulation des contrats, destruction du capital résiduel, aucun amortisseur
pour les créancières. Avec un marché réalisant un tour par entité et par pas,
le système produit alors de véritables cascades causales.

\figurine{t03_branchement_m4}{0.85}{\textbf{M4 : le rapport de branchement se
stabilise.} Une fois la règle annulation--destruction en place, la fraction
moyenne de morts non racines converge vers environ 0,30 et y reste. C'est ce
qui a justifié de poursuivre cette branche --- sans démontrer pour autant une
auto-organisation critique stricte, dont 0,30 est loin de la valeur critique 1
d'un processus de branchement simple.}

\figurine{t04_taille_finie_m4}{0.85}{\textbf{Et la coupure croît avec la taille
du système.} Signature de taille finie : la loi de puissance des avalanches est
tronquée, et la troncature recule quand le système grandit. C'est une condition
nécessaire d'un candidat auto-critique, pas une preuve.}

\subsection{M4B : la réduction minimale}

\fait{} M4B est la réduction validée de cette mécanique : un seul stock réel
productif par entité, un réseau de prêts nominaux perpétuels, des chocs
multiplicatifs, une production concave, de la dépréciation, et la règle
d'annulation-destruction. Rien d'autre.

C'est sur M4B qu'a été menée la \textbf{campagne de sensibilité}, le résultat
empirique le plus solide du stage : 594 simulations distinctes, environ
22 heures CPU, 8,6~Go de données, sans erreur comptable, sans censure, et avec
une phase de confirmation indépendante. Ses résultats forment la
partie~\ref{part:phenomenes}.

\subsection{De M4.2 à M4.4 : rendre le modèle pilotable, puis le perturber}

\fait{} \textbf{M4.2} généralise la production en $A\,K^\gamma$ et découvre que
l'effet naïvement attribué à la concavité $\gamma$ est en réalité dominé par un
\textbf{déplacement d'échelle} du capital. La dynamique possède une covariance
exacte sous remise à l'échelle, et c'est le rapport $K_0/K^*_{\mathrm{aut}}$ qui
devient la combinaison structurelle à contrôler.

\figurine{t05_effet_echelle_m42}{0.85}{\textbf{M4.2 : ce n'était pas la
concavité.} L'ablation d'échelle montre que l'effet attribué à $\gamma$ est
porté par le déplacement de l'échelle de capital. C'est la première fois que le
programme distingue un effet de mécanisme d'un effet d'unités --- distinction
qui reviendra, décisive, sur le rebond lui-même.}

\fait{} \textbf{M4.2B} étudie la distribution des intérêts reçus, et rencontre
le mur de la décidabilité (§\ref{sec:pareto}). \textbf{M4.3} trouve, sur une
branche à $\gamma$ compensé, le premier levier qui épaissit simultanément la
queue des intérêts et renforce les avalanches. \textbf{M4.3Live} rend $A$ et
$\gamma$ propres à chaque entité et modifiables \emph{pendant} le run : c'est
lui qui permet enfin la mesure du rebond (partie~\ref{part:rebond}).
\textbf{M4.3Live-v2} libère le sens du prêt et mène les campagnes A/B et de
rotation\footnote{Sa traçabilité contient 153 runs v2, chiffre soutenu par les
tables persistées et par l'index de l'interface locale. Le \code{README} de
cette lignée annonce 203 runs et son annexe 207 : cette liste reprend en réalité
plusieurs campagnes v1 que le journal dit ne pas avoir copiées. Le nombre 203 ne
doit pas être repris sans qualification.}. \textbf{M4.4Rebond} demande enfin
\emph{où} le surcroît de production se dépose dans la distribution.

\section{La chaîne de parité : pourquoi les résultats anciens tiennent encore}
\label{sec:parite}

\subsection{Le problème}

Sept moteurs successifs, sept campagnes. Si chaque moteur était une réécriture,
les résultats de M4B ne diraient rien de M4.4 et il faudrait tout rejouer à
chaque version --- ce qui est hors d'atteinte en temps de calcul.

\subsection{Ce qui a été fait}

\fait{} Chaque passage de version est couvert par un test qui fait tourner
\textbf{les deux moteurs côte à côte}, sur la même graine et les mêmes
paramètres, et compare leurs séries pas à pas. \textbf{Les sept passages ont
été rejoués pour cette note le 25 août 2026 ; tous sont verts.}

\begin{longtable}{@{}p{0.26\textwidth}p{0.30\textwidth}p{0.16\textwidth}p{0.20\textwidth}@{}}
\toprule
\textbf{passage} & \textbf{ce qui est comparé} & \textbf{tolérance} &
\textbf{état, 25 août} \\
\midrule
\endhead
M4 $\to$ M4B & séries complètes, deux régimes & exacte & vert, 11~s \\
M4B $\to$ M4.2 (à $\gamma = 1/2$, $A=1$) &
  discret exact ; flottant en tolérance relative &
  $10^{-9}$ ; mesuré à $5{,}7 \cdot 10^{-13}$\footnotemark &
  vert, 84~s \\
M4.2 $\to$ M4.2B (règle géométrique) & séries complètes, deux régimes &
  $10^{-9}$ ; mesuré à $0$ & vert, 181~s \\
M4.2B $\to$ M4.3 (deux institutions) & séries complètes, trois graines &
  égalité exacte & vert, 445~s\footnotemark \\
M4.3 $\to$ M4.3Live (régime homogène) & 26 colonnes, 500 pas &
  bit à bit & vert, 59~s \\
M4.3 $\to$ M4.4 (régime homogène) & 26 colonnes, \textbf{8000 pas} &
  bit à bit & vert, 763~s \\
M4.3Live-v1 $\to$ M4.4 (sens de prêt hérité) &
  34 colonnes, 300 pas, avec intervention & bit à bit & vert, 92~s \\
\bottomrule
\end{longtable}
\addtocounter{footnote}{-1}
\footnotetext{Seul maillon qui n'est pas bit à bit, et la raison est
documentée : les formes génériques $A\gamma K^{\gamma-1}$ et
$(A\gamma/r)^{1/(1-\gamma)}$ peuvent différer d'un dernier bit des formes
spécialisées $1/(2\sqrt{K})$ et $(1/2r)^2$ de M4B. Le cahier des charges
interdisait d'écrire une branche de compatibilité fabriquant l'égalité binaire.
Le discret --- naissances, morts, contrats, avalanches --- est, lui, identique,
et une sonde longue de 1000 pas ne trouve aucune divergence discrète pour un
écart flottant accumulé de $3{,}1 \cdot 10^{-12}$.}
\stepcounter{footnote}
\footnotetext{Ce test a un statut particulier : les fichiers de M4.3 sont des
copies octet pour octet de ceux de M4.2B, si bien qu'il ne peut pas détecter
d'écart aujourd'hui. Il existe comme garde-fou de \emph{régression} --- pour
qu'une modification future de M4.3 échoue bruyamment au lieu de dériver en
silence.}

\subsection{Ce que la parité autorise, et ce qu'elle n'autorise pas}

\inference{} \textbf{Ce qu'elle autorise.} Chaque moteur \emph{contient} le
précédent comme cas particulier : M4.2 à $\gamma=1/2$ et $A=1$ \emph{est} M4B ;
M4.4 avec \code{loan\_direction="richest\_lends"} \emph{est} M4.3Live-v1. Un
résultat établi sur un prédécesseur reste donc vrai du moteur actuel sur la
restriction correspondante --- non par ressemblance statistique, mais par
identité de trajectoire.

\incertitude{} \textbf{Ce qu'elle n'autorise pas.} Chaque parité vaut
\emph{sous sa restriction}. M4.2 hors de $\gamma = 1/2$ n'est pas M4B, et rien
ne garantit que les résultats de M4B y survivent. C'est précisément pourquoi la
campagne M4.4 a commencé par re-mesurer les élasticités déjà publiées avant
d'aller plus loin (§\ref{sec:ancrages}) --- et elle les a retrouvées.

\inference{} Cette chaîne n'était pas prévue au départ. Le passage de
\code{random.Random} à \code{numpy.random.default\_rng}, par exemple, n'avait
aucune motivation scientifique particulière. C'est ensuite que la conservation
exacte de la séquence de tirages est devenue cruciale : elle permet d'attribuer
une divergence de trajectoire à la \emph{seule} modification étudiée. Une
conservation des statistiques dans la variabilité inter-graines aurait parfois
suffi ; la parité exacte est la preuve la plus forte disponible.

\section{Simulation Lab : une méthode, pas seulement une interface}

\fait{} Une interface locale (\code{simulation\_lab/}) orchestre les runs,
les stocke, les indexe et --- surtout --- \textbf{engendre systématiquement une
batterie de figures pour chaque run}, sans qu'on ait à les demander. Plus de
2000 runs y sont enregistrés.

\inference{} Ce n'est pas un confort. C'est la colonne vertébrale pratique du
travail, pour quatre raisons : contrôler ce que produisent les moteurs et les
agents ; naviguer entre graines et cellules de campagne ; \textbf{repérer des
phénomènes non anticipés} ; et reformuler les questions à partir de ce qu'on
voit. Une grande partie des figures de cette note en vient directement.

\subsection{Le mode de travail, dit franchement}

\fait{} Une part importante du code, des balayages, des figures et de la
rédaction a été produite avec des agents conversationnels. Ils n'ont jamais
décidé du programme scientifique. La chaîne de contrôle est explicite :
question posée $\rightarrow$ proposition d'agent $\rightarrow$ vérification par
le code, les tests, les figures et des branches concurrentes $\rightarrow$
décision. Un énoncé d'agent ne devient un fait que recoupé par une source
adaptée ; son seul statut de rapport ne suffit pas\footnote{Ce principe a servi
plusieurs fois. Deux implémentations concurrentes du même prompt (M2) ont
fonctionné comme réplications adversariales. Et plusieurs erreurs d'agents ont
été trouvées par ce contrôle, dont une qui a inversé une conclusion : le
critère d'autosimilarité de M4.2B comparait le seuil à $\text{seuil}/2$, donc
redescendait dans le corps de la distribution (§\ref{sec:pareto}).}. Le gain
principal est le \textbf{débit expérimental} : pouvoir demander une
caractérisation de sensibilité complète sans écrire soi-même tout le harnais et
toutes les figures.
La contribution personnelle comprend la formulation des questions, les choix
de modèle, les garde-fous, la conception des expériences et leur critique.
L'auteur retient comme critère un code qu'il aurait pu écrire et comprendre
avec suffisamment de temps ; l'assistance augmente le débit, elle ne
transfère pas la responsabilité de validation. Simulation Lab rend ce contrôle
visuel répétable. La reconstitution historiographique détaillée des déclics
personnels est un chantier distinct : le récit n'en invente pas.

\clearpage
\part{Premier volet : quels phénomènes réels le système retrouve}
\label{part:phenomenes}

\section{Comment lire cette partie}

Cette partie est un \textbf{registre}, non une liste de succès. Chaque
phénomène candidat reçoit un statut, et les absences comptent autant que les
présences : un modèle qui reproduirait tout ne serait pas informatif, et un
modèle dont on ne dirait que les réussites ne serait pas honnête.

La source principale est la campagne de sensibilité M4B --- 594 simulations,
$\delta \in [0{,}01\,;0{,}10]$, $\sigma \in [0\,;1]$, $K_0 \in [2\,;200]$,
$k \in \{2,\ldots,10\}$, $\lambda \in [1\,;100]$ --- complétée par les
lignées ultérieures, qui en héritent par la chaîne de parité.

\begin{longtable}{@{}p{0.30\textwidth}p{0.18\textwidth}p{0.44\textwidth}@{}}
\caption{Registre des phénomènes. « Retrouvé » signifie : produit endogènement
par des règles locales, pas imposé.}\\
\toprule
\textbf{phénomène} & \textbf{statut} & \textbf{réserve principale} \\
\midrule
\endfirsthead
\toprule
\textbf{phénomène} & \textbf{statut} & \textbf{réserve principale} \\
\midrule
\endhead
Régime stationnaire borné avec renouvellement & \textbf{retrouvé} &
  aucune sur la plage explorée \\
Loi de Little (effectif = flux × durée de vie) & \textbf{retrouvé, exact} &
  c'est une identité de file d'attente, pas une prédiction \\
Différenciation créancière / débitrice sans types déclarés &
  \textbf{retrouvé} & c'est un cycle de vie, pas une structure de classes \\
Inégalité logée dans le bilan, pas dans la taille productive &
  \textbf{retrouvé} & robuste ; c'est le résultat distributionnel central \\
Survie croissante avec la taille & \textbf{retrouvé} &
  effet modéré, non calibré \\
Cascades de faillites causales, toutes tailles & \textbf{retrouvé} &
  pas de cible empirique comparable (voir §\ref{sec:absence}) \\
Loi de puissance tronquée des avalanches & \textbf{retrouvé} &
  exposant $\approx 2{,}31$, non calibré \\
Coupure croissant avec la taille du système & \textbf{retrouvé} &
  signature de taille finie, pas preuve de criticité \\
Activité agrégée à asymétrie négative (montée lente, chute brutale) &
  \textbf{retrouvé} & vit à l'échelle des cascades intra-pas, pas des cycles \\
Mémoire macroéconomique fixée par la démographie & \textbf{retrouvé} &
  corrélation de rang 0,93 entre cellules \\
Queue lourde sur les revenus d'intérêt et la valeur nette &
  \textbf{présent, non prouvé} & Pareto contre log-normale : indécidable
  (§\ref{sec:pareto}) \\
Corps exponentiel « thermique » de Boltzmann & \textbf{absent} &
  jamais obtenu, ni dans M4B ni dans le prototype dédié \\
Longues récessions & \textbf{absent} &
  durée moyenne 1,93 pas, maximum 10--11 \\
Mobilité sociale & \textbf{hors périmètre} &
  ce qu'on mesure est un renouvellement démographique \\
Deux classes économiques stables & \textbf{réfuté} &
  c'était un artefact de cohortes (§\ref{sec:trajectoire}) \\
Défaut de liquidité comme canal de contagion & \textbf{réfuté} &
  silencieux sur toute la plage principale \\
\bottomrule
\end{longtable}

\section{Le point de référence, chiffré}

\fait{} Au centre de campagne --- $\lambda = 30$, $\delta = 0{,}05$,
$\sigma = 0{,}25$, $K_0 = 25$, $k = 3$ --- sur cinq graines de confirmation
(11 à 15), fenêtre $[1000\,;4000]$. Les incertitudes sont des écarts-types
inter-graines\footnote{Source :
\code{recherche/sensibilite\_m4b/report/rapport\_final.tex},
table~« Le centre de campagne », et les tables engendrées
\code{results/summary/confirm.csv} et \code{nw\_confirm.csv} de la même
campagne. Le Gini de valeur nette et la profondeur du trou d'insolvabilité
viennent de l'extension \emph{post-hoc} « valeur nette », menée sur les mêmes
runs et confirmée sur graines séparées.}.

\begin{longtable}{@{}lr@{}}
\toprule
\textbf{grandeur} & \textbf{valeur} \\
\midrule
\endhead
Population / $\lambda$ & $14{,}75 \pm 0{,}02$ pas \\
Capital par tête & $98{,}8 \pm 0{,}3$ J \\
Prêts par tête & $8{,}88 \pm 0{,}03$ \\
Dette / capital & $1{,}39 \pm 0{,}05$ \\
Part d'entités endettées & $0{,}996 \pm 0{,}002$ \\
Coefficient de Gini du capital & $0{,}074 \pm 0{,}004$ \\
Coefficient de Gini de la valeur nette & $0{,}435 \pm 0{,}006$ \\
Rapport de branchement & $0{,}2968 \pm 0{,}0017$ \\
Exposant pur des avalanches, $s \ge 2$ & $2{,}248 \pm 0{,}003$ \\
Exposant tronqué & $2{,}041 \pm 0{,}009$ \\
Coupure ajustée & $51{,}6 \pm 1{,}8$ \\
Taille maximale d'avalanche & $54 \pm 2$ \\
\bottomrule
\end{longtable}

\section{Ce qui est retrouvé}

\subsection{Un système qui se tient tout seul}

\fait{} Sur tout l'espace exploré, \textbf{aucune trajectoire établie ne
s'éteint et aucune n'atteint le garde-fou de population}. Les réponses varient
continûment : il n'y a pas de transition extinction/emballement à
l'intérieur\footnote{L'extinction observée à très petit $\lambda$ survient
uniquement quand le premier tirage de Poisson est nul, de probabilité
$e^{-\lambda}$. C'est un artefact d'amorçage, pas une frontière dynamique. Il a
fallu le vérifier pour ne pas le publier comme une transition de phase.}.

\fait{} La \textbf{loi de Little} est vérifiée partout à mieux de 1~\% :
$\bar{N}/\lambda \simeq \mathbb{E}[\text{âge au décès}]$. L'effectif
stationnaire est donc le flux de naissance multiplié par une durée de vie
\emph{endogène} --- personne n'a fixé la taille du système, elle est le produit
de la mortalité que le système s'inflige.

\inference{} C'est le préalable à tout le reste. Le premier moteur du stage
divergeait pour une partie de son espace de paramètres, et une grande partie du
travail d'avril a consisté à cartographier la frontière entre divergence et
population bornée. Ici cette frontière n'existe plus.

\subsection{Des rôles durables, sans qu'aucun rôle n'ait été déclaré}

C'est le résultat qui répond le plus directement au mot « collaboratifs » de la
question dorsale.

\fait{} La position nette de crédit $C - D$ suit un \textbf{cycle de vie} :
plongée initiale en position débitrice --- une jeune entité emprunte pour
grandir --- puis remontée vers une position créancière à mesure que les
intérêts et les prêts s'accumulent. Au centre, environ 62~\% des vivantes sont
nettes débitrices, et la corrélation entre position nette et âge vaut 0,52.

\figurine{p02_structure_nw}{1.0}{\textbf{La structure de la valeur nette.}
Cycle de vie de la position nette, et profondeur du trou d'insolvabilité au
décès. Cette profondeur médiane, rapportée au capital par tête, vaut environ
$-8{,}6~\%$ et reste \emph{invariante} le long de $\delta$, $K_0$, $k$ et
$\lambda$ : elle ne dépend guère que de la volatilité. C'est une cible
analytique forte, non encore exploitée.}

\inference{} Il y a donc bien émergence de rôles économiques distincts ---
prêteuse nette, emprunteuse nette --- à partir de règles rigoureusement
identiques. Mais c'est un rôle \textbf{lié à l'âge}, non une position de classe
héritée. La différence importe : le premier programme avait pris ce genre de
structure pour deux classes économiques, et s'était trompé.

\subsection{L'inégalité est dans les bilans, pas dans les stocks}

\fait{} C'est le résultat distributionnel le plus important du stage. Au
centre :
\[
\mathrm{Gini}(NW) = 0{,}435 \pm 0{,}006, \qquad
\mathrm{Gini}(K) = 0{,}074 \pm 0{,}004.
\]
Un facteur six. La taille productive est presque égalitaire ; ce sont les
\emph{positions accumulées} qui sont inégales.

\fait{} Les deux coefficients répondent parfois \textbf{en sens opposé}. Quand
la volatilité augmente, l'inégalité de capital augmente mais celle de valeur
nette \emph{diminue}, parce que les chocs raccourcissent la vie des vieilles
créancières. Augmenter la taille du bassin de marché écrase l'inégalité de
capital sans presque changer celle du bilan : le marché égalise les stocks
productifs, pas les positions.

\figurine{p01_lorenz}{1.0}{\textbf{L'inégalité change de nature selon ce qu'on
regarde.} Courbes de Lorenz de la taille productive (à gauche, $G = 0{,}215$)
et du revenu d'intérêt (à droite, $G = 0{,}595$), avec les instantanés
successifs superposés. Le contraste est frappant, et il est stable dans le
temps. \emph{(Spécimen à $\sigma = 0{,}5$, distinct du point de référence de
campagne : $\mathrm{Gini}(K)=0{,}074$ et $\mathrm{Gini}(NW)=0{,}435$.
Cette dernière valeur concerne la valeur nette, non les intérêts.)}}

\figurine{p08_gini_renouvellement}{1.0}{\textbf{Et cela tient sur les graines
de confirmation.} Coefficients de Gini et renouvellement mesurés sur des
graines indépendantes de celles de la campagne principale.}

\inference{} Cela oriente toute comparaison empirique : si l'on veut confronter
ce modèle à des données d'inégalité, ce n'est pas $K$ qu'il faut regarder mais
$NW$, et le canal responsable est le crédit. C'est aussi ce que la dernière
campagne retrouvera par un chemin entièrement différent
(§\ref{sec:ou_vivent_les_queues}).

\figurine{rev_lorenz}{1.0}{\textbf{Comparer les observables au même instant.}
M4.4, contrôle, douze graines, instantané $t=4000$. À gauche : courbes de
Lorenz et IC95 entre graines, chaque observable ayant son propre classement.
À droite : survies des valeurs divisées par leur médiane dans le même
instantané, avec IC95 ; les IC atteignant zéro ne sont pas affichés.
La normalisation compare les formes, pas les niveaux. Ces instantanés finaux
ne sont pas les moyennes temporelles utilisées dans les statistiques de campagne.}

\subsection{Les grandes vivent plus longtemps}

\fait{} L'espérance de vie restante croît avec la taille instantanée --- d'un
facteur environ deux du bas au haut de l'échelle observée --- et le capital de
naissance $K_0$ agit essentiellement comme une protection des jeunes : la durée
de vie croît presque linéairement avec $\log K_0$.

\figurine{p03_vie_par_taille}{1.0}{\textbf{La survie croît avec la taille.} À
gauche, les couples (taille instantanée, vie restante) ; à droite, la moyenne
conditionnelle par classes équipopulées. L'effet est net mais modéré, et la
dispersion reste très large : \textbf{être grande ne met pas à l'abri}. Aucune
entité n'est immortelle, ce qui est une propriété voulue --- les premiers
moteurs en produisaient, et c'est une des raisons du retrait du mécanisme
bancaire spécial.}

\subsection{Des cascades causales, de toutes tailles}
\label{sec:avalanches}

C'est le mécanisme central, et celui qu'il a fallu le plus de temps à obtenir.

\fait{} \textbf{Définition.} Quand une emprunteuse meurt, chaque prêteuse perd
le principal correspondant : cela crée une arête dirigée de la morte vers la
prêteuse. Une avalanche est une composante faiblement connexe de ce graphe de
pertes. Deux racines indépendantes restent deux avalanches si aucune chaîne de
perte ne les relie\footnote{La mesure inter-pas a été explorée puis écartée des
conclusions : elle peut relier fortuitement des événements denses et fabriquer
une composante géante artificielle. C'est un exemple du type de piège que ce
travail a dû apprendre à éviter.}.

\figurine{p04_avalanches_structure}{1.0}{\textbf{C'est de la propagation, pas
de la synchronisation.} À gauche, la part de racines décroît régulièrement avec
la taille de l'avalanche --- une grande avalanche n'est pas un grand paquet de
morts indépendantes. À droite, la profondeur causale atteint 10. C'est
exactement ce que M3 ne produisait pas : là, le rapport racines/taille valait 1
et la profondeur ne dépassait pas 2.}

\fait{} Les tailles suivent une \textbf{loi de puissance tronquée}, qui bat la
loi pure dans chaque cellule et chaque graine de confirmation. L'exposant
converge vers environ $2{,}31$ à grande taille. Les scalings mesurés sont
\[
s_{\max} \propto \lambda^{0{,}58}, \qquad
\chi = \frac{\langle s^2 \rangle}{\langle s \rangle} \propto \lambda^{0{,}51}.
\]

\figurine{p05_cascades_rank_size}{1.0}{\textbf{La loi et sa coupure.} Densité
avec intervalles de Wilson, fonction de survie, et volume de créances détruites.
La coupure est visible dans la fenêtre observable --- c'est pour cela que la loi
tronquée bat la loi pure. À $\lambda = 100$, la coupure sort de la fenêtre et la
loi pure redevient meilleure : les deux constats sont cohérents, et c'est un
effet de taille finie, non un changement de mécanisme.}

\figurine{p06_scaling_avalanches}{1.0}{\textbf{Les lois d'échelle.} Taille
maximale et susceptibilité contre l'intensité de naissance. Les exposants
$0{,}58$ et $0{,}51$ sont mesurés sur les graines de confirmation.}

\incertitude{} \textbf{Le plateau proche de 0,30 concerne le régime M4B
étudié, pas toute la lignée.} M4.4 mesure au contrôle
$b_1=\BUnControle\pm\BUnControleIC$ ; ses ablations montrent notamment
l'importance de $\sigma$ et $\delta$ dans l'écart entre régimes.
Le branchement est une \textbf{grandeur endogène observée}, dépendante de
la configuration et de l'institution. La SOC reste non établie ; un plafond
universel à 0,30 ne peut pas servir d'argument.

\subsection{Une hiérarchie de paramètres nette, et séparable}

\fait{} Le criblage montre cinq rôles distincts :

\begin{itemize}
\item \textbf{la volatilité $\sigma$} pilote la démographie et l'exposition ;
      la dette rapportée au capital est non monotone en $\sigma$, avec un
      maximum vers $0{,}15$--$0{,}20$ ;
\item \textbf{le bassin de marché $k$} modifie peu la démographie mais
      fortement les inégalités et les cascades : le Gini du capital passe
      d'environ $0{,}14$ à $k=2$ à $0{,}01$ à $k=10$ ;
\item \textbf{la dotation $K_0$} protège les jeunes ;
\item \textbf{la dépréciation $\delta$} fixe l'échelle du capital --- le point
      fixe individuel sans crédit ni bruit vaut
      $K_{\mathrm{aut}}=[A(1-\delta)/\delta]^{1/(1-\gamma)}$, pour
      $A>0$, $0<\delta<1$ et $0<\gamma<1$. Le stock est mesuré après
      dépréciation : $K_{t+1}=(K_t+A K_t^\gamma)(1-\delta)$.
      Ce point fixe déterministe n'est pas la moyenne stationnaire avec bruit ;
\item \textbf{l'intensité $\lambda$} est un axe de taille finie : les grandeurs
      intensives ne bougent pas.
\end{itemize}

\figurine{p07_hierarchie_prcc}{1.0}{\textbf{La hiérarchie des paramètres.}
Corrélations partielles de rang sur le plan d'expérience par hypercube latin.
La séparation des rôles est nette --- c'est ce qui rend le modèle
\emph{pilotable}, et donc utilisable pour une expérience contrôlée. Un modèle
dont tous les paramètres feraient la même chose ne permettrait aucune ablation.}

\subsection{Une activité agrégée qui ressemble à un cycle économique --- en
partie}

\fait{} L'activité alterne des phases courtes de croissance et de contraction.
Au point central, sur cinq graines de confirmation et 3000 pas chacune, en
mesurant la contraction sur l'\textbf{énergie totale} $E_t = \sum_i K_i$
--- définition retenue en juillet comme critère principal de validation, parce
qu'elle porte sur le stock réel présent et non sur un flux\footnote{Source :
\code{recherche/sensibilite\_m4b/JOURNAL.md}, entrée du 27 juillet 2026, et
les tables \code{results/summary/cycles\_*.csv} de la même campagne. Ces
métriques ont été calculées \emph{sans nouveau run}, par extension du script
d'analyse aux séries déjà enregistrées ; le protocole historique sur la
production est conservé en parallèle. C'est une extension \emph{post-hoc},
étiquetée comme telle.} Les dispersions ci-dessous sont les écarts-types entre
les cinq graines, et non les IC95 des nouveaux graphiques. La fréquence ne
compte ici que les contractions, pas les épisodes des deux signes :

\begin{itemize}
\item fréquence : $252{,}4 \pm 4{,}4$ récessions par millier de pas ;
\item durée moyenne : $1{,}930 \pm 0{,}042$ pas, contre $2{,}035 \pm 0{,}033$
      pour les expansions ; les plus longues durent 10 à 11 pas ;
\item amplitude logarithmique moyenne : $0{,}0361 \pm 0{,}0007$, soit une perte
      pic--creux représentative d'environ $3{,}5~\%$ ; le quantile 90 est vers
      $7{,}5~\%$ et la plus forte baisse observée vers $23~\%$ ;
\item asymétrie des variations : $-0{,}233 \pm 0{,}021$, toujours négative.
\end{itemize}

\fait{} \textbf{Trois traits motivent une comparaison économique, sans la valider.} L'asymétrie est
négative : montée lente, chute brutale. La volatilité décroît comme
$\lambda^{-0{,}50}$, signature d'une somme d'aléas quasi indépendants, avec un
minimum vers $\sigma \simeq 0{,}1$ interprété comme la frontière entre bruit
endogène de cascades et bruit exogène de chocs. Et le temps de corrélation
intégré de l'activité \textbf{suit la durée de vie moyenne des entités}, avec
une corrélation de rang de $0{,}93$ entre cellules : la démographie est
l'horloge de la persistance macroéconomique.

\negatif{} \textbf{Un trait ne l'est pas.} Il n'y a \emph{pas de longue
récession} : les épisodes durent deux pas en moyenne et l'asymétrie disparaît
dès qu'on lisse la série à 25 pas. Elle vit à l'échelle des cascades
intra-pas, pas à celle de grands cycles. Un modèle qui produirait des cycles
longs endogènes produirait autre chose que celui-ci\footnote{\incertitude{} La
comparaison à des données réelles n'est pas faisable en l'état : elle exigerait
de fixer ce qu'un pas représente, ce qui n'est pas déterminé. Un test de
recalage temporel du 21 août a établi qu'un changement de pas exige de
transformer ensemble $\delta$, $\sigma$, $\lambda$, mais aussi les flux par pas
$A$ et $\rho$ --- et laisse un résidu de 6 à 17~\% parce que deux phases de
marché successives ne se composent pas comme une seule phase deux fois plus
intense. La question du passage à une limite continue n'a pas été traitée.}.

\figurine{p09_cycles_structure}{1.0}{\textbf{La structure des épisodes.}
Durée, amplitude et asymétrie des phases de croissance et de contraction.
Beaucoup d'épisodes, tous courts. \emph{Attention à l'observable} : cette
figure et la suivante sont tracées sur la \textbf{production totale}
$X_t = \sum_i A K_i^\gamma$, l'observable historique de la campagne, et non sur
l'énergie totale $E_t$ dont viennent les chiffres du paragraphe précédent. Sur
la production, le système alterne environ \textbf{487 épisodes par millier de
pas} au centre, de durée moyenne proche de deux pas et de maxima de 10 à
13 pas. Les deux observables donnent la même image qualitative --- épisodes
nombreux et courts, asymétrie négative --- et des chiffres différents. Une
figure dédiée à l'énergie manquait en juillet ; les deux figures suivantes
complètent cette lacune à partir des runs conservés.}

\figurine{p10_cycles_sensibilite}{1.0}{\textbf{Et ce qui les pilote.} Mémoire
de l'activité contre volatilité, durée de vie, et volatilité contre taille du
système avec la pente $-1/2$ attendue. Sur l'énergie totale, le criblage donne
une amplitude presque entièrement gouvernée par la volatilité
($\mathrm{PRCC} = +0{,}983$) et par la dotation de naissance ($-0{,}801$,
$R^2_{\mathrm{CV}} = 0{,}989$), tandis que la \emph{fréquence} n'est expliquée
de façon robuste par aucun paramètre ($R^2_{\mathrm{CV}} < 0$). Le modèle a
donc une fréquence de contraction quasi institutionnelle et une profondeur
pilotable --- dissociation qu'on ne soupçonnait pas avant de la mesurer.
\emph{Même réserve d'observable que la figure précédente} : les panneaux sont
sur la production.}


\figurine{rev_stock_trajectoire}{1.0}{\textbf{Stock et flux : deux lectures
d'une même réalisation.} M4B, centre de confirmation, graine 11, extrait fixé
$3000<t\le3200$. La ligne représente la trajectoire temporelle brute ;
les points orange sont les pas de baisse. Pas d'incertitude de moyenne
ajoutée à une réalisation : la variabilité entre graines est montrée ensuite.}

\figurine{rev_cycles_stock}{1.0}{\textbf{Les contractions du stock sont
maintenant représentées.} M4B, cinq graines 11--15, centre de confirmation,
$1000<t\le4000$, sans lissage. À gauche, fréquence des épisodes des deux
signes ; au centre, durée moyenne des contractions ; à droite, survie des
durées. IC95 de Student entre graines. La fréquence des seules contractions
est environ la moitié de celle des épisodes des deux signes. Aucun recalage
du pas sur une année ou une autre durée réelle n'est revendiqué.}

\section{Ce qui n'est pas retrouvé, et ce qui est réfuté}

\subsection{Une cible historique abandonnée : le corps exponentiel imposé}

\negatif{} La cible de Drăgulescu--Yakovenko --- un corps exponentiel
« thermique » couvrant la grande majorité de la distribution --- \textbf{n'est
obtenue nulle part}. Ni dans M4B, ni dans M3 dont la liquidité, pourtant la
variable la plus proche de celle visée par l'argument d'échange conservatif,
n'était pas Boltzmann--Gibbs. Un prototype dédié
(\chemin{recherche/conception_boltzmann_pareto_soc/}) obtient
provisoirement une queue stable et un renouvellement non démographique, mais
pas le corps recherché. Ce constat reste dans le récit des essais, mais
n'est plus une exigence de succès du modèle : l'auteur a abandonné l'idée
d'imposer ce corps face à la diversité des distributions et de leurs lectures.

\figurine{p12_revenu_empirique}{0.85}{\textbf{Le revenu brut, sans ajustement
imposé.} Volontairement tracé sans aucune loi superposée. Le corps est large et
plat, la décroissance rapide : ce n'est pas une exponentielle nette, et ce
n'est pas une loi de puissance. Publier un ajustement ici reviendrait à
choisir la conclusion.}

\subsection{Les queues de Pareto : indices visuels et tests incomplets}
\label{sec:pareto}

Les visualisations suggèrent des queues proches de Pareto, mais le temps
disponible et la puissance des tests ont limité leur caractérisation.
Il faut distinguer le manque de points au-delà d'un seuil, la difficulté de
départager des familles et les tests restés à mener.

\fait{} La queue se voit sur les graphiques d'intérêts reçus. Le premier test
d'autosimilarité conçu pour la trancher comparait le seuil à
$\mathrm{seuil}/2$ --- donc redescendait dans le corps de la distribution.
Le contrôle conceptuellement correct est à $\mathrm{seuil} \times 2$, et là
la taille de queue médiane tombe de \textbf{113 à 17 points} : trop peu pour
conclure.

\figurine{p13_queue_interets}{0.85}{\textbf{Ce qui se voit.} La distribution
des intérêts reçus, à la baseline de M4.2B. C'est cette forme qui a motivé le
test de Pareto.}

\figurine{p14_erreur_seuil}{0.85}{\textbf{Et ce que le test correct en dit.}
L'exemple travaillé du contrôle de seuil, après correction de son orientation.
La conclusion n'est pas « il n'y a pas de Pareto » : c'est « ce test n'a pas de
quoi trancher ». La distinction est essentielle et a été maintenue depuis.}

\fait{} La dernière campagne a repris la comparaison de familles avec une
couverture accrue (§\ref{sec:ou_vivent_les_queues}). Les tests employés ne
départagent pas Pareto et log-normale tronquée dans les régimes étudiés.
Cela ne démontre pas une impossibilité de discrimination pour tout volume
de données ou tout protocole.

\hyp{} L'existence des queues de Pareto reste donc une \textbf{hypothèse de
travail visible}, et les exposants sont \emph{caractérisés sous cette
hypothèse}, jamais démontrés.

\subsection{Le renouvellement du sommet n'est pas de la mobilité sociale}

\fait{} Le décile supérieur se renouvelle intégralement en moins de cent pas.

\figurine{p11_renouvellement_top}{0.72}{\textbf{Un renouvellement total, et
très rapide.} Part du décile supérieur initial encore présente, en fonction du
décalage temporel. La chute est immédiate.}

\incertitude{} \textbf{Cela ne doit pas s'appeler mobilité sociale.} Les
identités du sommet changent en grande partie parce que les vies sont courtes,
et parce qu'il n'existe ni héritage ni filiation dans le moteur. Une mesure de
mobilité sociale exigerait au minimum des positions initiales et finales
définies, des matrices de transition entre quantiles, des horizons comparables,
et --- selon l'objet visé --- un mécanisme intergénérationnel absent
aujourd'hui. C'est un ajout identifié, non réalisé.

\subsection{Le canal de liquidité est muet}

\negatif{} Sur la plage principale $\sigma \le 0{,}5$, \textbf{aucun défaut de
service n'est observé} après amorçage sur les 594 runs. Le canal ne s'active
qu'à volatilité extrême, vers $\sigma \ge 0{,}85$, et y reste marginal.

\inference{} La mortalité de ce modèle est donc, sur toute sa plage utile, une
\textbf{contagion de bilan} : insolvabilité directe ou perte de créances en
cascade. La machinerie de défaut de service peut rester comme garde-fou, mais
elle ne doit jamais être invoquée pour expliquer un résultat --- et cela vaut
pour toute lecture « crise de liquidité » qu'on serait tenté de faire.

\section{L'absence de données comparables, comme résultat}
\label{sec:absence}

\fait{} La recherche bibliographique n'a pas seulement fourni des valeurs de
référence : elle a révélé que plusieurs questions apparemment simples n'ont pas
de réponse académique unique directement utilisable.

\begin{itemize}
\item La « distribution des revenus » dépend de l'unité, du périmètre fiscal,
      de la fenêtre, du traitement des revenus du capital, de l'ajustement
      familial.
\item La « taille d'une firme » peut désigner l'emploi, les ventes, les actifs
      ou la capitalisation --- et la loi de son \emph{niveau} ne doit pas être
      confondue avec celle de son \emph{taux de croissance}.
\item Les travaux sur les faillites mesurent la taille des firmes défaillantes,
      la fréquence des disparitions ou des événements agrégés ; \textbf{presque
      jamais une chaîne causale racine--descendants} comparable aux avalanches
      d'ici.
\item La consommation d'énergie ou d'exergie varie avec la définition du
      consommateur, l'énergie primaire ou utile, l'échelle sectorielle.
\item Les durées de récession sont décrites par des lois différentes selon les
      jeux de données et les conventions de datation.
\end{itemize}

\inference{} Ce n'est pas un retard de bibliographie. C'est un résultat sur la
\emph{construction des faits empiriques} : les données ne répondent pas
spontanément à une question théorique, elles deviennent comparables après un
travail de définition, de mesure et de choix d'échelle. L'obstacle rencontré
n'est pas l'absence de littérature --- elle est abondante --- mais l'écart
entre cette abondance et la donnée précisément nécessaire pour réfuter une
hypothèse particulière.

\incertitude{} \textbf{La conséquence épistémique doit être formulée avec
soin.} En l'absence de cible solide, on ne peut ni valider le modèle par
ressemblance, ni le rejeter sur la base d'un exposant ou d'une famille choisis
arbitrairement. \textbf{Le verdict externe est suspendu.} Le statut du modèle
est \emph{non actuellement réfuté}, ce qui est informatif pour la suite du
programme, et ne doit jamais être converti en preuve de réalisme.

\inference{} Cela déplace une partie du travail futur : avant d'ajouter de la
complexité au moteur pour améliorer un ajustement, il faut construire ou
identifier l'objet empirique qui rendrait cet ajustement probant. L'absence de
données n'est pas une raison pour multiplier les paramètres.

\clearpage
\part{Second volet : reproduit-on un effet rebond ?}
\label{part:rebond}

\section{Le problème de définition, posé avant de mesurer}

\incertitude{} Sorrell et Dimitropoulos montrent que le terme « effet rebond »
recouvre des objets différents selon la définition du service énergétique, du
coût effectif et de la réponse comportementale. Cette précaution s'applique
directement ici : \textbf{une hausse de production ou de population après un
changement de paramètre n'est pas, à elle seule, une mesure de rebond}. Il faut
définir la ressource, le service, l'amélioration d'efficacité et le
contrefactuel sans réorganisation.

\fait{} \textbf{Définition de travail retenue.} L'amélioration d'efficacité est
une multiplication de la productivité $A$ par un facteur $m$, appliquée en
cours de run. La réponse est comparée à une trajectoire \emph{contrefactuelle}
partant du même état, avec le même état du générateur à la bifurcation.
Les comparaisons sont appariées par graine ; lorsque les populations divergent,
les chocs individuels ultérieurs ne sont pas nécessairement communs. Deux mesures :
\[
E = \frac{\text{variation relative observée}}{(m-1)\,p_{\mathrm{trait}}}, \qquad
\widehat\varepsilon_s =
\frac{\ln(\overline{Y}_{s,\text{traité}}/\overline{Y}_{s,\text{contrôle}})}{\ln m},
\]
où $p_{\mathrm{trait}}$ est la part de production traitée avant intervention,
distincte du partage contractuel $p$. Pour M4.4, $Y$ est la production totale,
$m=1{,}5$ et les moyennes portent sur $3000<t\le4000$ ; l'estimation publiée
est la moyenne des contrastes logarithmiques des douze graines.
Il s'agit d'une élasticité sur un choc fini, non d'une dérivée locale. $E > 1$ est
qualifié de rebond \emph{dans le modèle}.

\incertitude{} Cette définition permet une mesure interne rigoureuse. Elle ne
résout pas la correspondance empirique entre production, service utile et
ressource physique consommée --- le moteur ne sépare pas ces trois choses.
C'est la limite la plus sérieuse de ce volet.

\section{La première expérience, et pourquoi elle ne compte pas}

\fait{} La première expérience explicitement consacrée au rebond est bien plus
ancienne que la lignée actuelle : fin avril, la variante
\code{4-05-dynamique} compare le modèle courant à une hausse de 10~\% de la
productivité, et annonce $+5{,}1~\%$ d'extraction post-transitoire et
$+25{,}2~\%$ de prêts actifs finaux.

\negatif{} \textbf{Ce résultat ne peut pas être retenu.} Son interprétation
reposait sur les « banques » ensuite reconnues comme artefacts de cohorte, et
ses archives présentent une discordance entre trois graines annoncées et deux
archivées par variante. Il a nourri la réflexion ; il n'est pas une preuve. Le
programme a dû repartir de loin --- sans repartir de zéro, grâce à la chaîne de
parité.

\section{Le verdict de M4.3Live, qui dépend de trois choses}

\fait{} M4.3Live rend $A$ et $\gamma$ propres à chaque entité et modifiables
pendant le run\footnote{Tous les chiffres de cette section viennent de
\code{m4\_3live\_credit\_soc/report/rapport\_final.tex} et des tables
\code{results/analysis/} de cette lignée : $+36{,}1 \pm 0{,}4~\%$ et le
décompte $\times 0{,}754$ population, $\times 1{,}802$ production par entité ;
$E$ entre $2{,}3$ et $2{,}9$ à portée partielle ; $E = 2{,}53$ et
$\varepsilon = 2{,}02$ sous dotation compensée.}. Une intervention peut toucher toutes les entités, seulement les
nouvelles, ou une fraction des vivantes. \textbf{Le résultat ne se résume pas à
un verdict binaire} : il dépend de la portée, de l'horizon et de la convention
d'échelle.

\begin{longtable}{@{}p{0.30\textwidth}p{0.20\textwidth}p{0.42\textwidth}@{}}
\toprule
\textbf{configuration} & \textbf{verdict} & \textbf{ce qui se passe} \\
\midrule
\endhead
Portée partielle, court horizon & \textbf{rebond} &
  $E$ entre $2{,}3$ et $2{,}9$, sur deux amplitudes et trois intensités,
  à 3 à 48 écarts-types du bruit apparié \\
Portée partielle, long horizon & \textbf{non tranchable} &
  la cohorte traitée disparaît, sa part tend vers zéro : le dénominateur de
  la mesure s'éteint \\
Portée globale, dotation $K_0$ fixe & \textbf{sous-proportionnel} &
  $A \times 1{,}5$ donne $+36{,}1 \pm 0{,}4~\%$ de production, soit
  $E \simeq 0{,}72$ et $\varepsilon \simeq 0{,}76$\footnotemark : la production
  \emph{par entité} est multipliée par $1{,}80$, mais la population par
  $0{,}754$ \\
Portée globale, dotation $K_0$ compensée & \textbf{rebond fort} &
  la population ne se contracte plus ; $E \simeq 2{,}53$ et
  $\varepsilon \simeq 2{,}02$ \\
\bottomrule
\end{longtable}
\footnotetext{Les séries v1 et M4.4 ont été relues avec le même estimateur
fini, sur $3000<t\le4000$. Le bras v1 \code{all\_A150} donne
$0{,}7560\pm0{,}0203$ (IC95, cinq graines) ; \code{new\_A150} donne
$0{,}7639\pm0{,}0269$. Le contrôle M4.4 de la réponse globale donne
$0{,}7473\pm0{,}0106$ (douze graines). L'ancienne comparaison à
$0{,}765\pm0{,}008$ mélangeait bras et conventions d'incertitude :
elle est remplacée par le tableau de vérification ci-dessous.
La proximité des intervalles n'est pas présentée comme un test de parité.}

\figurine{r01_reponse_prod}{1.0}{\textbf{La mesure centrale du rebond.} Écart
relatif de production totale par rapport à la trajectoire contrefactuelle
appariée, en fonction de l'horizon après l'intervention. \textbf{En haut}, les
portées partielles : la réponse est forte au début --- jusqu'à $+130~\%$ pour
le levier d'exposant --- puis retombe dans la bande de bruit, parce que la
cohorte traitée meurt et n'est pas remplacée par des entités traitées. \textbf{En
bas}, les portées globales : la réponse s'établit à environ $+36~\%$ pour
$A \times 1{,}5$, très au-dessus de la bande de bruit, et \emph{y reste}. C'est
la différence entre un effet transitoire et un régime.}

\figurine{rev_reponse}{1.0}{\textbf{Même estimateur, bras distincts ; deux
conventions d'échelle.} À gauche : élasticités finies calculées depuis les séries,
sur $3000<t\le4000$, cinq graines v1 et douze M4.4. Points pâles : graines ;
barres : IC95 de Student. « Toutes » modifie vivantes et futures entrantes ;
« entrantes » ne traite pas les vivantes initiales. À droite : M4.4,
écarts relatifs par graine puis moyenne, lissage centré de 101 pas effectué
avant le calcul des IC95. La ligne grise est le contrefactuel proportionnel
de $+50\,\%$. Les deux réponses sont positives.}

\begin{footnotesize}
\begin{longtable}{@{}llrrr@{}}
\toprule
Lignée & Bras & $n$ & Moy. log-rapports (IC95) & Log moy. \\
\midrule
\endhead
M4.3Live-v1 & \texttt{all\_A150} & 5 & $0{,}7560\pm0{,}0203$ & $0{,}7560$ \\
M4.3Live-v1 & \texttt{new\_A150} & 5 & $0{,}7639\pm0{,}0269$ & $0{,}7639$ \\
M4.4 & \texttt{all\_A150} & 12 & $0{,}7473\pm0{,}0106$ & $0{,}7472$ \\
M4.4 & \texttt{new\_A150} & 12 & $0{,}7495\pm0{,}0147$ & $0{,}7495$ \\
\bottomrule
\end{longtable}

\end{footnotesize}
La dernière colonne est le logarithme du rapport des niveaux moyens entre
graines, divisé par $\ln(1{,}5)$ ; l'avant-dernière est la moyenne des
log-rapports par graine. Ces opérations ne sont pas identiques.
La colonne IC95 emploie le quantile de Student exact ; les macros historiques
M4.4 utilisaient sa valeur arrondie à $2{,}201$, sans effet au chiffre publié.

\subsection{La covariance d'échelle, et pourquoi elle explique le dernier cas}

\fait{} Le dernier résultat n'est pas un accident numérique : il se
\emph{démontre}. Le moteur possède une covariance exacte sous remise à
l'échelle du capital. Sous compensation de la dotation de naissance, le niveau
agrégé varie comme
\[
A^{1/(1-\gamma)}, \qquad \varepsilon = \frac{1}{1-\gamma}.
\]
Cette loi analytique concerne les équations homogènes sous rééchelonnement
cohérent des stocks et des contrats. Les tests numériques couvrent plusieurs
valeurs de $\gamma$, mais documentent des seuils absolus non rééchelonnés
(\code{MIN\_LOAN}, \code{ZERO\_TOL}) et des ambiguïtés d'arrondi sur le nombre
de créancières. Ils ne prouvent donc pas une covariance logicielle universelle.
La relaxation après un changement de $A$ est une question distincte de
la correspondance entre deux trajectoires entièrement rééchelonnées.

\figurine{r03_covariance_echelle}{1.0}{\textbf{La loi d'élasticité, vérifiée.}
La covariance d'échelle du moteur, confrontée aux mesures. L'accord est à la
précision machine. \textbf{C'est important méthodologiquement} : cela explique
pourquoi la réponse peut être super-proportionnelle \emph{sans invoquer
rétrospectivement un mécanisme comportemental absent du code}. Un modèle où
l'on ne pourrait pas faire cette vérification laisserait la porte ouverte à
n'importe quelle interprétation.}

\subsection{Le pivot : la dotation de naissance}

\fait{} Le verdict de long terme bascule sur une seule décision de
modélisation : \textbf{le capital fourni aux nouvelles entités reste-t-il fixe
en unités absolues, ou suit-il l'échelle technologique ?}

\figurine{r02_ablation_k0}{1.0}{\textbf{Le pivot du verdict.} Niveaux atteints
selon que la dotation de naissance est maintenue fixe ou compensée. Sous
dotation fixe, la population se contracte et absorbe le gain ; sous dotation
compensée, la contraction disparaît et le gain s'ajoute.}

\inference{} La lecture économique est directe, et c'est le point le plus
intéressant de tout le volet. Quand la technologie s'améliore mais que les
\emph{entrantes} arrivent avec la même dotation absolue qu'avant, elles sont
relativement plus petites dans un monde devenu plus grand. Elles meurent plus
vite. La population se contracte, et cette contraction mange une partie du
gain. Ce n'est \textbf{pas} un mécanisme comportemental de rebond : c'est un
effet de composition démographique.

\incertitude{} Aucune des deux conventions n'est « la bonne ». La question
réelle --- est-ce qu'une économie qui devient plus productive dote ses
nouveaux entrants proportionnellement plus ? --- est une question empirique
qui n'a pas été posée à des données. \textbf{Tant qu'elle ne l'est pas, le
caractère sous- ou super-proportionnel de long terme n'a pas de traduction
empirique univoque.}

\section{Le verdict, formulé honnêtement}

\fait{} \textbf{Oui, un phénomène qualifiable de rebond est reproduit},
sous une définition opérationnelle explicite, de façon reproductible, sur des
trajectoires appariées, avec une loi d'échelle démontrée.

\incertitude{} \textbf{Mais :}

\begin{enumerate}
\item son \textbf{écart à la proportionnalité} à long terme dépend d'une convention de naissance non
      tranchée ;
\item la portée partielle, qui est le cas le plus proche d'une politique réelle
      d'efficacité, produit un effet \textbf{transitoire} qui s'éteint faute de
      transmission du traitement aux entrantes ;
\item la mesure porte sur la \emph{production}, et le moteur ne sépare pas
      ressource, service et consommation finale : \textbf{on ne mesure donc pas
      le rebond au sens de Sorrell} ;
\item l'institution de crédit est particulière, le nombre de paramètres et de
      graines est restreint, et il n'y a aucune calibration empirique.
\end{enumerate}

\inference{} La formulation qui tient est donc : \textbf{le stage n'a pas fourni
une explication générale validée de l'effet rebond réel. Il a construit un
mécanisme candidat, et montré les conditions sous lesquelles une réponse
endogène apparaît ou s'inverse.} C'est un résultat plus modeste que
l'ambition initiale, et plus solide qu'une confirmation qu'on n'aurait pas
su interroger.

\inference{} Le résultat positif à retenir, formulé sans excès : \textbf{le
mécanisme de collaboration marchande modifie la réponse agrégée à une
amélioration technologique}. Multiplier $A$ par $1{,}5$ n'entraîne pas une
multiplication mécanique de la production totale par $1{,}5$ --- ni dans un
sens ni dans l'autre. C'est l'organisation qui décide, et on peut dire par quel
canal.

\clearpage

\part{Les résultats complets de la dernière campagne (M4.4)}
\label{part:m44}

\section{Ce que M4.4 demande, et sur quoi}

Les lignées précédentes avaient établi un énoncé sur des \emph{agrégats} :
rendre la technologie plus productive n'augmente la production totale que de
$\Epsilon$~\% par pour-cent, parce que la population se contracte en même
temps. M4.4 demande \textbf{où ce surcroît se dépose dans la distribution}.

\fait{} \textbf{Protocole.} Douze graines, un amorçage partagé par graine
jusqu'au pas 2000, puis des bras branchés sur ce même état. Chaque comparaison
est \emph{appariée} : graine par graine, le bras traité est confronté au
contrôle qui partage son amorçage. Les moyennes de fenêtre portent sur les mille
derniers pas ; les intervalles sont ceux d'un Student à onze degrés de liberté.
La campagne compte \textbf{372 runs} et environ 30 heures de calcul.

\fait{} Le diagnostic de dérive ne rejette aucun bras : le rapport de la
production du dernier quart de fenêtre au premier quart ne s'écarte pas
significativement de un, avec $|t| \le 1{,}83$ contre un seuil de
$2{,}201$\footnote{L'étendue par run, de $0{,}952$ à $1{,}047$, est le plancher
de bruit. C'est ce qui explique qu'une bande fixe de $\pm 1~\%$, employée par la
lignée précédente, rejetait onze graines sur douze d'un bras \emph{sans
intervention} : le critère était plus étroit que le bruit qu'il devait tolérer.}.
Ce test de moyenne entre graines ne prouve ni la stationnarité de chaque
trajectoire ni celle de toutes les observables.

\section{Résultat 1 --- la lignée est reproduite avant d'être prolongée}
\label{sec:ancrages}

Le moteur est un fork, les amorçages sont neufs, la campagne est neuve. Il
fallait donc vérifier que les grandeurs déjà publiées se retrouvent.

\begin{longtable}{@{}p{0.42\textwidth}p{0.24\textwidth}p{0.24\textwidth}@{}}
\toprule
\textbf{élasticité par rapport à la productivité} & \textbf{ce programme} &
\textbf{lignée précédente} \\
\midrule
\endhead
production totale ($\varepsilon$) & $\Epsilon \pm \EpsilonIC$ & $0{,}7473$ \\
effectif de fin de pas & $\ElasticitePop \pm \ElasticitePopIC$ & $-0{,}7043$ \\
effectif producteur & $\ElasticiteNProd \pm \ElasticiteNProdIC$ & $-0{,}6834$ \\
\bottomrule
\end{longtable}

\fait{} Aucune n'est distinguable de la valeur publiée. L'identité qui relie
ces trois grandeurs --- la définition même du capital équivalent --- se referme
à $\ResiduIdentite$.

\figurine{f01_ancrages}{0.86}{\textbf{Les ancrages tiennent.} Les trois
élasticités mesurées sur douze graines appariées, avec leur intervalle de
Student, et les valeurs publiées antérieurement en traits verticaux.}

\incertitude{} Une identité ne confirme rien : elle \emph{sépare} ce qui relève
de l'effectif de ce qui relève de l'échelle par entité. C'est une distinction
que le programme s'impose systématiquement, entre un résultat, une identité et
un corrélat.

\section{Résultat 2 --- des hausses de quantiles inégalement réparties}

\fait{} Rapport entre le bras traité et son contrôle apparié, quantile par
quantile. Ce sont des comparaisons de distributions : les populations et les
identités peuvent changer. Elles ne mesurent pas la captation du surplus par
un groupe fixe suivi dans le temps :

\begin{small}
% Généré : rapports par graine, IC95 Student (11 ddl).
\begin{longtable}{@{}lrrrr@{}}
\toprule
Grandeur & q10 & q50 & q99 & q99,9 \\
\midrule
\endhead
Production & $1{,}748 \pm 0{,}004$ & $1{,}806 \pm 0{,}003$ & $1{,}831 \pm 0{,}004$ & $1{,}816 \pm 0{,}008$ \\
Revenu brut & $1{,}755 \pm 0{,}004$ & $1{,}776 \pm 0{,}007$ & $1{,}741 \pm 0{,}070$ & $1{,}474 \pm 0{,}060$ \\
Valeur nette & $1{,}103 \pm 0{,}007$ & $1{,}486 \pm 0{,}016$ & $1{,}452 \pm 0{,}095$ & $1{,}169 \pm 0{,}065$ \\
Capital & $1{,}386 \pm 0{,}008$ & $1{,}451 \pm 0{,}005$ & $1{,}489 \pm 0{,}006$ & $1{,}460 \pm 0{,}013$ \\
Intérêts reçus & $0{,}111 \pm 0{,}025$ & $1{,}743 \pm 0{,}013$ & $1{,}726 \pm 0{,}081$ & $1{,}447 \pm 0{,}062$ \\
\bottomrule
\end{longtable}

\end{small}
Chaque $\pm$ est un IC95 de Student sur douze rapports appariés ; tableau et
figure proviennent du même calcul.

\fait{} La production monte de $1{,}748$ au premier décile à $1{,}831$ au
99\textsuperscript{e} centile : \textbf{moins de cinq pour-cent d'inclinaison
sur un facteur $1{,}8$}. Le Gini de la production ne bouge que de
$\DeltaGiniProd \pm \DeltaGiniProdIC$.

\fait{} \textbf{Les quantiles extrêmes augmentent relativement moins.} Au
99,9\textsuperscript{e} centile, le revenu total ne monte que de $1{,}47$ et la
valeur nette de $1{,}17$, contre $1{,}78$ et $1{,}49$ dans le corps. \textbf{Le
sommet augmente moins que la médiane} pour ces deux observables ; cela
ne quantifie pas la part du surplus reçue par les mêmes grandes entités.

\fait{} \textbf{Le bas du canal d'intérêt s'effondre.} Le premier décile du
revenu d'intérêt passe de $0{,}295$ à $0{,}033$ joule par pas --- un facteur
dix vers le bas --- pendant que la médiane augmente de $24{,}2$ à $42{,}1$.
Le premier décile reste non nul : les faibles intérêts diminuent fortement.

\figurine{rev_quantiles}{1.0}{\textbf{Les quantiles et leur incertitude.}
M4.4, $A\times1{,}5$ à $K_0$ fixe, $3000<t\le4000$. Points pâles :
douze rapports par graine ; points pleins et barres : moyenne et IC95 de Student.
Les centiles sont des catégories non équidistantes en probabilité, sans
segments d'interpolation. Les populations et identités ne sont pas constantes.}

\figurine{rev_deciles}{1.0}{\textbf{Parts par décile de capital.}
Mêmes bras et fenêtre. Reclassement par capital à chaque instantané dans
chaque bras ; dénominateur : total de l'observable du même instantané.
Les IC95 portent sur les parts moyennes de chaque run. Leur proximité
n'établit pas une absence de redistribution entre groupes fixes.}

\subsection{Producteurs ou rentiers}

\fait{} Dans le bras de contrôle, \textbf{la moitié du revenu passe par le
canal d'intérêt} --- exactement $0{,}533$ --- et $45~\%$ des entités reçoivent
plus d'intérêt qu'elles ne produisent.

\fait{} \textbf{L'inégalité de ce modèle est presque entièrement dans le canal
d'intérêt.} Le Gini vaut $0{,}513$ sur les intérêts reçus, contre $0{,}052$ sur
la production et $0{,}057$ sur le capital --- un facteur dix. Le décile
supérieur de capital ne détient que $11{,}4~\%$ du capital mais capte
$18{,}5~\%$ des intérêts.

\fait{} La part du revenu venant de l'intérêt se déplace de
$\DeltaPartInteretRevenu \pm \DeltaPartInteretRevenuIC$ sous le levier de
productivité : significatif mais minuscule. \textbf{Aucun des deux canaux ne
capte au détriment de l'autre.}

\figurine{f05_position_nette}{0.95}{\textbf{Producteurs ou rentiers.}
Répartition de l'effectif, de la production et du revenu net selon la position
nette, contrôle en teinte pleine et bras de rebond en teinte claire.}

\inference{} La signature distributionnelle du rebond est donc : une population
\textbf{plus petite et plus jeune} --- l'âge moyen baisse de $\DeltaAgeMoyen$
pas --- avec des niveaux de production individuelle plus élevés et un premier
décile d'intérêts fortement réduit. Ces déplacements distributionnels ne sont
pas une trajectoire individuelle commune à chaque membre.

\section{Résultat 3 --- une tension de la lignée se dénoue}
\label{sec:chaine}

\fait{} La contraction de population s'expliquait par une chaîne : le Gini du
bassin de marché gouverne la rotation du crédit, la rotation gouverne la
mortalité, la mortalité gouverne l'effectif. \textbf{Les trois maillons n'ont
pas le même statut}, et c'est le point à établir avant tout chiffre.

\begin{itemize}
\item \textbf{Le maillon de la rotation est vérifié.} L'élasticité de la
      rotation vaut $\ElasticiteRotation \pm \ElasticiteRotationIC$ et celle du
      Gini $\ElasticiteGini \pm \ElasticiteGiniIC$. Elles sont lues dans des
      \emph{fichiers différents} --- la série pour le volume, la sonde de marché
      pour le Gini --- donc leur accord mesure quelque chose.
\item \textbf{Le maillon de l'effectif est une identité.} À l'état
      stationnaire, les morts par pas valent le taux de naissance : la
      mortalité par entité vaut l'inverse de l'effectif \emph{par
      construction}. Mesuré : $\ElasticiteMortalite$ contre
      $-(\ElasticitePop)$. C'est un contrôle de stationnarité déguisé, et il
      passe.
\item \textbf{Le maillon du milieu est le seul à contenir un exposant.}
\end{itemize}

\figurine{f07_chaine}{0.86}{\textbf{La chaîne, maillon par maillon.}
Élasticité de chaque grandeur par rapport à la productivité, avec son
intervalle apparié, et le statut de chaque maillon.}

\fait{} Les deux lignées précédentes avaient publié pour cet exposant des
valeurs \textbf{dont les intervalles ne se recouvrent pas} : $1{,}337$ pour
l'une, $1{,}260 \pm 0{,}009$ pour l'autre. Mesuré ici bras par bras, en
contraste apparié :

\begin{longtable}{@{}lr@{}}
\toprule
\textbf{bras} & \textbf{exposant} \\
\midrule
\endhead
levier de productivité, portée globale &
  $\ExposantAallAUnCinqZero \pm \ExposantAallAUnCinqZeroIC$ \\
levier de productivité vers le bas, nouvelles &
  $\ExposantAnewAZeroSeptCinq \pm \ExposantAnewAZeroSeptCinqIC$ \\
levier de productivité, nouvelles &
  $\ExposantAnewAUnCinqZero \pm \ExposantAnewAUnCinqZeroIC$ \\
levier d'exposant de production, nouvelles &
  $\ExposantAnewgZeroSixZero \pm \ExposantAnewgZeroSixZeroIC$ \\
\bottomrule
\end{longtable}

\fait{} \textbf{L'exposant n'est pas universel : il dépend du levier}, de cinq
pour-cent environ. Et \textbf{les deux valeurs publiées antérieurement tombent
chacune sur un bras différent de cette même campagne}, avec des intervalles qui
se recouvrent.

\inference{} Les deux estimations étaient donc \emph{justes, chacune pour son
corpus}. Le désaccord n'était pas une erreur de mesure mais une différence de
mélange de leviers, et il ne pouvait pas se voir tant que l'exposant était
supposé unique. Aucune des deux lignées n'avait de quoi le découvrir : il
fallait plusieurs leviers mesurés séparément sur le même corpus.

\figurine{f06_mortalite_rotation}{1.0}{\textbf{L'exposant n'est pas universel.}
À gauche, chaque point est une graine, et les droites en tirets sont la loi
$\text{mortalité} \propto \text{rotation}^{a}$ passant par le contrôle. À
droite, les quatre exposants et leurs intervalles : ceux du levier de
productivité et du levier d'exposant \emph{ne se recouvrent pas}.}

\section{Résultat 4 --- la fragilité attribuée au rebond n'est pas la sienne}

\fait{} Le rapport de branchement se mesurait par la part des faillites qui ne
sont pas des racines. Cette lignée en ajoute un second, lu sur l'arbre causal :
la descendance moyenne par mort. Les deux coïncident si et seulement si chaque
mort non racine a exactement un parent.

\fait{} Ils ne coïncident pas : $b_1 = \BUnControle \pm \BUnControleIC$ et
$b_2 = \BDeuxControle$. L'écart mesure la \textbf{sur-détermination}. Et cet
écart est remarquablement constant --- de $0{,}0609$ à $0{,}0663$ sur les six
bras --- alors que $b_1$ lui-même varie de $0{,}72$ à $0{,}80$ : la
stabilité de cet excès d'arêtes suggère une propriété du mécanisme dans ces
six bras, sans établir son invariance dans tous les régimes.

\figurine{f09_b1_b2}{0.68}{\textbf{Les deux estimateurs de branchement.}
Chaque point est un run (bras et graine), non une moyenne de cellule.
L'écart à la diagonale représente l'excès d'arêtes par morte ; le nuage montre
la dispersion observée entre runs.}

\fait{} \textbf{Monter la productivité abaisse le rapport de branchement de
$\EffetABranchement$. Compenser en même temps l'échelle de dotation à la
naissance ne laisse que $\EffetResiduelBranchement$ --- non significatif : soit
$\PartAnnuleeKzero~\%$ de l'effet annulé.}

\inference{} C'est le même verdict que la lignée avait obtenu sur la
\emph{contraction de population}, retrouvé ici sur une grandeur entièrement
différente et par une chaîne de mesure entièrement différente. Ici, la métrique
est la part des morts non racines : elle diminue lorsque $A$ augmente à
dotation fixe. Une mortalité individuelle accrue peut coexister avec cette
baisse du branchement. Le mot « fragilité » ne doit pas fusionner risque de
décès, propagation et pertes de créances.

\figurine{f10_k0}{0.68}{\textbf{Le branchement dépend de la convention d'échelle.} Rapport
de branchement sous les trois bras, avec IC95 de Student des niveaux
sur douze graines ; les effets entre bras sont, eux, appariés.}

\figurine{f08_avalanches_ccdf}{0.78}{\textbf{Et la forme de la distribution
d'avalanches survit.} Fonction de survie empirique, quatre graines par bras. La
coupure de taille finie est visible au-delà de la centaine de faillites.}


\figurine{rev_avalanches}{1.0}{\textbf{Un diagnostic visuel avec incertitudes
et pentes par graine.} M4.4, douze graines par bras, $3000<t\le4000$.
Survies moyennes et IC95 de Student à gauche ; tendances descriptives sur
$2\le S\le30$ et ruban bootstrap de runs entiers. Les IC atteignant zéro
sont omis de l'axe logarithmique. À droite : pentes individuelles et IC95
de leur moyenne. Les points cumulatifs sont dépendants ; cette régression
visuelle ne prouve pas Pareto et ne remplace pas l'estimation de densité.}

\subsection{Un écart inter-lignées, partiellement expliqué}

\fait{} L'exposant des tailles d'avalanches vaut $\AlphaTailleControle$ ici,
contre environ $2{,}31$ pour M4B. L'écart est massif, et ne vient pas de
l'estimateur --- celui employé ici rend le même nombre que celui de M4B, à zéro
près, sur les mêmes tailles.

\fait{} L'ablation, un facteur à la fois :

\begin{longtable}{@{}lrrr@{}}
\toprule
\textbf{bras} & \textbf{$b_1$} & \textbf{exposant} &
\textbf{part de l'écart fermée} \\
\midrule
\endhead
départ (cette lignée) & $\BUnControle$ & $\AlphaTailleControle$ & --- \\
dépréciation de M4B & $\AblationDeltaCinqBUn$ & $\AblationDeltaCinqAlpha$ &
  $\AblationDeltaCinqPart~\%$ \\
volatilité de M4B & $\AblationSigmaVingtCinqBUn$ &
  $\AblationSigmaVingtCinqAlpha$ & $\AblationSigmaVingtCinqPart~\%$ \\
les deux ensemble & $\AblationMquatreBBUn$ & $\AblationMquatreBAlpha$ &
  $\AblationMquatreBPart~\%$ \\
\textit{M4B publié} & \textit{0,297} & \textit{$\approx$ 2,31} & --- \\
\bottomrule
\end{longtable}

\fait{} La volatilité est le facteur dominant : à elle seule elle ferme
$\AblationSigmaVingtCinqPart~\%$ de l'écart, et les deux ensemble
$\AblationMquatreBPart~\%$.

\incertitude{} \textbf{Un huitième reste non attribué}, et il est \emph{imputé}
à la seule différence de régime qu'on ne pouvait pas tester sans changer le
moteur --- la taille du bassin d'appariement, qui est un choix constitutif et
non un paramètre. \textbf{Cette imputation n'est pas vérifiée}\footnote{C'est
la réconciliation à retenir entre les deux nombres qui apparaissent dans cette
note : le $0{,}2968$ de M4B (partie~\ref{part:phenomenes}) et le
$\BUnControle$ de M4.4 ne sont pas contradictoires. Ce sont deux régimes
différents du même moteur, et l'ablation ferme $\AblationMquatreBPart~\%$ de
l'écart de façon documentée.}.

\figurine{f19_ablation}{1.0}{\textbf{Ce qui explique l'écart avec M4B.} À
gauche, chaque barre est une ablation d'un facteur à la fois par rapport au
contrôle --- ce ne sont pas les étapes d'une séquence. À droite, le balayage de
la volatilité seule, facteur dominant, à réponse monotone.}

\section{Résultat 5 --- où les queues apparaissent, et ce que les tests laissent ouvert}
\label{sec:ou_vivent_les_queues}

\fait{} \textbf{Le capital et la production n'ont pas de queue lourde.} Les
exposants ajustés valent $\CouvertureKAlpha$ et $\CouvertureProdAlpha$ : ce ne
sont pas des lois de puissance, c'est une décroissance quasi bornée dont
l'ajustement suit le bord de la distribution.

\fait{} Les queues lourdes vivent dans le \textbf{revenu d'intérêt}
($\CouvertureInteretAlpha$) et dans la \textbf{valeur nette}
($\CouvertureNWAlpha$) --- c'est-à-dire \textbf{dans les bilans}. C'est le
résultat que M4B avait établi par les coefficients de Gini, retrouvé ici par un
chemin entièrement différent, et cette fois avec le canal identifié.

\inference{} Ajuster des classes de lois sur le capital reviendrait donc à
ajuster du bruit.

\fait{} \textbf{La couverture n'est plus l'obstacle.} M4.2B bloquait sur le
nombre de points au-delà du seuil, de médiane $113$. Ici il vaut
$\CouvertureInteretNTail$, et une campagne dédiée où le taux de naissance passe
de 30 à 50 le porte à $365$ sur le revenu d'intérêt et $458$ sur la valeur
nette. \textbf{Et l'exposant ne bouge pas} quand l'échantillon grandit :
$3{,}898 \to 3{,}918$ et $2{,}761 \to 2{,}744$.

\figurine{f11_couverture}{1.0}{\textbf{La couverture n'est plus l'obstacle.} À
gauche, les points au-delà du seuil optimal sous les deux taux de naissance,
avec la cible visée et la valeur qui bloquait M4.2B. À droite, les exposants
correspondants : ils ne bougent pas. La bande claire marque le domaine des
queues lourdes --- seuls le revenu d'intérêt et la valeur nette y sont.}

\fait{} \textbf{Les ajustements révèlent une difficulté de discrimination.} Sur le revenu d'intérêt,
\DegenereesInteret{} ajustements log-normaux sur \AjustementsInteret{}
franchissent le seuil diagnostique de \emph{dégénérescence} --- le paramètre
de forme dépasse dix, signalant une proximité de la limite en puissance,
sans identité exacte à ce seuil --- et le test de Vuong rend
$z = \VuongInteretLN$. Sur la valeur nette, \DegenereesNW{} sur
\AjustementsNW{}, et $z = \VuongNWLN$. Aucun n'est concluant.

\fait{} Le couple loi de puissance contre exponentielle, lui, tranche :
$z = \VuongInteretExp$ et $z = \VuongNWExp$.

\figurine{rev_vuong}{1.0}{\textbf{Une discrimination non concluante dans ces données.}
Statistique de Vuong contre le paramètre de forme de la log-normale ajustée.
Les croix marquent le dépassement du seuil diagnostique $\sigma=10$, non une
frontière mathématique. Un point est un ajustement ; bras et graines sont
exportés. Les lignes horizontales $\pm1{,}96$ sont des repères usuels du diagnostic.}

\inference{} La proximité d'une limite de loi de puissance dans les ajustements
log-normaux tronqués peut expliquer la difficulté de discrimination observée.
Les données et les tests employés ne départagent pas ces familles dans les
régimes étudiés ; cela ne démontre pas une impossibilité universelle.
La présence de queues de Pareto reste une hypothèse motivée notamment par
les nombreuses visualisations. Le manque de temps a limité les vérifications :
les comparaisons de familles du lot C ne remplacent pas les tests d'adéquation
absolue, d'autosimilarité et de coupure restés à compléter.

\subsection{Trois échelles d'incertitude qu'il ne faut pas confondre}

\begin{longtable}{@{}lrrrr@{}}
\toprule
\textbf{grandeur} & \textbf{$\hat\alpha$} & \textbf{inter-graines} &
\textbf{inter-instantanés} & \textbf{intra-instantané} \\
\midrule
\endhead
revenu d'intérêt & $\AlphaInteret$ & $\AlphaInteretSdGraines$ &
  $\AlphaInteretSdInstantanes$ & $\AlphaInteretSdIntra$ \\
valeur nette & $\AlphaNW$ & $\AlphaNWSdGraines$ & $\AlphaNWSdInstantanes$ &
  $\AlphaNWSdIntra$ \\
\bottomrule
\end{longtable}

\incertitude{} L'incertitude d'un instantané vaut \textbf{plus de cinq fois}
celle de la moyenne de cellule. Confondre les deux ferait paraître un exposant
cinq fois plus précis qu'il n'est.

\fait{} \textbf{Et la queue ne vit que d'un côté du bilan.} Chez les entités
\emph{créancières nettes}, l'exposant de la valeur nette vaut
$\AlphaNWCreditrices$. Chez les \emph{débitrices nettes}, $\AlphaNWDebitrices$
--- il n'y a pas de queue. C'est la lecture la plus fine du résultat
« l'inégalité est dans les bilans » : elle est dans \emph{un seul côté} du
bilan.

\fait{} \textbf{Les leviers du rebond ne contrôlent pas les exposants.} Sur les
six bras, l'exposant de la valeur nette va de $\EtendueNWMin$ à
$\EtendueNWMax$ ($\EtendueNWPct~\%$ d'étendue) et celui du revenu d'intérêt de
$\EtendueInteretMin$ à $\EtendueInteretMax$ ($\EtendueInteretPct~\%$) ---
pendant que la population varie de vingt-cinq pour-cent et la production de
quatre-vingts.

\figurine{f12_alpha_bras}{0.95}{\textbf{Les exposants sont invariants sous les
leviers du rebond.} Exposant du revenu d'intérêt et de la valeur nette pour les
six bras, avec intervalle inter-graines. Les leviers déplacent massivement les
niveaux ; ils laissent la \emph{forme} de la queue là où elle est.}

\section{Résultat 6 --- le taux d'intérêt comme partage, et une découverte}
\label{sec:partage}

C'est une contribution supplémentaire importante, au-delà du résultat directeur sur la réponse technologique.

\fait{} La règle employée par toute la lignée --- la moyenne géométrique des
rendements marginaux --- n'avait jamais été présentée comme un partage. Elle en
réalise pourtant un, et on peut le mesurer : $p = (rq - L)/\Delta$ vaut en
moyenne $\TauxmarginalPartageImplique$ \textbf{par contrat}. Toute la lignée
tournait donc, sans que ce soit su, à un partage presque exactement équitable.

\fait{} Sur douze graines et cinq valeurs du partage, la réponse est monotone
et régulière : la population décroît de $\TauxpZeroPop$ à $\TauxpUnPop$ quand
le partage passe de l'altruisme ($p=0$) à l'asservissement ($p=1$).

\fait{} \textbf{Et voici l'anomalie.} La règle historique, qui partage à
$\TauxmarginalPartageImplique$, produit une population de $\TauxmarginalPop$
--- \emph{indiscernable} de celle du bras d'asservissement total (rapport
$\RapportPopMarginalPUn \pm \RapportPopMarginalPUnIC$). Alors que le bras qui
partage \emph{effectivement} à un demi donne vingt-sept pour-cent de population
en plus (rapport $\RapportPopMarginalPDemi \pm \RapportPopMarginalPDemiIC$).

\figurine{rev_institutions}{1.0}{\textbf{La moyenne contractuelle ne prédit
pas à elle seule la démographie.} M4.4, douze graines par bras, $3000<t\le4000$.
À gauche : effectif aux cinq partages imposés et règle historique positionnée
à son partage moyen observé ; les deux coordonnées de ce dernier ont leur
IC95. Centre et droite : rapports au bras $p=0$ calculés graine par graine,
puis moyenne et IC95. L'incertitude du dénominateur est ainsi incluse.
Les catégories ne sont pas reliées ; aucune population n'est prédite par
interpolation entre deux partages non simulés.}

\subsection{L'explication d'abord proposée était fausse}

\hyp{} La première explication proposée était : si le risque de mourir est
\emph{convexe} dans le fardeau porté, une même moyenne de fardeau tue davantage
lorsqu'elle est plus dispersée. C'est l'inégalité de Jensen.

\negatif{} \textbf{La mesure la réfute.} Le risque n'est pas convexe dans le
fardeau : il est \textbf{non monotone}. Il croît d'abord fortement, culmine à un
fardeau modeste, puis \emph{décroît} sur tout le reste de l'échelle --- une
entité très endettée relativement à sa production est une entité qui a
\emph{pu} emprunter, donc une entité solide. Et la courbure ajustée sur le
corps ne survit pas au conditionnement par le capital : sur
\ClassesCapitalTotal{} classes, \ClassesCapitalConvexes{} seule reste convexe.
\textbf{Le facteur dominant est le capital, pas le fardeau}\footnote{Il fallait
le conditionnement pour le voir. Sans lui, la convexité apparente n'aurait dit
que « les petites meurent » --- un énoncé vrai et sans contenu. C'est un
exemple précis du genre de piège que la discipline de mesure du programme
existe pour attraper.}.

\fait{} L'écart de Jensen, lui, est massif et se comporte comme l'argument le
voulait : la mortalité par insolvabilité vaut $\MortaliteObservee$ ; évaluée au
fardeau \emph{moyen}, la même courbe donne $\MortaliteAuFardeauMoyen$, un
facteur $\RapportJensen$. Et cet écart croît de façon monotone avec le partage,
de $\EcartJensenPZero$ sous l'altruisme à $\EcartJensenPUn$ sous
l'asservissement.

\incertitude{} La lecture contrefactuelle --- « si toutes portaient le fardeau
moyen, la mortalité serait six fois moindre » --- \textbf{n'est pas légitime} :
déplacer toutes les entités au fardeau moyen changerait le système entier, donc
la courbe elle-même. L'énoncé qui tient porte sur la \emph{fonction mesurée},
pas sur une intervention.

\figurine{f21_hasard}{1.0}{\textbf{La convexité posée, et ce que la mesure en
dit.} À gauche, le risque d'insolvabilité à dix pas contre le fardeau, avec
l'erreur binomiale : la courbe n'est pas convexe, elle est non monotone. Les
deux traits horizontaux sont la mortalité observée et celle qu'on lirait au
fardeau moyen. À droite, la courbure ajustée dans chaque classe de capital :
elle disparaît dès que le capital est tenu fixe.}

\subsection{Une identité de pondération, et son interprétation}

\fait{} Le partage implicite se lit de deux façons, et elles ne coïncident pas.
\textbf{Par contrat}, la règle historique partage à $\PartageParContrat$.
\textbf{Pondérée par le surplus en jeu}, elle partage à
$\PartagePondere \pm \PartagePondereIC$. L'écart n'est pas une estimation,
c'est une \textbf{identité} reliant ces statistiques :
\[
p_{\text{pondéré}} - p_{\text{par contrat}}
  = \frac{\operatorname{Cov}(p, \Delta)}{\mathbb{E}[\Delta]}
  = \CovariancePartage.
\]
La valeur numérique est mesurée sur les contrats ; la relation algébrique est
exacte. Elle ne démontre pas, seule, que la différence de population entre
institutions est causée par le partage pondéré. Cette explication reste une
inférence à distinguer des mesures et des interventions discriminantes.

\fait{} Le témoin valide l'estimateur avant qu'on le lise : sur les cinq bras à
partage \emph{fixe}, où les deux lectures doivent coïncider par construction,
l'écart mesuré vaut $\TemoinPartageFixe$.

\inference{} \textbf{La règle historique n'applique pas le même partage aux
gros et aux petits contrats. Elle asservit là où il y a beaucoup à prendre et
se montre altruiste là où il y a peu.} Le $\PartageParContrat$ n'était pas
faux : il comptait des \emph{contrats} là où le système compte des
\emph{joules}.

\figurine{f22_ponderation}{1.0}{\textbf{Compter des contrats ou compter des
joules.} À gauche, le partage pondéré par le surplus contre le partage moyen
par contrat. Les bras à partage fixe tombent sur la diagonale --- c'est le
témoin, nul au dix-millième. La règle historique en est à quarante et un
points, et cet écart \emph{est} la covariance entre la part prise et la taille
de ce sur quoi elle est prise. À droite, l'écart de Jensen par bras.}

\fait{} \textbf{La forme du partage est inscrite dans la règle.} Les trois fonctions
du moteur permettent d'examiner le partage sur une grille, sans lire un run.
Cette grille n'est pas la distribution des contrats effectivement conclus :
elle ne reproduit donc pas à elle seule la covariance de la campagne.
Évaluées sur une grille de paires en régime homogène, elles donnent un partage
impliqué qui vaut \textbf{exactement un demi à la limite des contrats
infinitésimaux} et croît monotonement avec le principal --- jusqu'à
\textbf{dépasser l'unité} sur les paires les plus inégales, c'est-à-dire à
prendre plus que la totalité du surplus.

\figurine{g03_partage_analytique}{1.0}{\textbf{La covariance, dérivée des
fonctions du moteur.} Aucun run n'est lu : les courbes sont l'évaluation
directe de $p = (rq - L)/\Delta$ sur une grille de paires. À gauche contre le
principal, à droite contre le surplus coopératif. La règle part du partage
équitable et franchit l'asservissement total, en fonction de la taille du
contrat.}

\inference{} \textbf{C'est le résultat à retenir de ce programme, et il déborde
le modèle.} Une institution peut être \emph{équitable en moyenne sur les
contrats} et \emph{asservissante sur la valeur} --- et c'est la seconde lecture
que le système subit. Il n'y a besoin d'aucune non-linéarité du dommage pour
l'obtenir : la covariance suffit.

\incertitude{} Ce qui manque est une \textbf{démonstration} que le partage
impliqué croît avec le principal pour \emph{toute} paire admissible. Ce qui est
établi est une évaluation sur une grille, pas une preuve.

\subsection{Deux conséquences, et une seconde invariance}

\fait{} \textbf{Un service plus léger n'achète pas la stabilité.} Le partage
altruiste rend le système \emph{plus} contagieux : le rapport de branchement
passe de $\TauxmarginalBUn$ à $\TauxpZeroBUn$, soit $+5{,}7~\%$. Il achète de
la population --- soixante-huit pour-cent de plus --- et cette population
achète de la contagion, parce qu'elle multiplie les chemins de cascade.

\figurine{f16_bargain_systeme}{1.0}{\textbf{Ce que le partage déplace dans le
système.} Branchement, tension, rotation du crédit et taux moyen du carnet
contre le partage. Le losange est la règle historique, à l'abscisse de son
partage moyen : sur les quatre grandeurs, elle se situe du côté de
l'asservissement et non au milieu.}

\fait{} \textbf{Le Gini de valeur nette ne bouge pas.} Il vaut
$\TauxmarginalGiniNW$ sous la règle historique et $\TauxpZeroGiniNW$ sous
l'altruisme --- \textbf{six \emph{dixièmes} de pour-cent d'étendue} sur tout le
balayage, pendant que la population varie de soixante-huit pour-cent et le
capital d'un facteur trois. Le Gini du capital, lui, passe de
$\TauxmarginalGiniK$ à $\TauxpZeroGiniK$, soit quarante pour-cent de moins.

La figure~\ref{fig:rev_institutions} montre simultanément la faible variation
du Gini de valeur nette et la variation des niveaux ; la stabilité d'un Gini
n'établit pas l'invariance de toute la distribution.

\inference{} \textbf{Deux familles de leviers entièrement différentes laissent
la même chose intacte : la forme de la distribution de valeur nette.} Ce modèle
a une structure d'inégalité que ni la technologie ni le partage du surplus ne
déplacent. Un seul levier des trois testés la déplace : \textbf{l'intensité de
la phase de marché}. Ce n'est ni ce qu'on produit ni comment on se partage le
surplus qui fixe la forme de la distribution, c'est \emph{combien on échange}.

\section{Résultat 7 --- ce qui contrôle les exposants}
\label{sec:controle}

\fait{} \textbf{Critère.} Un levier \emph{contrôle} une grandeur si et
seulement si : la réponse appariée a un signe constant sur toutes les graines ;
elle est monotone sur une étendue d'au moins un facteur trois ; et son exposant,
ajusté \emph{graine par graine} et jamais en régression groupée, est
significativement non nul et de signe constant. Tout ce qui est plus faible est
une corrélation, et doit être appelé ainsi.

\begin{longtable}{@{}lrrrl@{}}
\toprule
\textbf{grandeur} & \textbf{bas} & \textbf{haut} & \textbf{exposant} &
\textbf{contrôlé} \\
\midrule
\endhead
$\hat\alpha$ du revenu d'intérêt & $\RhoAlphaInteretBas$ &
  $\RhoAlphaInteretHaut$ &
  $\RhoAlphaInteretExposant \pm \RhoAlphaInteretExposantIC$ &
  \RhoAlphaInteretControle \\
$\hat\alpha$ de la valeur nette & $\RhoAlphaNWBas$ & $\RhoAlphaNWHaut$ &
  $\RhoAlphaNWExposant \pm \RhoAlphaNWExposantIC$ & \RhoAlphaNWControle \\
rapport de branchement $b_1$ & $\RhoBUnBas$ & $\RhoBUnHaut$ &
  $\RhoBUnExposant \pm \RhoBUnExposantIC$ & \RhoBUnControle \\
descendance moyenne $b_2$ & $\RhoBDeuxBas$ & $\RhoBDeuxHaut$ &
  $\RhoBDeuxExposant \pm \RhoBDeuxExposantIC$ & \RhoBDeuxControle \\
Gini du bassin & $\RhoGiniBas$ & $\RhoGiniHaut$ &
  $\RhoGiniExposant \pm \RhoGiniExposantIC$ & \RhoGiniControle \\
\bottomrule
\end{longtable}

\fait{} \textbf{L'intensité de marché contrôle les deux exposants de queue et
les deux estimateurs de branchement}, au sens strict ci-dessus. Le contraste
avec les leviers du rebond est net : ceux-ci laissaient les exposants à quatre
pour-cent près.

\figurine{rev_rho}{1.0}{\textbf{L'intensité de marché contrôle ce
qu'aucun autre levier ne déplace.} Les deux exposants de queue, le rapport de
branchement et le Gini du bassin contre l'intensité, en échelle logarithmique,
avec des ajustements log-log par graine. Points : moyennes et IC95 sur douze
graines, sans segments entre niveaux ; courbe et ruban : moyenne et IC95 des
courbes ajustées, non intervalle de prédiction.}

\fait{} \textbf{Un point où les deux estimateurs divergent qualitativement.} Au
niveau le plus haut du levier, la descendance moyenne franchit l'unité
($\RhoBDeuxHaut$) tandis que le rapport classique reste à $\RhoBUnHaut$.

\inference{} Ce n'est \textbf{pas} une supercriticité. Le premier estimateur
vaut $1 - \text{racines}/\text{morts}$ : il est borné par l'unité \emph{par
construction}. Le second compte les couples causals et le peut, dès que la
sur-détermination est forte --- la part de morts à plusieurs parents passe ici
de sept à plus de vingt pour-cent. Ce que le franchissement signale, c'est que
\textbf{l'attribution à parent unique devient intenable dans ce régime}. C'est
le seul endroit du programme où les deux estimateurs racontent une histoire
différente, et c'est précisément pour cela qu'il en fallait deux.

\subsection{Une relation héritée qui ne survit pas}

\fait{} La lignée précédente avait établi que la rotation du crédit vaut
l'intensité de marché multipliée par le Gini du bassin. Vérifiée ici quand le
levier est la productivité --- les deux élasticités coïncident à
$1{,}3 \cdot 10^{-4}$ --- elle \textbf{ne survit pas} quand le levier est
l'intensité elle-même : la relation prédit $\RotationPredite$, le mesuré vaut
$\RotationMesuree$, soit plus de cent erreurs-types d'écart.

\inference{} Le mécanisme est lisible : le Gini mesuré \emph{après} la phase de
marché tombe de $0{,}0986$ à $0{,}0168$ sur l'étendue du balayage, contre
$0{,}1238 \to 0{,}0514$ pour celui d'avant. La part de dispersion consommée à
l'intérieur d'un même pas passe de trente-quatre à soixante-sept pour-cent. La
relation n'est donc exacte que dans la limite où la phase de marché n'épuise
pas la dispersion --- la limite où la lignée précédente travaillait.

\figurine{f18_rotation_rho}{0.72}{\textbf{Une relation héritée qui ne survit
pas au changement de levier.} La rotation mesurée contre la prédiction
$\rho \times \bar{G}$. Les deux coïncident au bas du balayage et divergent
franchement au-delà, parce que les rondes tardives trouvent un bassin déjà
égalisé par les premières.}

\section{Résultat 8 --- sur-détermination et suffisance ne se confondent pas}

Le résultat~4 mesurait que plusieurs faillites amont atteignent souvent une
même victime. On ne pouvait pas en conclure que chacune aurait \emph{suffi}.

\fait{} \textbf{L'encadrement est exact.} Pour une victime de cascade, la
valeur nette d'avant cascade est positive --- sinon elle serait une racine ---
et vaut la valeur nette de mort, plus les créances perdues, \emph{moins} les
dettes effacées. Ce dernier terme n'est pas tracé, mais il est positif : la
valeur nette d'avant est donc \emph{au plus} la somme des deux premiers, et
l'on peut certifier qu'une arête aurait suffi dès que son principal dépasse ce
majorant. Le compte obtenu est un \textbf{plancher}, jamais un plafond.

\fait{} \textbf{Le témoin s'est retourné en mesure.} Sur les victimes à une
seule arête, le plancher devrait valoir un exactement ; il vaut $0{,}981$.
L'écart n'est pas un défaut du calcul : c'est l'effacement de dette pris sur le
fait --- $\IncidenceEffacement~\%$ des victimes meurent avec une valeur nette
\emph{positive}, parce qu'une de leurs créancières est morte dans le même pas,
après que la file eut été constituée.

\fait{} \textbf{Deux conventions d'attribution, deux réponses.} En attribution
stricte génération à génération, $\PartMultiParentStricte~\%$ des morts non
racines ont plusieurs parents. En comptant toutes les créances perdues \emph{au
pas de la mort}, $\PartSurDeterminees~\%$ des victimes en ont plusieurs. C'est
le second compte qui vaut ici : une créance perdue ampute la victime, que sa
source soit ou non de la génération précédente\footnote{Les deux nombres sont
corrects, et ils ne se contredisent pas : ils comptent deux choses différentes.
Les publier ensemble sans dire cela produirait une contradiction apparente.}.

\fait{} Sur ces victimes sur-déterminées, le plancher de suffisance vaut
$\PlancherSuffisance \pm \PlancherSuffisanceIC~\%$ : \textbf{la majorité avait
bien une cause qui, seule, aurait suffi}, et le choc est très concentré --- la
plus grosse créance perdue en porte $\ConcentrationChoc$ en moyenne. Mais la forme
forte est rare : $\ToutesCausesSuffisantes~\%$ seulement avaient \emph{toutes}
leurs causes individuellement suffisantes, pour une moyenne de
$\CausesSuffisantesMoyen$ cause suffisante par victime.

\inference{} La réserve est donc levée \textbf{dans les deux sens}. Compter les
parents surestime largement le nombre de causes qui auraient suffi. Mais
l'attribution à parent unique n'est pas arbitraire non plus : dans la majorité
des cas il existe bien une perte dominante certifiée suffisante par le majorant. Ce qui est faux, c'est de croire que \emph{chaque} parent était une
cause suffisante.

\figurine{rev_suffisance}{1.0}{\textbf{Plusieurs causes, rarement chacune
suffisante.} À gauche, la part du choc portée par la plus grosse créance perdue, sur
les victimes sur-déterminées : très concentrée. À droite, par bras, la part des
victimes recevant plusieurs chocs et, parmi elles, la part dont au moins une
cause était certainement suffisante. \textbf{La seconde série est un plancher.}
Les dénominateurs diffèrent ; IC95 inter-graines à droite. L'histogramme
regroupe les victimes, sans supposer leur indépendance.}

\section{Résultat annexe --- la mémoire du carnet de crédit}

\fait{} Deux rampes de cinq paliers parcourent la même échelle de partage,
l'une en montant, l'autre en descendant. Les écarts significatifs sont de un à
deux pour-cent --- $\RampeQuart \pm \RampeQuartIC$ au quart et
$\RampeDemi \pm \RampeDemiIC$ à la moitié --- et leur signe suit la direction
d'approche. Aux deux derniers paliers, l'écart n'est plus significatif.

\inference{} Ce que la rampe mesure est donc un \textbf{retard de relaxation},
et non une dépendance au chemin. Le mécanisme est chiffrable : un contrat
disparaît dès que l'une de ses deux extrémités meurt, soit environ quatre
pour-cent par pas --- une demi-vie de dix-sept pas. Après deux cents pas de
palier, il ne reste rien du palier précédent. \textbf{La mémoire du carnet est
réelle, mais elle s'efface bien avant que le palier ne finisse} ; ce qui relaxe
lentement, c'est la population.

\incertitude{} Une dépendance au chemin pourrait exister à des paliers plus
courts que la demi-vie du carnet. \textbf{Ce protocole ne peut pas le dire.}

\figurine{f20_rampes}{1.0}{\textbf{Montée et descente se superposent.}
Effectif et nombre de contrats vivants à chaque palier, selon le sens de
parcours. Les deux trajectoires coïncident à l'intervalle près.}

\clearpage
\part{Ce que rien de tout cela n'établit}

\section{Les limites structurelles}

Elles ne sont pas des réserves de politesse : chacune bloque une conclusion
qu'on pourrait être tenté de tirer.

\begin{enumerate}
\item \textbf{Le modèle n'est pas calibré.} Aucun paramètre n'est ajusté sur
      une économie réelle. Les ordres de grandeur des exposants simulés sont
      parfois proches de ceux de la littérature --- De Vries et Toda situent
      les exposants de Pareto du revenu du capital principalement entre 1 et 3
      sur 475 observations pays--année --- mais cette proximité ne valide rien
      et exige une conversion explicite : ces auteurs emploient l'exposant de
      fonction de survie, tandis que certains rapports du stage donnent un
      exposant de densité.
\item \textbf{Le joule est une unité comptable.} Le postulat monnaie--exergie
      justifie la lecture énergétique de tout le programme, mais il n'est ni
      imposé par les équations, ni validé par les simulations.
\item \textbf{L'entité est volontairement générique.} C'est ce qui donne au
      modèle sa portée formelle, et c'est aussi ce qui rend le choix d'une
      cible empirique moins immédiat : toute validation externe doit préciser
      le système concret auquel elle se rapporte.
\item \textbf{Les applications « personne » et « entreprise » sont deux
      paramétrisations du même moteur}, pas deux populations coexistantes. Le
      travail ne construit pas une économie articulant ménages et firmes.
\item \textbf{L'analogie de la machine n'est pas encodée.} Le moteur ne sépare
      ni outil, ni énergie incorporée, ni maintenance, ni service utile, ni
      classes ouvrière et capitaliste.
\item \textbf{Les prêts perpétuels, nominaux et sans recouvrement} sont des
      choix institutionnels forts, utiles pour isoler un mécanisme, non des
      descriptions universelles du crédit.
\item \textbf{La règle de rencontre} --- la plus riche et la plus pauvre d'un
      échantillon --- a une portée empirique qui reste à discuter.
\item \textbf{Les horizons numériques} ne permettent pas d'exclure des
      métastabilités très longues.
\item \textbf{La campagne de sensibilité n'est pas une analyse de Sobol}, et
      les extensions valeur nette et cycles sont des analyses \emph{post-hoc}
      menées sur les mêmes runs --- confirmées sur graines séparées, mais
      explicitement étiquetées comme telles.
\item \textbf{L'étiquette temporelle d'un pas n'est pas fixée}, et la longueur
      de la fenêtre d'agrégation peut modifier la forme d'une distribution de
      flux. Cette robustesse reste à tester.
\end{enumerate}

\section{Les six choses que la dernière campagne n'établit pas}

\incertitude{} Reprises telles quelles du rapport de campagne, parce qu'il vaut
mieux les écrire que les laisser supposer.

\begin{enumerate}
\item \textbf{L'existence des queues de Pareto n'est pas prouvée}.
      Les visualisations la suggèrent ; les tests et le temps disponible
      n'ont pas permis de conclure. Les exposants sont caractérisés sous
      hypothèse, sans impossibilité universelle de discrimination affirmée.
\item \textbf{L'identité de pondération est établie, pas toute la causalité
      démographique du partage}. La croissance du partage implicite sur une
      grille ne constitue pas non plus une démonstration pour toute paire.
\item \textbf{Un huitième de l'écart de branchement avec M4B reste non
      attribué}, et il est imputé --- sans vérification --- à la taille du
      bassin d'appariement.
\item \textbf{Le taux de suffisance n'est connu que par un plancher.}
      L'effacement de dette intra-cascade n'est pas tracé par
      l'instrumentation ; son incidence est mesurée à
      $\IncidenceEffacement~\%$, donc le plancher est serré, sans être exact.
\item \textbf{L'absence de dépendance au chemin n'est établie qu'à l'échelle
      des paliers employés.}
\item \textbf{Le canal de liquidité est muet} dans ce régime : la contagion est
      purement de bilan, et tout raisonnement sur la liquidité y est hors sujet.
\end{enumerate}

\section{Les travaux identifiés et non réalisés}

\incertitude{} Ils ne sont pas listés comme des projets, mais comme des trous
connus dans ce qui précède.

\begin{itemize}
\item \textbf{Les tests d'adéquation absolue, d'autosimilarité et l'analyse
      de coupure restent à compléter.} La comparaison de familles par Vuong
      a bien été menée au lot C ; ce n'est pas le même test. Il faut être
      exact sur le motif : ce n'était pas une décision théorique que ces
      questions soient sans intérêt, c'est qu'elles n'ont pas pu être menées
      proprement dans le temps disponible\footnote{Le plan de la campagne
      qualifiait les résultats antérieurs d'« acquis qu'il est interdit de
      refaire ». Cette formulation, relue le 24 août, décrit un périmètre
      imposé par le temps et non une clôture épistémologique.}.
\item \textbf{La limite continue n'a pas été étudiée.} Faire tendre le pas vers
      une petite échelle, tester la convergence vers une dynamique continue et
      vérifier si les simulations actuelles représentent ce régime est une
      étude ouverte. Le recalage temporel actuel laisse un résidu de 6 à 17~\%.
\item \textbf{La mobilité intergénérationnelle} exige une filiation absente
      du moteur. Les transitions de rang parmi les survivantes peuvent,
      elles, être étudiées sans filiation, en séparant les décès.
\item \textbf{La cible empirique de l'exposant des cascades d'entreprises}
      doit encore être fixée à partir d'une définition et d'un jeu de données
      compatibles avec la définition \emph{causale} des avalanches d'ici. C'est
      une condition falsificatrice explicite : si l'architecture impose un
      exposant durablement incompatible avec les données, une restructuration
      des hypothèses sera nécessaire.
\item \textbf{Les contractions d'énergie totale ont une baseline interne mais
      pas de cible alignée.} Une comparaison quantitative de durée et de
      fréquence exigera d'abord une correspondance temporelle.
\item \textbf{Une campagne sur le trio hétérogène $(K_{0,i},A_i,\gamma_i)$}
      est la suite identifiée : les lois de naissance, leurs corrélations et
      leur dépendance à l'état restent à définir, avec contrôle des échelles.
\end{itemize}

\section{Questions ouvertes et expériences discriminantes}
\label{sec:perspectives}
Le premier résultat à prolonger est la réponse collective à la technologie ;
le partage du surplus constitue une seconde voie de recherche. Les pistes
suivantes sont des intentions scientifiques, pas des résultats acquis ni
des protocoles définitivement arrêtés.

\paragraph{Hétérogénéité du trio $(K_0,A,\gamma)$.}
Il faut étudier ensemble la dotation initiale, la productivité et la concavité :
trio identique à toutes les naissances, tirage dans une distribution bornée,
corrélations entre composantes, ou dépendance à l'état de la simulation.
Comparer les conventions à $\kappa_0$ contrôlé puis libre permettrait de séparer
hétérogénéité technologique et déplacement d'échelle. Si $\gamma_i$ varie,
la dimension de $A_i$ varie aussi : un tirage uniforme de valeurs numériques
de $A_i$ n'a pas de sens indépendant de l'unité choisie. Il faut fixer une
échelle commune $K_{\rm ref}$ et, par exemple, spécifier
$a_i=A_i K_{\rm ref}^{\gamma_i-1}$, puis documenter la loi jointe de
$(K_{0,i}/K_{\rm ref},a_i,\gamma_i)$.
Les sorties à suivre incluent survie des cohortes, production par entité,
population, réseau de crédit et distributions de bilan ; quels indicateurs
endogènes résument leur évolution micro- et macroscopique reste une question.

\paragraph{Causalité entre pas et hypothèse de criticité.}
Les avalanches enregistrées ferment la cascade \emph{dans un pas}.
Cette convention de mesure ne signifie pas que la causalité du système
s'y limite : le carnet de dettes, les stocks et les cohortes persistent et
modifient les risques futurs. Une entité classée « racine » aujourd'hui peut
donc porter une fragilité accumulée hier. Inversement, la criticité n'est pas
définie en général comme un phénomène limité à un seul pas.
La SOC reste ouverte : il faudrait définir une réponse spatio-temporelle à
une perturbation, des horizons, un contrefactuel et une analyse de taille finie.
Ni $b_1$ intra-pas, ni une corrélation retardée, ni une droite log--log ne
suffisent à trancher seuls.

\paragraph{Sur-altruisme et temporalité du contrat.}
L'extension proposée consiste à accepter $p<0$ dans $rq=L+p\Delta$ :
la prêteuse reçoit initialement moins que sa perte productive de référence.
La dissipation et l'évolution des deux entités pourraient changer ce bilan
ultérieurement ; cette inversion est une hypothèse à tester, pas une
conséquence automatique du signe de $p$.\footnote{Le moteur de la campagne
impose actuellement $0\leq p\leq1$ pour la règle \texttt{bargain}.
Un prolongement devra spécifier l'admissibilité du service $rq$, le maintien
ou non du taux nominal, l'horizon, les faillites et le contrefactuel sans prêt.
Un gain productif instantané négatif n'est pas une valeur actualisée nette
négative : le modèle ne fixe pas de taux d'actualisation.}
Un premier test isolerait une paire, avec et sans contrat, puis répéterait
la comparaison dans une population ; seuls les seconds essais renseigneraient
les rétroactions démographiques et la propagation.

\paragraph{Distributions et comparaison empirique.}
Pareto demeure une hypothèse suggérée par les graphiques et imparfaitement
testée faute de temps. Le corps exponentiel n'est plus une cible obligatoire :
l'objectif est d'identifier la forme et le mécanisme, non de reproduire un
gabarit universel supposé. Comparer les familles sur le même échantillon et le
même domaine exige de tenir compte de leur flexibilité.\footnote{Pour deux
familles emboîtées ajustées correctement, la vraisemblance maximale ne diminue
pas lorsqu'on ajoute un paramètre ; pour des familles non emboîtées, le seul
nombre de paramètres ne classe pas les ajustements. Une comparaison par
AIC/BIC, validation hors échantillon ou test adapté complète l'adéquation
absolue. Les observations de plusieurs instantanés d'un même run ne sont pas
des réplications indépendantes.}
La priorité empirique est une distribution de bilan ou d'intérêts dont
l'unité statistique, la fenêtre et les conventions comptables soient
compatibles ; l'entité générique ne sera assimilée à un acteur réel qu'à
cette condition. Robustesse des queues, seuils de coupure, limite continue
et autres règles de crédit restent des expériences possibles.


\section{Où en est la question dorsale}

\paragraph{Premier volet --- l'interprétabilité.} Le modèle retrouve
endogènement une dizaine de traits qualitatifs d'une économie réelle, dont
plusieurs n'étaient pas visés : la localisation de l'inégalité dans les bilans,
la survie croissante avec la taille, l'asymétrie négative de l'activité, et la
démographie comme horloge de la persistance macroéconomique. \textbf{Cela
soutient l'interprétabilité du modèle sans la démontrer.} Il en manque autant,
et les manques sont instructifs : pas de longue
récession, pas de mobilité intergénérationnelle représentée, et une SOC non
établie. Le branchement est endogène : le plateau historique M4B proche de
$0{,}30$ n'est pas une borne universelle, comme le montre le contrôle M4.4.

\paragraph{Second volet --- le rebond.} Une amplification qualifiée de rebond
au sens interne du modèle est mesurée ; elle n'est pas un rebond énergétique
empiriquement calibré. \textbf{La réponse de long terme est sous- ou
super-proportionnelle selon la convention de dotation}, sans que la réponse
de production elle-même change nécessairement de signe. Le cas le
plus proche d'une politique réelle d'efficacité --- traiter une fraction des
acteurs --- produit un effet qui \emph{s'éteint}. Ce qui est acquis, et qui est
un vrai résultat, est que \textbf{l'organisation marchande modifie la réponse
agrégée à une amélioration technologique}, et qu'on peut dire par quel canal :
la contraction démographique induite par une dotation de naissance qui ne suit
pas l'échelle.

\paragraph{Et un résultat qui déborde la question.} L'institution de crédit
employée par toute la lignée paraissait partager le surplus de coopération à
parts égales. Elle le fait \emph{par contrat} ($\PartageParContrat$) et pas
\emph{par joule} ($\PartagePondere$), et l'écart est exactement
$\operatorname{Cov}(p,\Delta)/\mathbb{E}[\Delta]$. C'est une propriété de la
règle de taux, dérivable des seules fonctions du moteur. \textbf{Une
institution peut être équitable en moyenne sur les contrats et asservissante
sur la valeur} --- et c'est la seconde lecture que subit le système.

\clearpage
\appendix

\section{Le dictionnaire de correspondance modèle $\leftrightarrow$ données}
\label{sec:dictionnaire}

\incertitude{} \textbf{À lire avant toute comparaison externe.} Il existe pour
éviter les homonymies trompeuses : la priorité n'est pas de faire coïncider
deux histogrammes portant le même nom, mais de vérifier que l'unité
statistique, le bilan, le flux et la fenêtre temporelle décrivent des objets
réellement comparables.

\begin{longtable}{@{}p{0.15\textwidth}p{0.27\textwidth}p{0.24\textwidth}p{0.26\textwidth}@{}}
\toprule
\textbf{variable} & \textbf{sens interne exact} &
\textbf{analogue empirique possible} & \textbf{limite principale} \\
\midrule
\endhead
$K_i$ & taille intrinsèque productive, réelle et dépréciable &
  capacité ou actif productif abstrait &
  \textbf{pas} le patrimoine ni les capitaux propres comptables \\
$NW_i = K_i + C_i - D_i$ & valeur nette interne après créances et dettes &
  patrimoine net, capitaux propres &
  valorisation nominale simplifiée, sans prix de marché \\
production $A K_i^\gamma$ &
  puissance extractrice : capacité à grossir en fonction de la taille et de la
  technologie ; terme productif positif, distinct de la croissance nette &
  énergie ou travail utile ; historiquement proxy du chiffre d'affaires &
  érosion, intérêts, transferts et défauts modifient ensuite le bilan ; pas de
  mesure séparée du service final \\
intérêts reçus \code{int\_in} & transfert entrant versé par les emprunteuses &
  revenu monétaire usuel &
  ne doit pas être lu comme le seul revenu financier réel \\
principal $q$, créances $C$, dettes $D$ &
  mise à disposition de stock productif et droits nominaux associés &
  investissement dans une machine, financement d'un équipement &
  ni machine séparée, ni contrat réel, ni prix de l'actif \\
$E_t = \sum_i K_i$ & stock réel total présent dans le système &
  taille ou capacité énergétique agrégée &
  ne mesure \textbf{pas} la consommation finale d'énergie ou d'exergie \\
avalanche & morts reliées causalement par destruction de créances &
  cascade de faillites &
  les données réelles identifient rarement les liens racine--descendants \\
rapport $b$ & fraction moyenne de morts non racines &
  rapport de branchement d'une contagion &
  proximité de 1 = criticité de branchement, \textbf{pas} preuve de SOC \\
\bottomrule
\end{longtable}

\paragraph{Une clarification d'ontologie, faite en cours de route.} Le
\textbf{revenu monétaire usuel} d'une entité est représenté par la somme des
intérêts qu'elle perçoit sur une fenêtre. La production représente le versant
\emph{physique}. Des analyses antérieures ont utilisé la production cumulée
comme proxy du chiffre d'affaires et « production + intérêts reçus » comme
proxy du revenu brut ; ces grandeurs restent des \textbf{proxys statistiques
historiques}, valides comme définitions de leurs statistiques, non comme
identifications économiques démontrées. Les anciens résultats ne sont pas
renommés rétroactivement : leur définition exacte doit rester dans leurs
légendes.

\section{Provenance des figures}
\label{sec:provenance}

\fait{} Les figures historiques ont été produites par les campagnes ou
Simulation Lab, puis copiées. Les nouveaux graphiques de rédaction sont
recalculés par \chemin{recherche/revision_documents/build_visuals.py} depuis
les données conservées ; leurs points et incertitudes sont exportés.
Le tableau historique ci-dessous et l'annexe complémentaire donnent leur provenance. L'empreinte sous chaque nom est celle du fichier source au
moment de la copie : elle permet de vérifier qu'une figure de cette note est
bien celle du dossier d'origine.

\incertitude{} Les figures issues de \code{simulation\_lab\_data/} viennent
d'un répertoire \emph{non versionné}. C'est la raison de la copie plutôt que du
lien : une note qui cesserait de compiler parce qu'un run a été nettoyé ne
serait pas un document autonome.

% Engendré par scripts/collect_figures.py — NE PAS ÉDITER À LA MAIN.
% 40 figures.

{\setlength{\tabcolsep}{4pt}
\begin{longtable}{@{}>{\raggedright\arraybackslash}p{0.18\textwidth}>{\raggedright\arraybackslash}p{0.27\textwidth}>{\raggedright\arraybackslash}p{0.45\textwidth}@{}}
\toprule
\textbf{figure} & \textbf{provenance} & \textbf{ce qu'elle établit} \\
\midrule
\endhead
{\small\texttt{m01\_\allowbreak{}reseau\_\allowbreak{}credit}}\newline{\scriptsize\texttt{daa5f01fe9a4}} & Simulation Lab, run 20260718\_\allowbreak{}033746\_\allowbreak{}8b0c5acb (M4B, point central, graine 14) & Le réseau de crédit final : aucun type déclaré, et 182 entités sur 228 sont simultanément prêteuses et emprunteuses. \\
{\small\texttt{m02\_\allowbreak{}vue\_\allowbreak{}macro}}\newline{\scriptsize\texttt{2e48959eb0d1}} & Simulation Lab, run 20260718\_\allowbreak{}033746\_\allowbreak{}8b0c5acb & Capital, crédit nouveau, population et levier sur 4000 pas : le système est borné et stationnaire, sans dérive ni extinction. \\
{\small\texttt{t01\_\allowbreak{}artefact\_\allowbreak{}cohorte}}\newline{\scriptsize\texttt{e093fd184fac}} & recherche/\allowbreak{}analyse\_\allowbreak{}distributions\_\allowbreak{}taille\_\allowbreak{}revenu, campagne du WIP du 27 avril (163 runs stationnaires sur 805) & Le résultat négatif fondateur : la bimodalité lue comme « banques et travailleurs » est une structure d'âge, pas deux classes économiques. \\
{\small\texttt{t02\_\allowbreak{}derive\_\allowbreak{}queue}}\newline{\scriptsize\texttt{4408f6ce8620}} & même campagne & La première queue lourde n'était pas stationnaire : son exposant dérive avec l'horizon d'observation. \\
{\small\texttt{t03\_\allowbreak{}branchement\_\allowbreak{}m4}}\newline{\scriptsize\texttt{7a57ca8f11c0}} & m4\_\allowbreak{}credit\_\allowbreak{}soc\_\allowbreak{}fable, rapport 01\_\allowbreak{}soc\_\allowbreak{}final & M4 : le rapport de branchement se stabilise vers 0,30 une fois la règle annulation–destruction en place. \\
{\small\texttt{t04\_\allowbreak{}taille\_\allowbreak{}finie\_\allowbreak{}m4}}\newline{\scriptsize\texttt{26a7595238c4}} & m4\_\allowbreak{}credit\_\allowbreak{}soc\_\allowbreak{}fable, rapport 01\_\allowbreak{}soc\_\allowbreak{}final & M4 : la coupure des avalanches croît avec la taille du système — signature de taille finie, condition nécessaire d'un candidat SOC. \\
{\small\texttt{t05\_\allowbreak{}effet\_\allowbreak{}echelle\_\allowbreak{}m42}}\newline{\scriptsize\texttt{dcdd296f22ed}} & m4\_\allowbreak{}2\_\allowbreak{}credit\_\allowbreak{}soc, campagne M4.2 & M4.2 : l'effet attribué naïvement à la concavité γ est dominé par un déplacement d'échelle du capital. \\
{\small\texttt{p01\_\allowbreak{}lorenz}}\newline{\scriptsize\texttt{39a211d47ecf}} & Simulation Lab, run 20260718\_\allowbreak{}033746\_\allowbreak{}8b0c5acb & L'inégalité ne vit pas dans la taille productive (G = 0,215) mais dans le revenu d'intérêt (G = 0,595). \\
{\small\texttt{p02\_\allowbreak{}structure\_\allowbreak{}nw}}\newline{\scriptsize\texttt{a26543aae3ee}} & recherche/\allowbreak{}sensibilite\_\allowbreak{}m4b, campagne de 594 simulations & Structure de la valeur nette : cycle de vie débitrice puis créancière, et trou d'insolvabilité au décès. \\
{\small\texttt{p03\_\allowbreak{}vie\_\allowbreak{}par\_\allowbreak{}taille}}\newline{\scriptsize\texttt{8887ab830575}} & Simulation Lab, run 20260718\_\allowbreak{}033746\_\allowbreak{}8b0c5acb & L'espérance de vie restante croît avec la taille : les grandes entités survivent plus longtemps, sans être immortelles. \\
{\small\texttt{p04\_\allowbreak{}avalanches\_\allowbreak{}structure}}\newline{\scriptsize\texttt{4f7ede37f1bf}} & Simulation Lab, run 20260718\_\allowbreak{}033746\_\allowbreak{}8b0c5acb & Structure causale des avalanches : la part des racines décroît avec la taille et la profondeur atteint 10 — c'est de la propagation, pas de la synchronisation. \\
{\small\texttt{p05\_\allowbreak{}cascades\_\allowbreak{}rank\_\allowbreak{}size}}\newline{\scriptsize\texttt{95e15b6fec17}} & Simulation Lab, run 20260718\_\allowbreak{}033746\_\allowbreak{}8b0c5acb & Densité, CCDF et volume des cascades : loi de puissance tronquée, coupure visible dans la fenêtre observable. \\
{\small\texttt{p06\_\allowbreak{}scaling\_\allowbreak{}avalanches}}\newline{\scriptsize\texttt{b983649476fb}} & recherche/\allowbreak{}sensibilite\_\allowbreak{}m4b, phase de confirmation & Les lois d'échelle des avalanches en fonction de la taille du système : $s_{\max} \propto \lambda^{0,58}$ et $\chi \propto \lambda^{0,51}$. \\
{\small\texttt{p07\_\allowbreak{}hierarchie\_\allowbreak{}prcc}}\newline{\scriptsize\texttt{e27175372621}} & recherche/\allowbreak{}sensibilite\_\allowbreak{}m4b, criblage LHS & Hiérarchie des paramètres : $\sigma$, $k$, $K_0$, $\delta$ et $\lambda$ ont des rôles distincts et séparables. \\
{\small\texttt{p08\_\allowbreak{}gini\_\allowbreak{}renouvellement}}\newline{\scriptsize\texttt{99507cfc5529}} & recherche/\allowbreak{}sensibilite\_\allowbreak{}m4b, phase de confirmation & Gini et renouvellement sur graines de confirmation : les deux inégalités, capital et valeur nette, répondent parfois en sens opposé. \\
{\small\texttt{p09\_\allowbreak{}cycles\_\allowbreak{}structure}}\newline{\scriptsize\texttt{5f8d465986c5}} & recherche/\allowbreak{}sensibilite\_\allowbreak{}m4b, extension cycles (post-audit) & Structure des épisodes d'activité : courts, asymétriques, sans longue récession endogène. \\
{\small\texttt{p10\_\allowbreak{}cycles\_\allowbreak{}sensibilite}}\newline{\scriptsize\texttt{03e540350fef}} & recherche/\allowbreak{}sensibilite\_\allowbreak{}m4b, extension cycles (post-audit) & Ce qui pilote les épisodes : l'amplitude suit σ et K₀, la fréquence n'est expliquée de façon robuste par aucun paramètre. \\
{\small\texttt{p11\_\allowbreak{}renouvellement\_\allowbreak{}top}}\newline{\scriptsize\texttt{1fb0a1f62b71}} & Simulation Lab, run 20260718\_\allowbreak{}033746\_\allowbreak{}8b0c5acb & Le décile supérieur se renouvelle intégralement en moins de cent pas : c'est un renouvellement démographique, et non une mobilité sociale comparable aux données. \\
{\small\texttt{p12\_\allowbreak{}revenu\_\allowbreak{}empirique}}\newline{\scriptsize\texttt{daec7099081e}} & Simulation Lab, run 20260718\_\allowbreak{}033746\_\allowbreak{}8b0c5acb & Le revenu brut du modèle, sans ajustement imposé : ni corps exponentiel net, ni queue de Pareto convaincante sur cette observable. \\
{\small\texttt{p13\_\allowbreak{}queue\_\allowbreak{}interets}}\newline{\scriptsize\texttt{3d7090de2469}} & m4\_\allowbreak{}2b\_\allowbreak{}credit\_\allowbreak{}soc, campagne M4.2B & La queue d'intérêts telle qu'elle se voit sur les graphiques — l'observation qui a motivé le test de Pareto. \\
{\small\texttt{p14\_\allowbreak{}erreur\_\allowbreak{}seuil}}\newline{\scriptsize\texttt{7eae899c2345}} & m4\_\allowbreak{}2b\_\allowbreak{}credit\_\allowbreak{}soc, campagne M4.2B & Le contrôle d'autosimilarité mal orienté, puis corrigé : à seuil×2 la queue médiane tombe de 113 à 17 points. L'existence de Pareto reste une hypothèse. \\
{\small\texttt{r01\_\allowbreak{}reponse\_\allowbreak{}prod}}\newline{\scriptsize\texttt{13fa2cce2157}} & m4\_\allowbreak{}3live\_\allowbreak{}credit\_\allowbreak{}soc, campagne M4.3Live & La mesure centrale du rebond : à portée partielle la réponse est forte puis s'éteint avec la cohorte traitée ; à portée globale elle s'établit à +36 \% pour A×1,5, donc sous-proportionnelle. \\
{\small\texttt{r02\_\allowbreak{}ablation\_\allowbreak{}k0}}\newline{\scriptsize\texttt{3841b13b9646}} & m4\_\allowbreak{}3live\_\allowbreak{}credit\_\allowbreak{}soc, campagne M4.3Live & Le pivot du verdict : selon que la dotation de naissance K₀ suit ou non l'échelle technologique, la population se contracte ou non. \\
{\small\texttt{r03\_\allowbreak{}covariance\_\allowbreak{}echelle}}\newline{\scriptsize\texttt{90a37fc95450}} & m4\_\allowbreak{}3live\_\allowbreak{}credit\_\allowbreak{}soc, campagne M4.3Live & La covariance d'échelle du moteur, vérifiée numériquement : sous compensation de K₀, l'élasticité vaut exactement 1/(1−γ). \\
{\small\texttt{f01\_\allowbreak{}ancrages}}\newline{\scriptsize\texttt{380d0c4ce946}} & m4\_\allowbreak{}4\_\allowbreak{}rebond\_\allowbreak{}credit\_\allowbreak{}soc, engendrée par scripts/\allowbreak{}make\_\allowbreak{}figures.py & Les trois élasticités de la lignée sont reproduites sur une campagne neuve. \\
{\small\texttt{f05\_\allowbreak{}position\_\allowbreak{}nette}}\newline{\scriptsize\texttt{2a343fbc63d3}} & m4\_\allowbreak{}4\_\allowbreak{}rebond\_\allowbreak{}credit\_\allowbreak{}soc, engendrée par scripts/\allowbreak{}make\_\allowbreak{}figures.py & La position nette de crédit par âge, et son déplacement sous intervention. \\
{\small\texttt{f06\_\allowbreak{}mortalite\_\allowbreak{}rotation}}\newline{\scriptsize\texttt{d68a937c47b2}} & m4\_\allowbreak{}4\_\allowbreak{}rebond\_\allowbreak{}credit\_\allowbreak{}soc, engendrée par scripts/\allowbreak{}make\_\allowbreak{}figures.py & L'exposant qui relie mortalité et rotation du crédit n'est pas universel : il dépend du levier. \\
{\small\texttt{f07\_\allowbreak{}chaine}}\newline{\scriptsize\texttt{fc2e1cb4c39e}} & m4\_\allowbreak{}4\_\allowbreak{}rebond\_\allowbreak{}credit\_\allowbreak{}soc, engendrée par scripts/\allowbreak{}make\_\allowbreak{}figures.py & La chaîne causale de la contraction, maillon par maillon, avec la valeur de chaque élasticité. \\
{\small\texttt{f08\_\allowbreak{}avalanches\_\allowbreak{}ccdf}}\newline{\scriptsize\texttt{3e58411cb80c}} & m4\_\allowbreak{}4\_\allowbreak{}rebond\_\allowbreak{}credit\_\allowbreak{}soc, engendrée par scripts/\allowbreak{}make\_\allowbreak{}figures.py & Les CCDF d'avalanches par bras : la forme survit au changement de productivité. \\
{\small\texttt{f09\_\allowbreak{}b1\_\allowbreak{}b2}}\newline{\scriptsize\texttt{3a3b08fdc4ac}} & m4\_\allowbreak{}4\_\allowbreak{}rebond\_\allowbreak{}credit\_\allowbreak{}soc, engendrée par scripts/\allowbreak{}make\_\allowbreak{}figures.py & Les deux estimateurs de branchement, et pourquoi il en faut deux. \\
{\small\texttt{f10\_\allowbreak{}k0}}\newline{\scriptsize\texttt{2c1316cb7b87}} & m4\_\allowbreak{}4\_\allowbreak{}rebond\_\allowbreak{}credit\_\allowbreak{}soc, engendrée par scripts/\allowbreak{}make\_\allowbreak{}figures.py & La fragilité attribuée au rebond n'est pas la sienne : 97,9 \% de l'effet est annulé par la compensation de dotation. \\
{\small\texttt{f12\_\allowbreak{}alpha\_\allowbreak{}bras}}\newline{\scriptsize\texttt{62c50e375bc3}} & m4\_\allowbreak{}4\_\allowbreak{}rebond\_\allowbreak{}credit\_\allowbreak{}soc, engendrée par scripts/\allowbreak{}make\_\allowbreak{}figures.py & Les exposants de queue ne bougent pas d'un bras à l'autre. \\
{\small\texttt{f11\_\allowbreak{}couverture}}\newline{\scriptsize\texttt{c1408a137036}} & m4\_\allowbreak{}4\_\allowbreak{}rebond\_\allowbreak{}credit\_\allowbreak{}soc, engendrée par scripts/\allowbreak{}make\_\allowbreak{}figures.py & La couverture de queue n'est plus l'obstacle, et l'exposant ne bouge pas quand l'échantillon grandit de moitié. \\
{\small\texttt{f16\_\allowbreak{}bargain\_\allowbreak{}systeme}}\newline{\scriptsize\texttt{701670760a09}} & m4\_\allowbreak{}4\_\allowbreak{}rebond\_\allowbreak{}credit\_\allowbreak{}soc, engendrée par scripts/\allowbreak{}make\_\allowbreak{}figures.py & Ce que le partage déplace : branchement, tension, rotation et taux du carnet. \\
{\small\texttt{f18\_\allowbreak{}rotation\_\allowbreak{}rho}}\newline{\scriptsize\texttt{f8e3cc4e7bfe}} & m4\_\allowbreak{}4\_\allowbreak{}rebond\_\allowbreak{}credit\_\allowbreak{}soc, engendrée par scripts/\allowbreak{}make\_\allowbreak{}figures.py & Une relation héritée qui ne survit pas au changement de levier. \\
{\small\texttt{f19\_\allowbreak{}ablation}}\newline{\scriptsize\texttt{79aee50ff832}} & m4\_\allowbreak{}4\_\allowbreak{}rebond\_\allowbreak{}credit\_\allowbreak{}soc, engendrée par scripts/\allowbreak{}make\_\allowbreak{}figures.py & Ce qui explique l'écart de branchement avec la lignée M4B : 87,8 \% attribués, un huitième résiduel. \\
{\small\texttt{f20\_\allowbreak{}rampes}}\newline{\scriptsize\texttt{ecb7337f793c}} & m4\_\allowbreak{}4\_\allowbreak{}rebond\_\allowbreak{}credit\_\allowbreak{}soc, engendrée par scripts/\allowbreak{}make\_\allowbreak{}figures.py & Montée et descente du partage se superposent : retard de relaxation, pas dépendance au chemin. \\
{\small\texttt{f21\_\allowbreak{}hasard}}\newline{\scriptsize\texttt{e4d5f553d85a}} & m4\_\allowbreak{}4\_\allowbreak{}rebond\_\allowbreak{}credit\_\allowbreak{}soc, engendrée par scripts/\allowbreak{}make\_\allowbreak{}figures.py & La convexité posée ne tient pas : le risque n'est pas convexe dans le fardeau, et la courbure meurt au conditionnement. \\
{\small\texttt{f22\_\allowbreak{}ponderation}}\newline{\scriptsize\texttt{c5db42797aab}} & m4\_\allowbreak{}4\_\allowbreak{}rebond\_\allowbreak{}credit\_\allowbreak{}soc, engendrée par scripts/\allowbreak{}make\_\allowbreak{}figures.py & Le mécanisme, et c'est une identité : 0,53 par contrat, 0,94 par joule, l'écart est Cov(p,Δ)/E[Δ]. \\
{\small\texttt{g03\_\allowbreak{}partage\_\allowbreak{}analytique}}\newline{\scriptsize\texttt{692f06393086}} & m4\_\allowbreak{}4\_\allowbreak{}rebond\_\allowbreak{}credit\_\allowbreak{}soc, engendrée par scripts/\allowbreak{}make\_\allowbreak{}figures.py & Le partage dérivé des seules fonctions du moteur : ½ à la limite infinitésimale, au-delà de 1 sur les paires très inégales. \\
\bottomrule
\end{longtable}}


\subsection*{Traçabilité par identifiant de run}

\fait{} Règle permanente du projet : toute simulation qui apparaît dans une
figure, un tableau ou le texte doit être retrouvable dans l'interface locale par
son identifiant. Le tableau ci-dessous rattache chaque figure à son run ou à sa
campagne, avec les paramètres et le chemin où retrouver le détail run par run.

\incertitude{} Une figure qui agrège une campagne entière ne \emph{peut} pas
porter un identifiant unique : elle porte celui de la campagne et renvoie à
l'annexe engendrée de celle-ci, qui liste un \code{run\_id} par ligne ---
693~lignes pour la campagne M4B, 372 pour M4.4. Fabriquer un faux identifiant
unique serait pire que de renvoyer.

% Engendré par scripts/collect_figures.py — NE PAS ÉDITER À LA MAIN.

{\footnotesize\setlength{\tabcolsep}{4pt}
\begin{longtable}{@{}>{\raggedright\arraybackslash}p{0.19\textwidth}>{\raggedright\arraybackslash}p{0.27\textwidth}>{\raggedright\arraybackslash}p{0.22\textwidth}>{\raggedright\arraybackslash}p{0.26\textwidth}@{}}
\toprule
\textbf{identifiant} & \textbf{figures de la note} & \textbf{paramètres} & \textbf{où retrouver le détail} \\
\midrule
\endhead
\code{20260718\_\allowbreak{}033746\_\allowbreak{}8b0c5acb} & \code{m01\_\allowbreak{}reseau\_\allowbreak{}credit}, \code{m02\_\allowbreak{}vue\_\allowbreak{}macro}, \code{p01\_\allowbreak{}lorenz}, \code{p03\_\allowbreak{}vie\_\allowbreak{}par\_\allowbreak{}taille}, \code{p04\_\allowbreak{}avalanches\_\allowbreak{}structure}, \code{p05\_\allowbreak{}cascades\_\allowbreak{}rank\_\allowbreak{}size}, \code{p11\_\allowbreak{}renouvellement\_\allowbreak{}top}, \code{p12\_\allowbreak{}revenu\_\allowbreak{}empirique} & $\lambda = 30$, $\delta = 0{,}05$, $\sigma = 0{,}5$, $K_0 = 25$, $k = 2$, $T = 4000$, graine 14, statut \code{ok}, population finale 235 & \code{simulation\_\allowbreak{}lab\_\allowbreak{}data/\allowbreak{}runs/\allowbreak{}}\code{20260718\_\allowbreak{}033746\_\allowbreak{}8b0c5acb} --- \code{run.json}, \code{config.json}, \code{series.csv} \\
\addlinespace
campagne \code{sensibilite\_\allowbreak{}m4b} (594 runs) & \code{p02\_\allowbreak{}structure\_\allowbreak{}nw}, \code{p06\_\allowbreak{}scaling\_\allowbreak{}avalanches}, \code{p07\_\allowbreak{}hierarchie\_\allowbreak{}prcc}, \code{p08\_\allowbreak{}gini\_\allowbreak{}renouvellement}, \code{p09\_\allowbreak{}cycles\_\allowbreak{}structure}, \code{p10\_\allowbreak{}cycles\_\allowbreak{}sensibilite} & centre : $\lambda = 30$, $\delta = 0{,}05$, $\sigma = 0{,}25$, $K_0 = 25$, $k = 3$ ; graines confirmatoires 11--15, $T = 4000$ & annexe engendrée \code{recherche/\allowbreak{}sensibilite\_\allowbreak{}m4b/\allowbreak{}report/\allowbreak{}annexe\_\allowbreak{}tracabilite\_\allowbreak{}finale.tex} --- un \code{run\_\allowbreak{}id} par ligne \\
\addlinespace
campagne du WIP du 27 avril (805 runs, 163 stationnaires) & \code{t01\_\allowbreak{}artefact\_\allowbreak{}cohorte}, \code{t02\_\allowbreak{}derive\_\allowbreak{}queue} & filtre \code{bounded\_tail} et \code{flow\_balanced}, calculés par \code{detect\_regime\_}\allowbreak{}\code{diagnostics} du harnais & \code{recherche/\allowbreak{}analyse\_\allowbreak{}distributions\_\allowbreak{}taille\_\allowbreak{}revenu/\allowbreak{}} \\
\addlinespace
campagne \code{m4\_\allowbreak{}4\_\allowbreak{}rebond\_\allowbreak{}credit\_\allowbreak{}soc} (372 runs) & \code{f01\_\allowbreak{}ancrages}, \code{f05\_\allowbreak{}position\_\allowbreak{}nette}, \code{f06\_\allowbreak{}mortalite\_\allowbreak{}rotation}, \code{f07\_\allowbreak{}chaine}, \code{f08\_\allowbreak{}avalanches\_\allowbreak{}ccdf}, \code{f09\_\allowbreak{}b1\_\allowbreak{}b2}, \code{f10\_\allowbreak{}k0}, \code{f11\_\allowbreak{}couverture}, \code{f12\_\allowbreak{}alpha\_\allowbreak{}bras}, \code{f16\_\allowbreak{}bargain\_\allowbreak{}systeme}, \code{f18\_\allowbreak{}rotation\_\allowbreak{}rho}, \code{f19\_\allowbreak{}ablation}, \code{f20\_\allowbreak{}rampes}, \code{f21\_\allowbreak{}hasard}, \code{f22\_\allowbreak{}ponderation}, \code{g03\_\allowbreak{}partage\_\allowbreak{}analytique} & 12 graines appariées, amorçage partagé jusqu'au pas 2000, fenêtre sur les 1000 derniers pas & \code{m4\_\allowbreak{}4\_\allowbreak{}rebond\_\allowbreak{}credit\_\allowbreak{}soc/\allowbreak{}results/\allowbreak{}analysis/\allowbreak{}traceability.csv} et \code{simulation\_\allowbreak{}lab\_\allowbreak{}index.csv} \\
\addlinespace
campagne \code{m4\_\allowbreak{}credit\_\allowbreak{}soc\_\allowbreak{}fable} (lots 100) & \code{t03\_\allowbreak{}branchement\_\allowbreak{}m4}, \code{t04\_\allowbreak{}taille\_\allowbreak{}finie\_\allowbreak{}m4} & balayage en $\lambda$ pour la taille finie ; règle annulation--destruction & \code{m4\_\allowbreak{}credit\_\allowbreak{}soc\_\allowbreak{}fable/\allowbreak{}reports/\allowbreak{}01\_\allowbreak{}soc\_\allowbreak{}final/\allowbreak{}} \\
\addlinespace
campagne \code{m4\_\allowbreak{}2\_\allowbreak{}credit\_\allowbreak{}soc} & \code{t05\_\allowbreak{}effet\_\allowbreak{}echelle\_\allowbreak{}m42} & ablation d'échelle à $\gamma$ balayé, $K_0/K^*_{\mathrm{aut}}$ tenu fixe ou libre & \code{m4\_\allowbreak{}2\_\allowbreak{}credit\_\allowbreak{}soc/\allowbreak{}results/\allowbreak{}} et rapport de lignée \\
\addlinespace
campagne \code{m4\_\allowbreak{}2b\_\allowbreak{}credit\_\allowbreak{}soc} & \code{p13\_\allowbreak{}queue\_\allowbreak{}interets}, \code{p14\_\allowbreak{}erreur\_\allowbreak{}seuil} & queue des intérêts reçus ; seuils de Hill, contrôle d'autosimilarité à seuil$\times 2$ & \code{m4\_\allowbreak{}2b\_\allowbreak{}credit\_\allowbreak{}soc/\allowbreak{}results/\allowbreak{}} et rapport de lignée \\
\addlinespace
campagne \code{m4\_\allowbreak{}3live\_\allowbreak{}credit\_\allowbreak{}soc} & \code{r01\_\allowbreak{}reponse\_\allowbreak{}prod}, \code{r02\_\allowbreak{}ablation\_\allowbreak{}k0}, \code{r03\_\allowbreak{}covariance\_\allowbreak{}echelle} & 5 graines appariées, intervention à $t_0 = 2000$, horizon 2000 pas ; $A \times 1{,}5$ et $A \times 1{,}25$ & \code{m4\_\allowbreak{}3live\_\allowbreak{}credit\_\allowbreak{}soc/\allowbreak{}results/\allowbreak{}analysis/\allowbreak{}} et rapport de lignée \\
\addlinespace
\bottomrule
\end{longtable}}


\section{Les campagnes, et ce qu'elles ont coûté}

\begin{longtable}{@{}p{0.20\textwidth}p{0.16\textwidth}p{0.58\textwidth}@{}}
\toprule
\textbf{lignée} & \textbf{volume} & \textbf{objet} \\
\midrule
\endhead
WIP du 27 avril & 805 runs, 163 stationnaires &
  cartographie divergence/bornage ; c'est là qu'est établi l'artefact de
  cohortes \\
M4B --- sensibilité & 594 runs, $\approx$ 22 h CPU, 8,6 Go &
  régime stationnaire unique, hiérarchie des paramètres, avalanches, valeur
  nette, cycles \\
M4.2 & --- & effet d'échelle contre concavité \\
M4.2B & --- & queue des intérêts, non-décidabilité de Pareto,
  trois échelles d'incertitude \\
M4.3 & --- & levier qui épaissit la queue et renforce les avalanches \\
M4.3Live & --- & mesure du rebond par portée, horizon et convention d'échelle \\
M4.3Live-v2 & 153 runs & sens du prêt libéré ; campagnes A/B et de rotation \\
M4.4Rebond & 372 runs, $\approx$ 30 h &
  où va le surcroît dans la distribution ; le taux comme partage ;
  sur-détermination contre suffisance \\
\bottomrule
\end{longtable}

\fait{} Chaque run de la dernière campagne porte sur disque sa série complète,
ses tables d'événements, son arbre causal de cascades, ses panneaux par entité
et son checkpoint de fin. Le détail run par run --- graine, paramètres, statut,
durée, présence de chaque table --- est dans
\chemin{m4_4_rebond_credit_soc/results/analysis/traceability.csv}, engendré
par script depuis les répertoires de runs eux-mêmes\footnote{Le script
\emph{refuse} de produire l'annexe si une famille de runs est sans rôle
déclaré : un run qui a tourné, qui occupe du disque, et dont personne ne sait à
quelle question il répond, est un run qui n'aurait pas dû tourner.}. La suite
de tests du programme compte 19 tests, tous verts, pour 2204 secondes.

\section{Configurations, figures et données de reprise}
\label{sec:sources_revision}
Les références de figure ci-dessous sont celles de ce document. Les copies
locales permettent sa compilation autonome ; reproduire les mesures exige
aussi les données sources, dont les répertoires ne sont pas tous versionnés.
\subsection{Configurations de lecture}
\begin{longtable}{@{}p{.19\linewidth}p{.75\linewidth}@{}}
\toprule Référence & Paramètres et protocole \\
\midrule\endhead
M4B central & $\lambda=30$, $\sigma=0{,}25$, $\delta=0{,}05$, $K_0=25$, $k=3$ ; confirmation : graines 11--15, $T=4000$, analyse après $t=1000$. \\
M4B spécimen & $\lambda=30$, $\sigma=0{,}5$, $\delta=0{,}05$, $K_0=25$, $k=2$, graine 14 ; run \texttt{20260718\_033746\_8b0c5acb}. Ne pas substituer ses Gini à ceux du centre. \\
Live-v1 & $A=1$, $\gamma=1/2$, $\lambda=30$, $\sigma=\delta=0{,}01$, $K_0=25$ ; graines 0--4 ; intervention à $t=2001$ ; régime étudié $3000<t\le4000$. \\
Live-v2 / M4.4 & Référence homogène identique en paramètres physiques ; sens libre du prêt. M4.4 : graines 0--11 ; amorçage jusqu'à $2000$ ; bras technologiques et partage analysés sur $3000<t\le4000$, ablations sur $5000<t\le6000$. Les paramètres initiaux et interventions sont distingués ci-dessous. \\
\bottomrule\end{longtable}
\subsection{Chaque figure, sa source et son protocole}
{\footnotesize\begin{longtable}{@{}p{.12\linewidth}p{.37\linewidth}p{.45\linewidth}@{}}
\toprule Figure & Fichier source / données & Modèle, sélection, incertitude \\
\midrule\endhead
\ref{fig:m01_reseau_credit} & simulation\_\allowbreak{}lab\_\allowbreak{}data/\allowbreak{}runs/\allowbreak{}20260718\_\allowbreak{}033746\_\allowbreak{}8b0c5acb/\allowbreak{}figures/\allowbreak{}loan\_\allowbreak{}network\_\allowbreak{}final.png & M4B spécimen : paramètres ci-dessus ; graine 14. Image Simulation Lab,\allowbreak{}  non moyenne de campagne. \\
\ref{fig:m02_vue_macro} & simulation\_\allowbreak{}lab\_\allowbreak{}data/\allowbreak{}runs/\allowbreak{}20260718\_\allowbreak{}033746\_\allowbreak{}8b0c5acb/\allowbreak{}figures/\allowbreak{}macro\_\allowbreak{}overview.png & M4B spécimen : paramètres ci-dessus ; graine 14. Image Simulation Lab,\allowbreak{}  non moyenne de campagne. \\
\ref{fig:t01_artefact_cohorte} & recherche/\allowbreak{}analyse\_\allowbreak{}distributions\_\allowbreak{}taille\_\allowbreak{}revenu/\allowbreak{}figures/\allowbreak{}age\_\allowbreak{}cohorte.png & recherche/\allowbreak{}analyse\_\allowbreak{}distributions\_\allowbreak{}taille\_\allowbreak{}revenu,\allowbreak{}  campagne du WIP du 27 avril (163 runs stationnaires sur 805) \\
\ref{fig:t02_derive_queue} & recherche/\allowbreak{}analyse\_\allowbreak{}distributions\_\allowbreak{}taille\_\allowbreak{}revenu/\allowbreak{}figures/\allowbreak{}derive\_\allowbreak{}T.png & même campagne \\
\ref{fig:t03_branchement_m4} & m4\_\allowbreak{}credit\_\allowbreak{}soc\_\allowbreak{}fable/\allowbreak{}reports/\allowbreak{}01\_\allowbreak{}soc\_\allowbreak{}final/\allowbreak{}figures/\allowbreak{}lot100\_\allowbreak{}branching\_\allowbreak{}ratio.png & m4\_\allowbreak{}credit\_\allowbreak{}soc\_\allowbreak{}fable,\allowbreak{}  rapport 01\_\allowbreak{}soc\_\allowbreak{}final \\
\ref{fig:t04_taille_finie_m4} & m4\_\allowbreak{}credit\_\allowbreak{}soc\_\allowbreak{}fable/\allowbreak{}reports/\allowbreak{}01\_\allowbreak{}soc\_\allowbreak{}final/\allowbreak{}figures/\allowbreak{}lots\_\allowbreak{}scaling\_\allowbreak{}finite\_\allowbreak{}size.png & m4\_\allowbreak{}credit\_\allowbreak{}soc\_\allowbreak{}fable,\allowbreak{}  rapport 01\_\allowbreak{}soc\_\allowbreak{}final \\
\ref{fig:t05_effet_echelle_m42} & m4\_\allowbreak{}2\_\allowbreak{}credit\_\allowbreak{}soc/\allowbreak{}figures/\allowbreak{}ablation\_\allowbreak{}scale.png & m4\_\allowbreak{}2\_\allowbreak{}credit\_\allowbreak{}soc,\allowbreak{}  campagne M4.2 \\
\ref{fig:p02_structure_nw} & recherche/\allowbreak{}sensibilite\_\allowbreak{}m4b/\allowbreak{}figures/\allowbreak{}nw\_\allowbreak{}structure.png & recherche/\allowbreak{}sensibilite\_\allowbreak{}m4b,\allowbreak{}  campagne de 594 simulations \\
\ref{fig:p01_lorenz} & simulation\_\allowbreak{}lab\_\allowbreak{}data/\allowbreak{}runs/\allowbreak{}20260718\_\allowbreak{}033746\_\allowbreak{}8b0c5acb/\allowbreak{}figures/\allowbreak{}gini\_\allowbreak{}lorenz\_\allowbreak{}snapshots.png & M4B spécimen : paramètres ci-dessus ; graine 14. Image Simulation Lab,\allowbreak{}  non moyenne de campagne. \\
\ref{fig:p08_gini_renouvellement} & recherche/\allowbreak{}sensibilite\_\allowbreak{}m4b/\allowbreak{}figures/\allowbreak{}confirm\_\allowbreak{}gini\_\allowbreak{}renewal.png & recherche/\allowbreak{}sensibilite\_\allowbreak{}m4b,\allowbreak{}  phase de confirmation \\
\ref{fig:rev_lorenz} & data\_\allowbreak{}revision/\allowbreak{} : rev\_\allowbreak{}lorenz\_\allowbreak{}ginis.csv,\allowbreak{}  rev\_\allowbreak{}lorenz\_\allowbreak{}points.csv & M4.4/\allowbreak{}free/\allowbreak{}control/\allowbreak{}seed0--11 ; panels.npz ; instantané final t=4000. IC95 entre 12 instantanés indépendants ; pas les moyennes temporelles publiées ; Lorenz interpolé,\allowbreak{}  survie en points \\
\ref{fig:p03_vie_par_taille} & simulation\_\allowbreak{}lab\_\allowbreak{}data/\allowbreak{}runs/\allowbreak{}20260718\_\allowbreak{}033746\_\allowbreak{}8b0c5acb/\allowbreak{}figures/\allowbreak{}instantaneous\_\allowbreak{}life\_\allowbreak{}expectancy\_\allowbreak{}by\_\allowbreak{}size.png & M4B spécimen : paramètres ci-dessus ; graine 14. Image Simulation Lab,\allowbreak{}  non moyenne de campagne. \\
\ref{fig:p04_avalanches_structure} & simulation\_\allowbreak{}lab\_\allowbreak{}data/\allowbreak{}runs/\allowbreak{}20260718\_\allowbreak{}033746\_\allowbreak{}8b0c5acb/\allowbreak{}figures/\allowbreak{}soc\_\allowbreak{}avalanche\_\allowbreak{}structure.png & M4B spécimen : paramètres ci-dessus ; graine 14. Image Simulation Lab,\allowbreak{}  non moyenne de campagne. \\
\ref{fig:p05_cascades_rank_size} & simulation\_\allowbreak{}lab\_\allowbreak{}data/\allowbreak{}runs/\allowbreak{}20260718\_\allowbreak{}033746\_\allowbreak{}8b0c5acb/\allowbreak{}figures/\allowbreak{}cascades\_\allowbreak{}rank\_\allowbreak{}size.png & M4B spécimen : paramètres ci-dessus ; graine 14. Image Simulation Lab,\allowbreak{}  non moyenne de campagne. \\
\ref{fig:p06_scaling_avalanches} & recherche/\allowbreak{}sensibilite\_\allowbreak{}m4b/\allowbreak{}figures/\allowbreak{}confirm\_\allowbreak{}scaling.png & recherche/\allowbreak{}sensibilite\_\allowbreak{}m4b,\allowbreak{}  phase de confirmation \\
\ref{fig:p07_hierarchie_prcc} & recherche/\allowbreak{}sensibilite\_\allowbreak{}m4b/\allowbreak{}figures/\allowbreak{}lhs\_\allowbreak{}prcc.png & recherche/\allowbreak{}sensibilite\_\allowbreak{}m4b,\allowbreak{}  criblage LHS \\
\ref{fig:p09_cycles_structure} & recherche/\allowbreak{}sensibilite\_\allowbreak{}m4b/\allowbreak{}figures/\allowbreak{}cycles\_\allowbreak{}structure.png & recherche/\allowbreak{}sensibilite\_\allowbreak{}m4b,\allowbreak{}  extension cycles (post-audit) \\
\ref{fig:p10_cycles_sensibilite} & recherche/\allowbreak{}sensibilite\_\allowbreak{}m4b/\allowbreak{}figures/\allowbreak{}cycles\_\allowbreak{}sensibilite.png & recherche/\allowbreak{}sensibilite\_\allowbreak{}m4b,\allowbreak{}  extension cycles (post-audit) \\
\ref{fig:rev_stock_trajectoire} & data\_\allowbreak{}revision/\allowbreak{} : rev\_\allowbreak{}stock\_\allowbreak{}trajectoire\_\allowbreak{}points.csv & 20260718\_\allowbreak{}012820\_\allowbreak{}3a2e041c ; M4B centre\_\allowbreak{}lam30 ; graine 11 ; 3000<t<=3200. Trajectoire brute,\allowbreak{}  aucune incertitude de mesure ajoutée ; variabilité inter-graines dans rev\_\allowbreak{}cycles\_\allowbreak{}stock \\
\ref{fig:rev_cycles_stock} & data\_\allowbreak{}revision/\allowbreak{} : rev\_\allowbreak{}cycles\_\allowbreak{}episodes.csv,\allowbreak{}  rev\_\allowbreak{}cycles\_\allowbreak{}points.csv & 20260718\_\allowbreak{}012820\_\allowbreak{}3a2e041c; 20260718\_\allowbreak{}012820\_\allowbreak{}8832c88c; 20260718\_\allowbreak{}012820\_\allowbreak{}9e1b2711; 20260718\_\allowbreak{}012820\_\allowbreak{}5296e7ec; 20260718\_\allowbreak{}012820\_\allowbreak{}ec1bf58b ; M4B centre\_\allowbreak{}lam30 ; 1000<t<=4000. IC95 Student entre 5 graines (11--15) ; épisodes corrélés non traités comme répétitions indépendantes \\
\ref{fig:p12_revenu_empirique} & simulation\_\allowbreak{}lab\_\allowbreak{}data/\allowbreak{}runs/\allowbreak{}20260718\_\allowbreak{}033746\_\allowbreak{}8b0c5acb/\allowbreak{}figures/\allowbreak{}soc\_\allowbreak{}revenue\_\allowbreak{}fits.png & M4B spécimen : paramètres ci-dessus ; graine 14. Image Simulation Lab,\allowbreak{}  non moyenne de campagne. \\
\ref{fig:p13_queue_interets} & m4\_\allowbreak{}2b\_\allowbreak{}credit\_\allowbreak{}soc/\allowbreak{}report/\allowbreak{}figures/\allowbreak{}sim\_\allowbreak{}interets\_\allowbreak{}baseline.png & m4\_\allowbreak{}2b\_\allowbreak{}credit\_\allowbreak{}soc,\allowbreak{}  campagne M4.2B \\
\ref{fig:p14_erreur_seuil} & m4\_\allowbreak{}2b\_\allowbreak{}credit\_\allowbreak{}soc/\allowbreak{}report/\allowbreak{}figures/\allowbreak{}fig7\_\allowbreak{}seuil\_\allowbreak{}exemple\_\allowbreak{}travaille.png & m4\_\allowbreak{}2b\_\allowbreak{}credit\_\allowbreak{}soc,\allowbreak{}  campagne M4.2B \\
\ref{fig:p11_renouvellement_top} & simulation\_\allowbreak{}lab\_\allowbreak{}data/\allowbreak{}runs/\allowbreak{}20260718\_\allowbreak{}033746\_\allowbreak{}8b0c5acb/\allowbreak{}figures/\allowbreak{}soc\_\allowbreak{}top\_\allowbreak{}decile\_\allowbreak{}renewal.png & M4B spécimen : paramètres ci-dessus ; graine 14. Image Simulation Lab,\allowbreak{}  non moyenne de campagne. \\
\ref{fig:r01_reponse_prod} & m4\_\allowbreak{}3live\_\allowbreak{}credit\_\allowbreak{}soc/\allowbreak{}report/\allowbreak{}figures/\allowbreak{}response\_\allowbreak{}prod.png & m4\_\allowbreak{}3live\_\allowbreak{}credit\_\allowbreak{}soc,\allowbreak{}  campagne M4.3Live \\
\ref{fig:rev_reponse} & data\_\allowbreak{}revision/\allowbreak{} : rev\_\allowbreak{}reponse\_\allowbreak{}courbes.csv,\allowbreak{}  rev\_\allowbreak{}reponse\_\allowbreak{}estimateurs.csv,\allowbreak{}  rev\_\allowbreak{}reponse\_\allowbreak{}graines.csv & M4.3Live-v1/\allowbreak{}{control,\allowbreak{} all\_\allowbreak{}A150,\allowbreak{} new\_\allowbreak{}A150}/\allowbreak{}seed0--4 ; M4.4/\allowbreak{}free/\allowbreak{}{control,\allowbreak{} all\_\allowbreak{}A150,\allowbreak{} new\_\allowbreak{}A150,\allowbreak{} all\_\allowbreak{}A150\_\allowbreak{}K0comp}/\allowbreak{}seed0--11. IC95 Student entre graines ; trajectoires : moyenne mobile centrée de 101 pas par graine \\
\ref{fig:r03_covariance_echelle} & m4\_\allowbreak{}3live\_\allowbreak{}credit\_\allowbreak{}soc/\allowbreak{}report/\allowbreak{}figures/\allowbreak{}scaling\_\allowbreak{}covariance.png & m4\_\allowbreak{}3live\_\allowbreak{}credit\_\allowbreak{}soc,\allowbreak{}  campagne M4.3Live \\
\ref{fig:r02_ablation_k0} & m4\_\allowbreak{}3live\_\allowbreak{}credit\_\allowbreak{}soc/\allowbreak{}report/\allowbreak{}figures/\allowbreak{}ablation\_\allowbreak{}k0\_\allowbreak{}levels.png & m4\_\allowbreak{}3live\_\allowbreak{}credit\_\allowbreak{}soc,\allowbreak{}  campagne M4.3Live \\
\ref{fig:f01_ancrages} & m4\_\allowbreak{}4\_\allowbreak{}rebond\_\allowbreak{}credit\_\allowbreak{}soc/\allowbreak{}results/\allowbreak{}analysis/\allowbreak{}lotB\_\allowbreak{}elasticities.csv & arms/\allowbreak{}free ; contrastes résiduels,\allowbreak{}  12 graines \\
\ref{fig:rev_quantiles} & data\_\allowbreak{}revision/\allowbreak{} : rev\_\allowbreak{}quantiles\_\allowbreak{}graines.csv,\allowbreak{}  rev\_\allowbreak{}quantiles\_\allowbreak{}points.csv & M4.4/\allowbreak{}free/\allowbreak{}{control,\allowbreak{} all\_\allowbreak{}A150}/\allowbreak{}seed0--11 ; 3000<t<=4000. Student sur 12 rapports par graine \\
\ref{fig:rev_deciles} & data\_\allowbreak{}revision/\allowbreak{} : rev\_\allowbreak{}deciles\_\allowbreak{}points.csv & M4.4/\allowbreak{}free/\allowbreak{}{control,\allowbreak{} all\_\allowbreak{}A150}/\allowbreak{}seed0--11 ; 3000<t<=4000. Student des parts moyennes par run,\allowbreak{}  12 graines \\
\ref{fig:f05_position_nette} & m4\_\allowbreak{}4\_\allowbreak{}rebond\_\allowbreak{}credit\_\allowbreak{}soc/\allowbreak{}results/\allowbreak{}analysis/\allowbreak{}lotB\_\allowbreak{}distribution.csv & arms/\allowbreak{}free ; moyennes par run,\allowbreak{}  3000<t<=4000 \\
\ref{fig:f07_chaine} & m4\_\allowbreak{}4\_\allowbreak{}rebond\_\allowbreak{}credit\_\allowbreak{}soc/\allowbreak{}results/\allowbreak{}analysis/\allowbreak{}lotB\_\allowbreak{}elasticities.csv & arms/\allowbreak{}free ; all\_\allowbreak{}A150 ; 12 graines \\
\ref{fig:f06_mortalite_rotation} & m4\_\allowbreak{}4\_\allowbreak{}rebond\_\allowbreak{}credit\_\allowbreak{}soc/\allowbreak{}results/\allowbreak{}analysis/\allowbreak{}lotB\_\allowbreak{}gate.json & arms/\allowbreak{}free ; ajustements par graine \\
\ref{fig:f09_b1_b2} & m4\_\allowbreak{}4\_\allowbreak{}rebond\_\allowbreak{}credit\_\allowbreak{}soc/\allowbreak{}results/\allowbreak{}analysis/\allowbreak{}lotD\_\allowbreak{}avalanches.csv & arms/\allowbreak{}free et bargain ; un point par run,\allowbreak{}  pas par moyenne \\
\ref{fig:f10_k0} & m4\_\allowbreak{}4\_\allowbreak{}rebond\_\allowbreak{}credit\_\allowbreak{}soc/\allowbreak{}results/\allowbreak{}analysis/\allowbreak{}lotD\_\allowbreak{}avalanches.csv & arms/\allowbreak{}free/\allowbreak{}{control,\allowbreak{} all\_\allowbreak{}A150,\allowbreak{} all\_\allowbreak{}A150\_\allowbreak{}K0comp} ; niveaux,\allowbreak{}  12 graines \\
\ref{fig:f08_avalanches_ccdf} & m4\_\allowbreak{}4\_\allowbreak{}rebond\_\allowbreak{}credit\_\allowbreak{}soc/\allowbreak{}avalanches.csv dans les runs & arms/\allowbreak{}free/\allowbreak{}{control,\allowbreak{} all\_\allowbreak{}A150,\allowbreak{} new\_\allowbreak{}g060} ; quatre premiers dossiers triés ; 3000<t<=4000 \\
\ref{fig:rev_avalanches} & data\_\allowbreak{}revision/\allowbreak{} : rev\_\allowbreak{}avalanches\_\allowbreak{}comptages.csv,\allowbreak{}  rev\_\allowbreak{}avalanches\_\allowbreak{}fits.csv,\allowbreak{}  rev\_\allowbreak{}avalanches\_\allowbreak{}points.csv & M4.4/\allowbreak{}free/\allowbreak{}{control,\allowbreak{} all\_\allowbreak{}A150}/\allowbreak{}seed0--11 ; 3000<t<=4000. IC95 Student entre 12 survies ; ruban bootstrap des 12 runs (2000 réplications),\allowbreak{}  non test de Pareto \\
\ref{fig:f19_ablation} & m4\_\allowbreak{}4\_\allowbreak{}rebond\_\allowbreak{}credit\_\allowbreak{}soc/\allowbreak{}results/\allowbreak{}analysis/\allowbreak{}lotD\_\allowbreak{}avalanches.json & ablation ; 5000<t<=6000 ; autres bras 3000<t<=4000 \\
\ref{fig:f11_couverture} & m4\_\allowbreak{}4\_\allowbreak{}rebond\_\allowbreak{}credit\_\allowbreak{}soc/\allowbreak{}results/\allowbreak{}analysis/\allowbreak{}lotC0\_\allowbreak{}coverage.csv & arms/\allowbreak{}free et coverage ; voir groupe et instantané de chaque ligne \\
\ref{fig:rev_vuong} & data\_\allowbreak{}revision/\allowbreak{} : rev\_\allowbreak{}vuong\_\allowbreak{}points.csv & M4.4/\allowbreak{}arms/\allowbreak{}free et coverage ; lotC\_\allowbreak{}families.csv ; un point par ajustement,\allowbreak{}  bras et graine exportés. Observations individuelles,\allowbreak{}  sans IC de moyenne ajouté ; $\sigma$=10 est un seuil diagnostique,\allowbreak{}  non une frontière mathématique de validité \\
\ref{fig:f12_alpha_bras} & m4\_\allowbreak{}4\_\allowbreak{}rebond\_\allowbreak{}credit\_\allowbreak{}soc/\allowbreak{}results/\allowbreak{}analysis/\allowbreak{}lotC\_\allowbreak{}alpha.csv & arms/\allowbreak{}free et coverage ; exposants conditionnels aux seuils ajustés \\
\ref{fig:rev_institutions} & data\_\allowbreak{}revision/\allowbreak{} : rev\_\allowbreak{}institutions\_\allowbreak{}points.csv & M4.4/\allowbreak{}bargain/\allowbreak{}{p=0,\allowbreak{} p=0.25,\allowbreak{} p=0.5,\allowbreak{} p=0.75,\allowbreak{} p=1,\allowbreak{} marginal}/\allowbreak{}seed0--11 ; 3000<t<=4000. IC95 Student ; normalisations recalculées graine par graine,\allowbreak{}  incertitude du dénominateur incluse \\
\ref{fig:f21_hasard} & m4\_\allowbreak{}4\_\allowbreak{}rebond\_\allowbreak{}credit\_\allowbreak{}soc/\allowbreak{}results/\allowbreak{}analysis/\allowbreak{}lotI\_\allowbreak{}convexity.json & bargain ; risque à dix pas ; erreurs conditionnelles et IC par graine distincts \\
\ref{fig:f22_ponderation} & m4\_\allowbreak{}4\_\allowbreak{}rebond\_\allowbreak{}credit\_\allowbreak{}soc/\allowbreak{}results/\allowbreak{}analysis/\allowbreak{}lotI\_\allowbreak{}convexity.json & bargain ; pondération et risque agrégés par graine \\
\ref{fig:g03_partage_analytique} & m4\_\allowbreak{}4\_\allowbreak{}rebond\_\allowbreak{}credit\_\allowbreak{}soc/\allowbreak{}scripts/\allowbreak{}make\_\allowbreak{}figures.py,\allowbreak{}  fonction g03 & Grille analytique,\allowbreak{}  aucune simulation ni incertitude d'échantillonnage \\
\ref{fig:f16_bargain_systeme} & m4\_\allowbreak{}4\_\allowbreak{}rebond\_\allowbreak{}credit\_\allowbreak{}soc/\allowbreak{}results/\allowbreak{}analysis/\allowbreak{}bargain\_\allowbreak{}runs.csv & bargain ; p=0,\allowbreak{} 0.25,\allowbreak{} 0.5,\allowbreak{} 0.75,\allowbreak{} 1 et marginal \\
\ref{fig:rev_rho} & data\_\allowbreak{}revision/\allowbreak{} : rev\_\allowbreak{}rho\_\allowbreak{}fits.csv,\allowbreak{}  rev\_\allowbreak{}rho\_\allowbreak{}graines.csv,\allowbreak{}  rev\_\allowbreak{}rho\_\allowbreak{}points.csv & M4.4/\allowbreak{}control\_\allowbreak{}rho ; 5 niveaux ; graines 0--11 ; 3000<t<=4000. IC95 Student entre graines ; ruban de la moyenne des courbes ajustées par graine,\allowbreak{}  non intervalle de prédiction \\
\ref{fig:f18_rotation_rho} & m4\_\allowbreak{}4\_\allowbreak{}rebond\_\allowbreak{}credit\_\allowbreak{}soc/\allowbreak{}results/\allowbreak{}analysis/\allowbreak{}lotE\_\allowbreak{}rho\_\allowbreak{}runs.csv & control\_\allowbreak{}rho ; mêmes runs \\
\ref{fig:rev_suffisance} & data\_\allowbreak{}revision/\allowbreak{} : rev\_\allowbreak{}suffisance\_\allowbreak{}histogramme.csv,\allowbreak{}  rev\_\allowbreak{}suffisance\_\allowbreak{}points.csv & M4.4/\allowbreak{}arms/\allowbreak{}free ; contrôle et bras technologiques ; 12 graines ; 3000<t<=4000. Histogramme regroupé des victimes,\allowbreak{}  sans IC indépendant par victime ; à droite IC95 inter-graines de l'analyse source \\
\ref{fig:f20_rampes} & m4\_\allowbreak{}4\_\allowbreak{}rebond\_\allowbreak{}credit\_\allowbreak{}soc/\allowbreak{}results/\allowbreak{}analysis/\allowbreak{}lotT\_\allowbreak{}ramp.csv & bargain/\allowbreak{}ramp\_\allowbreak{}up et ramp\_\allowbreak{}down ; paliers temporels \\
\bottomrule\end{longtable}}
\subsection{Runs M4B des nouveaux graphiques de stock}
\begin{longtable}{@{}rl@{}}\toprule Graine & Identifiant Simulation Lab \\\midrule\endhead
11 & 20260718\_\allowbreak{}012820\_\allowbreak{}3a2e041c \\
12 & 20260718\_\allowbreak{}012820\_\allowbreak{}8832c88c \\
13 & 20260718\_\allowbreak{}012820\_\allowbreak{}9e1b2711 \\
14 & 20260718\_\allowbreak{}012820\_\allowbreak{}5296e7ec \\
15 & 20260718\_\allowbreak{}012820\_\allowbreak{}ec1bf58b \\
\bottomrule\end{longtable}
\subsection{Répertoire des configurations M4.4}
Chaque ligne décrit des runs effectivement retrouvés sur disque. Le fichier
local \chemin{data_revision/runs_m44.csv} conserve les 372 identifiants exacts,
paramètres, interventions et empreintes des métadonnées.
Les douze métadonnées d'amorçage coverage/burn ont un dictionnaire de paramètres vide : les valeurs absentes ne sont pas reconstruites par supposition. Les chemins sont
relatifs à \texttt{m4\_4\_rebond\_credit\_soc/results/campaign/}.
{\scriptsize\begin{longtable}{@{}p{.27\linewidth}p{.14\linewidth}p{.53\linewidth}@{}}
\toprule Dossier / bras & Graines / fin & Paramètres initiaux et interventions \\\midrule\endhead
ablation/\allowbreak{}delta005 & 0--11 ; T=6000 & lam=30.0,\allowbreak{}  sigma=0.01,\allowbreak{}  delta=0.05,\allowbreak{}  K0=25.0,\allowbreak{}  A=1.0,\allowbreak{}  gamma=0.5,\allowbreak{}  rho=1.0,\allowbreak{}  rate\_\allowbreak{}rule=marginal ; t=2001 : delta=0.05 (all) \\
ablation/\allowbreak{}m4b\_\allowbreak{}like & 0--11 ; T=6000 & lam=30.0,\allowbreak{}  sigma=0.25,\allowbreak{}  delta=0.05,\allowbreak{}  K0=25.0,\allowbreak{}  A=1.0,\allowbreak{}  gamma=0.5,\allowbreak{}  rho=1.0,\allowbreak{}  rate\_\allowbreak{}rule=marginal ; t=2001 : sigma=0.25 (all) ; t=2001 : delta=0.05 (all) \\
ablation/\allowbreak{}sigma005 & 0--11 ; T=6000 & lam=30.0,\allowbreak{}  sigma=0.05,\allowbreak{}  delta=0.01,\allowbreak{}  K0=25.0,\allowbreak{}  A=1.0,\allowbreak{}  gamma=0.5,\allowbreak{}  rho=1.0,\allowbreak{}  rate\_\allowbreak{}rule=marginal ; t=2001 : sigma=0.05 (all) \\
ablation/\allowbreak{}sigma010 & 0--11 ; T=6000 & lam=30.0,\allowbreak{}  sigma=0.1,\allowbreak{}  delta=0.01,\allowbreak{}  K0=25.0,\allowbreak{}  A=1.0,\allowbreak{}  gamma=0.5,\allowbreak{}  rho=1.0,\allowbreak{}  rate\_\allowbreak{}rule=marginal ; t=2001 : sigma=0.1 (all) \\
ablation/\allowbreak{}sigma025 & 0--11 ; T=6000 & lam=30.0,\allowbreak{}  sigma=0.25,\allowbreak{}  delta=0.01,\allowbreak{}  K0=25.0,\allowbreak{}  A=1.0,\allowbreak{}  gamma=0.5,\allowbreak{}  rho=1.0,\allowbreak{}  rate\_\allowbreak{}rule=marginal ; t=2001 : sigma=0.25 (all) \\
arms/\allowbreak{}free/\allowbreak{}all\_\allowbreak{}A150 & 0--11 ; T=4000 & lam=30.0,\allowbreak{}  sigma=0.01,\allowbreak{}  delta=0.01,\allowbreak{}  K0=25.0,\allowbreak{}  A=1.0,\allowbreak{}  gamma=0.5,\allowbreak{}  rho=1.0,\allowbreak{}  rate\_\allowbreak{}rule=marginal ; t=2001 : A=1.5 (all) \\
arms/\allowbreak{}free/\allowbreak{}all\_\allowbreak{}A150\_\allowbreak{}K0comp & 0--11 ; T=4000 & lam=30.0,\allowbreak{}  sigma=0.01,\allowbreak{}  delta=0.01,\allowbreak{}  K0=56.25,\allowbreak{}  A=1.0,\allowbreak{}  gamma=0.5,\allowbreak{}  rho=1.0,\allowbreak{}  rate\_\allowbreak{}rule=marginal ; t=2001 : A=1.5 (all) ; t=2001 : K0=56.25 (all) \\
arms/\allowbreak{}free/\allowbreak{}control & 0--11 ; T=4000 & lam=30.0,\allowbreak{}  sigma=0.01,\allowbreak{}  delta=0.01,\allowbreak{}  K0=25.0,\allowbreak{}  A=1.0,\allowbreak{}  gamma=0.5,\allowbreak{}  rho=1.0,\allowbreak{}  rate\_\allowbreak{}rule=marginal ; sans intervention \\
arms/\allowbreak{}free/\allowbreak{}new\_\allowbreak{}A075 & 0--11 ; T=4000 & lam=30.0,\allowbreak{}  sigma=0.01,\allowbreak{}  delta=0.01,\allowbreak{}  K0=25.0,\allowbreak{}  A=1.0,\allowbreak{}  gamma=0.5,\allowbreak{}  rho=1.0,\allowbreak{}  rate\_\allowbreak{}rule=marginal ; t=2001 : A=0.75 (new) \\
arms/\allowbreak{}richest\_\allowbreak{}lends/\allowbreak{}new\_\allowbreak{}A075 & 0--11 ; T=4000 & lam=30.0,\allowbreak{}  sigma=0.01,\allowbreak{}  delta=0.01,\allowbreak{}  K0=25.0,\allowbreak{}  A=1.0,\allowbreak{}  gamma=0.5,\allowbreak{}  rho=1.0,\allowbreak{}  rate\_\allowbreak{}rule=marginal ; t=2001 : A=0.75 (new) \\
arms/\allowbreak{}free/\allowbreak{}new\_\allowbreak{}A150 & 0--11 ; T=4000 & lam=30.0,\allowbreak{}  sigma=0.01,\allowbreak{}  delta=0.01,\allowbreak{}  K0=25.0,\allowbreak{}  A=1.0,\allowbreak{}  gamma=0.5,\allowbreak{}  rho=1.0,\allowbreak{}  rate\_\allowbreak{}rule=marginal ; t=2001 : A=1.5 (new) \\
arms/\allowbreak{}richest\_\allowbreak{}lends/\allowbreak{}new\_\allowbreak{}A150 & 0--11 ; T=4000 & lam=30.0,\allowbreak{}  sigma=0.01,\allowbreak{}  delta=0.01,\allowbreak{}  K0=25.0,\allowbreak{}  A=1.0,\allowbreak{}  gamma=0.5,\allowbreak{}  rho=1.0,\allowbreak{}  rate\_\allowbreak{}rule=marginal ; t=2001 : A=1.5 (new) \\
arms/\allowbreak{}free/\allowbreak{}new\_\allowbreak{}g060 & 0--11 ; T=4000 & lam=30.0,\allowbreak{}  sigma=0.01,\allowbreak{}  delta=0.01,\allowbreak{}  K0=25.0,\allowbreak{}  A=1.0,\allowbreak{}  gamma=0.5,\allowbreak{}  rho=1.0,\allowbreak{}  rate\_\allowbreak{}rule=marginal ; t=2001 : gamma=0.6 (new) \\
arms/\allowbreak{}richest\_\allowbreak{}lends/\allowbreak{}new\_\allowbreak{}g060 & 0--11 ; T=4000 & lam=30.0,\allowbreak{}  sigma=0.01,\allowbreak{}  delta=0.01,\allowbreak{}  K0=25.0,\allowbreak{}  A=1.0,\allowbreak{}  gamma=0.5,\allowbreak{}  rho=1.0,\allowbreak{}  rate\_\allowbreak{}rule=marginal ; t=2001 : gamma=0.6 (new) \\
bargain/\allowbreak{}marginal & 0--11 ; T=4000 & lam=30.0,\allowbreak{}  sigma=0.01,\allowbreak{}  delta=0.01,\allowbreak{}  K0=25.0,\allowbreak{}  A=1.0,\allowbreak{}  gamma=0.5,\allowbreak{}  rho=1.0,\allowbreak{}  rate\_\allowbreak{}rule=marginal ; sans intervention \\
bargain/\allowbreak{}p=0 & 0--11 ; T=4000 & lam=30.0,\allowbreak{}  sigma=0.01,\allowbreak{}  delta=0.01,\allowbreak{}  K0=25.0,\allowbreak{}  A=1.0,\allowbreak{}  gamma=0.5,\allowbreak{}  rho=1.0,\allowbreak{}  rate\_\allowbreak{}rule=bargain,\allowbreak{}  p=0.0 ; sans intervention \\
bargain/\allowbreak{}p=0.25 & 0--11 ; T=4000 & lam=30.0,\allowbreak{}  sigma=0.01,\allowbreak{}  delta=0.01,\allowbreak{}  K0=25.0,\allowbreak{}  A=1.0,\allowbreak{}  gamma=0.5,\allowbreak{}  rho=1.0,\allowbreak{}  rate\_\allowbreak{}rule=bargain,\allowbreak{}  p=0.25 ; sans intervention \\
bargain/\allowbreak{}p=0.5 & 0--11 ; T=4000 & lam=30.0,\allowbreak{}  sigma=0.01,\allowbreak{}  delta=0.01,\allowbreak{}  K0=25.0,\allowbreak{}  A=1.0,\allowbreak{}  gamma=0.5,\allowbreak{}  rho=1.0,\allowbreak{}  rate\_\allowbreak{}rule=bargain,\allowbreak{}  p=0.5 ; sans intervention \\
bargain/\allowbreak{}p=0.75 & 0--11 ; T=4000 & lam=30.0,\allowbreak{}  sigma=0.01,\allowbreak{}  delta=0.01,\allowbreak{}  K0=25.0,\allowbreak{}  A=1.0,\allowbreak{}  gamma=0.5,\allowbreak{}  rho=1.0,\allowbreak{}  rate\_\allowbreak{}rule=bargain,\allowbreak{}  p=0.75 ; sans intervention \\
bargain/\allowbreak{}p=1 & 0--11 ; T=4000 & lam=30.0,\allowbreak{}  sigma=0.01,\allowbreak{}  delta=0.01,\allowbreak{}  K0=25.0,\allowbreak{}  A=1.0,\allowbreak{}  gamma=0.5,\allowbreak{}  rho=1.0,\allowbreak{}  rate\_\allowbreak{}rule=bargain,\allowbreak{}  p=1.0 ; sans intervention \\
bargain/\allowbreak{}ramp\_\allowbreak{}down & 0--11 ; T=4000 & lam=30.0,\allowbreak{}  sigma=0.01,\allowbreak{}  delta=0.01,\allowbreak{}  K0=25.0,\allowbreak{}  A=1.0,\allowbreak{}  gamma=0.5,\allowbreak{}  rho=1.0,\allowbreak{}  rate\_\allowbreak{}rule=bargain,\allowbreak{}  p=0.0 ; t=2001 : bargain\_\allowbreak{}p=1.0 (all) ; t=2401 : bargain\_\allowbreak{}p=0.75 (all) ; t=2801 : bargain\_\allowbreak{}p=0.5 (all) ; t=3201 : bargain\_\allowbreak{}p=0.25 (all) ; t=3601 : bargain\_\allowbreak{}p=0.0 (all) \\
bargain/\allowbreak{}ramp\_\allowbreak{}up & 0--11 ; T=4000 & lam=30.0,\allowbreak{}  sigma=0.01,\allowbreak{}  delta=0.01,\allowbreak{}  K0=25.0,\allowbreak{}  A=1.0,\allowbreak{}  gamma=0.5,\allowbreak{}  rho=1.0,\allowbreak{}  rate\_\allowbreak{}rule=bargain,\allowbreak{}  p=1.0 ; t=2001 : bargain\_\allowbreak{}p=0.0 (all) ; t=2401 : bargain\_\allowbreak{}p=0.25 (all) ; t=2801 : bargain\_\allowbreak{}p=0.5 (all) ; t=3201 : bargain\_\allowbreak{}p=0.75 (all) ; t=3601 : bargain\_\allowbreak{}p=1.0 (all) \\
burn & 0--11 ; T=2000 & lam=30.0,\allowbreak{}  sigma=0.01,\allowbreak{}  delta=0.01,\allowbreak{}  K0=25.0,\allowbreak{}  A=1.0,\allowbreak{}  gamma=0.5,\allowbreak{}  rho=1.0,\allowbreak{}  rate\_\allowbreak{}rule=marginal ; sans intervention \\
control\_\allowbreak{}rho/\allowbreak{}rho=0.5 & 0--11 ; T=4000 & lam=30.0,\allowbreak{}  sigma=0.01,\allowbreak{}  delta=0.01,\allowbreak{}  K0=25.0,\allowbreak{}  A=1.0,\allowbreak{}  gamma=0.5,\allowbreak{}  rho=0.5,\allowbreak{}  rate\_\allowbreak{}rule=marginal ; t=2001 : rho=0.5 (all) \\
control\_\allowbreak{}rho/\allowbreak{}rho=1 & 0--11 ; T=4000 & lam=30.0,\allowbreak{}  sigma=0.01,\allowbreak{}  delta=0.01,\allowbreak{}  K0=25.0,\allowbreak{}  A=1.0,\allowbreak{}  gamma=0.5,\allowbreak{}  rho=1.0,\allowbreak{}  rate\_\allowbreak{}rule=marginal ; sans intervention \\
control\_\allowbreak{}rho/\allowbreak{}rho=1.5 & 0--11 ; T=4000 & lam=30.0,\allowbreak{}  sigma=0.01,\allowbreak{}  delta=0.01,\allowbreak{}  K0=25.0,\allowbreak{}  A=1.0,\allowbreak{}  gamma=0.5,\allowbreak{}  rho=1.5,\allowbreak{}  rate\_\allowbreak{}rule=marginal ; t=2001 : rho=1.5 (all) \\
control\_\allowbreak{}rho/\allowbreak{}rho=2 & 0--11 ; T=4000 & lam=30.0,\allowbreak{}  sigma=0.01,\allowbreak{}  delta=0.01,\allowbreak{}  K0=25.0,\allowbreak{}  A=1.0,\allowbreak{}  gamma=0.5,\allowbreak{}  rho=2.0,\allowbreak{}  rate\_\allowbreak{}rule=marginal ; t=2001 : rho=2.0 (all) \\
control\_\allowbreak{}rho/\allowbreak{}rho=3 & 0--11 ; T=4000 & lam=30.0,\allowbreak{}  sigma=0.01,\allowbreak{}  delta=0.01,\allowbreak{}  K0=25.0,\allowbreak{}  A=1.0,\allowbreak{}  gamma=0.5,\allowbreak{}  rho=3.0,\allowbreak{}  rate\_\allowbreak{}rule=marginal ; t=2001 : rho=3.0 (all) \\
coverage/\allowbreak{}all\_\allowbreak{}A150 & 0--11 ; T=4000 & lam=50.0,\allowbreak{}  sigma=0.01,\allowbreak{}  delta=0.01,\allowbreak{}  K0=25.0,\allowbreak{}  A=1.0,\allowbreak{}  gamma=0.5,\allowbreak{}  rho=1.0,\allowbreak{}  rate\_\allowbreak{}rule=marginal ; t=2001 : A=1.5 (all) \\
coverage/\allowbreak{}burn & 0,\allowbreak{}  0,\allowbreak{}  0,\allowbreak{}  0,\allowbreak{}  0,\allowbreak{}  0,\allowbreak{}  0,\allowbreak{}  0,\allowbreak{}  0,\allowbreak{}  0,\allowbreak{}  0,\allowbreak{}  0 ; T= & lam=non renseigne,\allowbreak{}  sigma=non renseigne,\allowbreak{}  delta=non renseigne,\allowbreak{}  K0=non renseigne,\allowbreak{}  A=non renseigne,\allowbreak{}  gamma=non renseigne,\allowbreak{}  rho=non renseigne,\allowbreak{}  rate\_\allowbreak{}rule=non renseigne ; sans intervention \\
coverage/\allowbreak{}control & 0--11 ; T=4000 & lam=50.0,\allowbreak{}  sigma=0.01,\allowbreak{}  delta=0.01,\allowbreak{}  K0=25.0,\allowbreak{}  A=1.0,\allowbreak{}  gamma=0.5,\allowbreak{}  rho=1.0,\allowbreak{}  rate\_\allowbreak{}rule=marginal ; sans intervention \\
\bottomrule\end{longtable}}
\subsection{Reproduire et interpréter les incertitudes}
Les nouveaux calculs sont produits par \texttt{recherche/revision\_documents/build\_visuals.py} ;
l'annexe par \texttt{build\_sources.py}. Le manifeste local conserve les empreintes SHA-256 des sources.
Les IC95 de Student reposent sur les graines indépendantes, pas sur une indépendance
supposée des entités ou des pas. Les rubans de régression de cascades rééchantillonnent
les runs entiers. Les images historiques conservent leur convention propre :
dispersion des observations, erreur binomiale conditionnelle, erreur-type ou IC ;
une barre ne doit pas être interprétée automatiquement comme un IC95.
L'environnement de vérification est inventorié dans \texttt{data\_revision/environnement.json}.
Cet inventaire ne constitue pas une installation propre testée sur une autre machine.


%% Références citées dans la note. Les identifiants Sxx renvoient au registre
%% recherche/memo_stage/registre_sources_scientifiques.md, qui donne pour
%% chacune son fichier canonique, son empreinte et ce qu'elle établit
%% exactement — ainsi que les précautions de lecture.
\section{Références}
\label{sec:biblio}

\begingroup\emergencystretch=2em
\begin{description}[leftmargin=1.6em,style=nextline,itemsep=0.35em]

\item[Hendrick, Rinaldo et Manoli (2025)] Martin Hendrick, Andrea Rinaldo et
  Gabriele Manoli, « A stochastic theory of urban metabolism », \emph{PNAS},
  122(33), e2501224122. DOI \code{10.1073/pnas.2501224122}. \textsf{[S01]}\\
  \emph{Déclencheur de l'hypothèse directrice.} L'article établit un
  \emph{collapse} de distributions urbaines après remise à l'échelle. Il ne
  démontre ni auto-organisation critique, ni mécanisme de rebond : le passage
  à « une société économique pourrait être auto-critique » est une hypothèse
  formulée à partir de cette lecture.

\item[Bak, Chen, Scheinkman et Woodford (1992)] « Self-Organized Criticality
  and Fluctuations in Economics », Santa Fe Institute Working Paper
  1992-04-018. \textsf{[S02]}\\
  \emph{Antériorité à reconnaître.} Un modèle-jouet d'économie de production
  auto-critique existe depuis 1992. Ce travail ne revendique pas l'idée
  générale d'appliquer l'auto-organisation critique à l'économie.

\item[Drăgulescu et Yakovenko (2001)] « Exponential and power-law probability
  distributions of wealth and income in the United Kingdom and the United
  States », \emph{Physica A}, 299, 213--221. DOI
  \code{10.1016/S0378-4371(01)00298-9}. \textsf{[S03]}\\
  \emph{Cible distributionnelle initiale} : corps exponentiel « thermique »,
  queue de Pareto « superthermique ». Le corps exponentiel n'a jamais été
  obtenu dans ce modèle.

\item[Yakovenko et Rosser (2009)] « Colloquium: Statistical mechanics of money,
  wealth, and income », \emph{Reviews of Modern Physics}, 81, 1703--1725. DOI
  \code{10.1103/RevModPhys.81.1703}. \textsf{[S04]}\\
  Synthèse du programme précédent. Dans les données fiscales américaines
  qu'elle rassemble, la queue concerne environ 3~\% des déclarations en 1997 et
  son exposant varie dans le temps : \textbf{il n'existe pas de cible
  universelle} pour les revenus.

\item[Herbert, Giraud, Louis-Napoléon et Goupil (2022)] « Macroeconomic
  Dynamics in a Finite World: the Thermodynamic Potential Approach »,
  prépublication arXiv:2204.02038v2. \textsf{[S05]}\\
  \emph{Continuité théorique du postulat monnaie--exergie.} Ces auteurs fondent
  une macrodynamique de monde fini sur le couplage de sphères physique et
  économique \textbf{distinctes} : leur travail soutient la dépendance
  matérielle, \emph{pas} une identité monnaie--exergie.

\item[Dupont, Jeanmart et Possoz (2017)] « Vers un monde 100~\% renouvelable ? »,
  \emph{Regards économiques}, 135, p. 5--16. \textsf{[S06]}\\
  \emph{Illustration du couplage PIB--énergie} entre 1965 et 2015. La source
  mesure l'\textbf{énergie} et non l'exergie, et documente aussi un découplage
  relatif : elle illustre le postulat sans le valider.

\item[Ormerod et Mounfield (2001)] « Power law distribution of the duration and
  magnitude of recessions in capitalist economies: Breakdown of scaling »,
  \emph{Physica A}, 293(3--4), 573--582. DOI
  \code{10.1016/S0378-4371(01)00108-X}. \textsf{[S09]}\\
  \emph{Cible externe potentielle pour les récessions.} Elle impose de comparer
  les \textbf{distributions complètes} de durée et d'amplitude, et leurs
  ruptures éventuelles, pas seulement des moyennes. Cette comparaison n'a pas
  été faite : elle exige d'abord une correspondance temporelle.

\item[de Vries et Toda (2021)] « Capital and Labor Income Pareto Exponents
  across Time and Space », prépublication arXiv:2006.03441v3. \textsf{[S12]}\\
  \emph{Repère de comparaison le plus large disponible} : sur 475 observations
  pays--année couvrant 52 pays entre 1967 et 2018, les exposants du revenu du
  capital se situent principalement entre 1 et 3, contre 2 à 5 pour le revenu
  du travail. \textbf{Précaution obligatoire} : ces auteurs emploient
  l'exposant de fonction de survie, tandis que certains rapports de ce travail
  donnent un exposant de densité. La conversion doit être explicite avant toute
  comparaison.

\item[Sorrell et Dimitropoulos (2008)] Steve Sorrell et John Dimitropoulos,
  « The rebound effect: Microeconomic definitions, limitations and extensions »,
  \emph{Ecological Economics}, 65(3), 636--649.
  \href{https://doi.org/10.1016/j.ecolecon.2007.08.013}{DOI de l'article}.\\
  Distingue les définitions et hypothèses de mesure du rebond direct ;
  motive la séparation entre service utile, ressource et production du modèle.

\item[Wright (2005)] Ian Wright, « The duration of recessions follows an
  exponential not a power law », \emph{Physica A}, 345(3--4), 608--610.
  \href{https://doi.org/10.1016/j.physa.2004.07.035}{DOI : 10.1016/j.physa.2004.07.035}.
  Prépublication \href{https://arxiv.org/abs/cond-mat/0311585}{cond-mat/0311585}.
  \textsf{[S11]}\\
  Recompare les durées de récession d'Ormerod et Mounfield : l'exponentielle
  s'ajuste mieux aux données complètes de leur échantillon. Le traitement des
  épisodes d'un an change la comparaison. Ne pas confondre ce texte court avec
  l'autre article de Wright, \emph{The Social Architecture of Capitalism}.

\item[Reed (2003) ; Reed et Jorgensen (2004)] « The Pareto law of incomes---an
  explanation and an extension », \emph{Physica A}, 319, 469--486, DOI
  \code{10.1016/S0378-4371(02)01507-8} ; « The Double Pareto-Lognormal
  Distribution », \emph{Communications in Statistics}, 33(8), 1733--1753, DOI
  \code{10.1081/STA-120037438}.\\
  \emph{Loi de référence du processus de croissance-renouvellement} auquel M2
  se ramène. Son intérêt n'est pas d'être flexible à quatre paramètres : c'est
  la loi de l'état d'un mouvement brownien géométrique partant d'un état
  log-normal, observé après une durée exponentielle. C'est cette dérivation
  \emph{générative} qui la rend pertinente ici.

\item[Clauset, Shalizi et Newman (2009)] « Power-Law Distributions in
  Empirical Data », \emph{SIAM Review}, 51(4), 661--703.
  DOI \chemin{10.1137/070710111}.\\
  \emph{Méthode d'identification des lois de puissance} (maximum de
  vraisemblance et test de Kolmogorov--\allowbreak Smirnov) dont s'inspirent les scripts de
  test de queue de ce travail. C'est la source de la mise en garde
  « une régression log-log à bon $r^2$ ne suffit pas ».

\item[Wildauer et Heck (version septembre 2024)] Rafael Wildauer et Ines Heck,
  « Was Pareto right? Do the US wealth and income distributions follow power
  laws? », \emph{Greenwich Papers in Political Economy}, GPERC92
  (page de garde de série : 2023 ; version du texte : septembre 2024).
  \href{https://gala.gre.ac.uk/id/eprint/38597/13/38597\%20WILDAUER_Was_Pareto_right_Is_the_distribution_of\%20wealth_thick_tailed_(REVISED)_2023.pdf}{Archive institutionnelle de Greenwich}.\\
  Distingue Pareto strict et comportements asymptotiques ; les alternatives
  ne sont pas interchangeables. Le soutien à une queue asymptotiquement en
  puissance ne valide pas une loi de Pareto stricte sur tout le domaine.

\item[Shaikh, Papanikolaou et Wiener (2014)] Anwar Shaikh, Nikolaos
  Papanikolaou et Noe Wiener, « Race, gender and the econophysics of income
  distribution in the USA », \emph{Physica A}, 415, 54--60.
  \href{https://doi.org/10.1016/j.physa.2014.07.043}{DOI : 10.1016/j.physa.2014.07.043}.\\
  Étudie le programme exponentiel sur des sous-groupes de revenus du travail
  américains. Cette reprise ne démontre pas un consensus ni une loi universelle.

\item[Bouchaud (2024)] Jean-Philippe Bouchaud, « The Self-Organized Criticality
  Paradigm in Economics \& Finance », \href{https://arxiv.org/abs/2407.10284v2}{arXiv:2407.10284v2},
  6 septembre 2024. \textsf{[S07]}\\
  Distingue propagation critique et organisation endogène vers la criticité ;
  un rapport de branchement observé ne suffit pas au second énoncé.

\item[Dickman, Vespignani et Zapperi (1998)] Ronald Dickman, Alessandro
  Vespignani et Stefano Zapperi, « Self-organized criticality as an
  absorbing-state phase transition », \emph{Physical Review E}, 57(5),
  5095--5105. \href{https://doi.org/10.1103/PhysRevE.57.5095}{DOI : 10.1103/PhysRevE.57.5095}.
  \textsf{[S08]}\\
  Le substrat peut garder mémoire d'épisodes précédents ; la dynamique rapide
  d'une avalanche et l'évolution du système ne sont pas le même objet.

\item[Ribeiro et Rybski (2023)] Fabiano L. Ribeiro et Diego Rybski,
  « Mathematical models to explain the origin of urban scaling laws »,
  \emph{Physics Reports}, 1012, 1--39.
  \href{https://doi.org/10.1016/j.physrep.2023.02.002}{DOI : 10.1016/j.physrep.2023.02.002}.
  \textsf{[S10]}\\
  Plusieurs mécanismes produisent des lois d'échelle : la forme seule ne
  permet pas d'identifier le mécanisme de ce modèle.

\item[West, Brown et Enquist (1997)] Geoffrey B. West, James H. Brown et
  Brian J. Enquist, « A General Model for the Origin of Allometric Scaling
  Laws in Biology », \emph{Science}, 276(5309), 122--126.
  \href{https://doi.org/10.1126/science.276.5309.122}{DOI de l'article}.\\
  Modèle de transport en réseaux biologiques ; il éclaire la motivation
  allométrique sans imposer un exposant universel aux entités économiques.

\end{description}

Le registre complet, dans le dossier \chemin{recherche/memo_stage/}, fichier
\chemin{registre_sources_scientifiques.md}, donne pour
chaque source son fichier canonique, son empreinte, ce qu'elle établit
exactement et ce qu'elle \emph{n'établit pas}. Les entrées ci-dessus sont
identifiables sans consulter ce registre. Les dates de version des
prépublications sont conservées pour ne pas confondre versions et publications.
\endgroup


\section{Où trouver quoi}

\begin{longtable}{@{}p{0.42\textwidth}p{0.54\textwidth}@{}}
\toprule
\textbf{objet} & \textbf{emplacement} \\
\midrule
\endhead
Dossier de recherche complet du stage &
  \chemin{recherche/memo_stage/memo_recherche_complet.md} \\
Corrections non encore intégrées à ce dossier &
  \chemin{recherche/memo_stage/note_consolidation_questionnaire_2026-08-24.md} \\
Provenance des codes, données et matériaux &
  \chemin{recherche/memo_stage/inventaire_sources_preuves.md} \\
Rapport de la campagne de sensibilité M4B &
  \chemin{recherche/sensibilite_m4b/report/rapport_final.pdf} \\
Mécanique détaillée du moteur de référence &
  \chemin{m4b_credit_soc_mini/report/mecanique_m4b.pdf} \\
Rapport de la dernière campagne &
  \chemin{m4_4_rebond_credit_soc/report/rapport_final.pdf} \\
Rapport de conception de la dernière campagne &
  \chemin{m4_4_rebond_credit_soc/report/conception_m4_4.pdf} \\
Mesure du rebond &
  \chemin{m4_3live_credit_soc/report/} \\
Interface locale et runs &
  \chemin{simulation_lab/}, \chemin{simulation_lab_data/} \\
Reprise pratique et questions scientifiques ouvertes &
  \chemin{recherche/note_communication_encadrants/REPRISE.md},
  section~\ref{sec:perspectives} \\
\bottomrule
\end{longtable}

\end{document}
